\begin{solution}{3.01}
(1) Gauss's law no longer holds. It suffices to examine the field of a point charge to show this. Take the spherical surface of radius $r$ centred on the point charge $q$ as the Gaussian surface and calculate the flux of the electric field through this surface. From the statement of the problem,
\[
\oint \vec{E}\cdot\dd\vec{S}=E\cdot 4\pi r^2\propto \frac{q}{r^3}\cdot r^2=\frac{q}{r},
\]
which no longer depends only on the charge inside the sphere.

(2) The concept of potential can still be introduced. Although the force between the charges is inversely proportional to the cube of the distance, it is still only a function of the distance; hence the electric force is still a conservative central force, and therefore the concept of potential can still be introduced. But the potential of a point charge is no longer inversely proportional to the distance (if the potential zero is still taken at infinity), but inversely proportional to the square of the distance.

(3) In this case an electric field will exist inside a uniformly positively charged spherical surface, and the direction of the field is along the radius towards the centre of the sphere. Consider the field at any point $P$ inside the sphere. Take $P$ as the vertex and construct two small opposite cones (see the figure); they cut out the surface elements $\dd S_1$ and $\dd S_2$ on the sphere. From the geometry, $\dd S_1\propto r_1^2$, $\dd S_2\propto r_2^2$. From the statement of the problem,
\[
\dd E_1\propto\frac{\dd S_1}{r_1^3}\propto\frac{r_1^2}{r_1^3}=\frac{1}{r_1},
\qquad
\dd E_2\propto\frac{1}{r_2}.
\]
It is thus seen that $\dd E_1>\dd E_2$. From the superposition principle and the symmetry it follows that $E_P$ points towards the centre of the sphere.

(4) The charge is no longer distributed on the surface. There is also a charge distribution inside the conductor. Taking a conducting sphere as an example, from the discussion in part (3) it is known that if the charge were distributed only on the spherical surface, the field inside the sphere would not be zero; hence the charge will move towards the interior of the sphere until the field inside the sphere is zero everywhere.

\begin{figbox}{3.01}
\begin{tikzpicture}[scale=0.95,>=latex]
  % sphere
  \draw[thickline] (0,0) circle (2.2);
  \fill (0,0) circle (1.0pt);
  \node[below left] at (-0.05,-0.12) {$O$};
  % point P
  \fill (0.6,0.7) circle (1.2pt);
  \node[left] at (0.5,0.85) {$P$};
  % upper-right cone
  \coordinate (A1) at (1.414,1.685);
  \coordinate (A2) at (1.802,1.262);
  \draw[dashedline] (0.6,0.7)--(A1);
  \draw[dashedline] (0.6,0.7)--(A2);
  \draw[thickline] (A1)--(A2);
  \node[right] at (1.72,1.62) {$\dd S_1$};
  \node[below left] at (1.15,1.28) {$r_1$};
  % lower-left cone
  \coordinate (B1) at (-1.802,-1.262);
  \coordinate (B2) at (-2.067,-0.752);
  \draw[dashedline] (0.6,0.7)--(B1);
  \draw[dashedline] (0.6,0.7)--(B2);
  \draw[thickline] (B1)--(B2);
  \node[left] at (-2.12,-1.02) {$\dd S_2$};
  \node[above right] at (-0.65,-0.05) {$r_2$};
\end{tikzpicture}
\end{figbox}

\begin{commentary}
This problem involves two basic implications of Coulomb's law: the central force and the inverse square. The former is the basis for introducing the potential, the latter the foundation for deriving Gauss's law. Hence the inverse-cube assumption of this problem does not affect the introduction of the potential, but it destroys the conditions under which Gauss's law holds, and also destroys certain results connected with Gauss's law.
\end{commentary}
\end{solution}

\begin{solution}{3.02}
(1) The electric field produced by dipole $A$ at the position of dipole $B$ is
\begin{align*}
E_r&=\frac{1}{4\pi\epsilon_0}\frac{2p_1\cos\theta}{r^3},\\
E_\theta&=\frac{1}{4\pi\epsilon_0}\frac{p_1\sin\theta}{r^3}.
\end{align*}
Hence the interaction energy of the two dipoles is
\begin{align*}
W&=-\vec{p}_2\cdot\vec{E}
=-\bigl[p_2\cos(\varphi-\theta)\,\hat{r}+p_2\sin(\varphi-\theta)\,\hat{\theta}\bigr]
\cdot\bigl(E_r\,\hat{r}+E_\theta\,\hat{\theta}\bigr)\\
&=-\frac{p_1p_2}{4\pi\epsilon_0r^3}\bigl[2\cos\theta\cos(\varphi-\theta)+\sin\theta\sin(\varphi-\theta)\bigr]\\
&=-\frac{p_1p_2}{4\pi\epsilon_0r^3}\bigl\{3\cos\theta\cos(\varphi-\theta)
-\bigl[\cos\theta\cos(\varphi-\theta)-\sin\theta\sin(\varphi-\theta)\bigr]\bigr\}\\
&=-\frac{p_1p_2}{4\pi\epsilon_0r^3}\bigl[3\cos\theta\cos(\varphi-\theta)-\cos\varphi\bigr].
\end{align*}
Therefore
\[
F_r=-\frac{\partial W}{\partial r}
=-\frac{3p_1p_2}{4\pi\epsilon_0r^4}\bigl[3\cos\theta\cos(\varphi-\theta)-\cos\varphi\bigr].
\]
\[
F_\theta=-\frac{1}{r}\frac{\partial W}{\partial\theta}
=-\frac{3p_1p_2}{4\pi\epsilon_0r^4}\bigl[\sin\theta\cos(\varphi-\theta)-\cos\theta\sin(\varphi-\theta)\bigr]
=-\frac{3p_1p_2}{4\pi\epsilon_0r^4}\sin(2\theta-\varphi).
\]

(2) The electric field produced by the electric quadrupole at the point $P$ may be regarded as the superposition of the electric fields produced at $P$ by two oppositely directed dipoles. The two oppositely directed dipoles have dipole moments equal in magnitude, $p_+=p_-=ql=p$, but opposite in direction. For each dipole the potential is
\[
\Phi_+=\frac{1}{4\pi\epsilon_0}\frac{p_+\cos\theta_+}{r_+^2},
\qquad
\Phi_-=\frac{1}{4\pi\epsilon_0}\frac{p_-\cos\theta_-}{r_-^2},
\]
where $\theta_+=\theta+\Delta\theta$, $\theta_-=\theta-\Delta\theta$, $r_+=r-\dfrac{l}{2}\cos\theta$, $r_-=r+\dfrac{l}{2}\cos\theta$. Hence the potential at the point $P$ is
\begin{align*}
\Phi&=\Phi_+-\Phi_-
=\frac{p}{4\pi\epsilon_0}\left[\frac{\cos(\theta+\Delta\theta)}{\left(r-\frac{l}{2}\cos\theta\right)^2}
-\frac{\cos(\theta-\Delta\theta)}{\left(r+\frac{l}{2}\cos\theta\right)^2}\right]\\
&=\frac{p}{4\pi\epsilon_0r^2}\left[(\cos\theta-\sin\theta\,\Delta\theta)\left(1+\frac{l}{r}\cos\theta\right)
-(\cos\theta+\sin\theta\cdot\Delta\theta)\left(1-\frac{l}{r}\cos\theta\right)\right]\\
&=\frac{p}{4\pi\epsilon_0r^2}\left[-\sin\theta\cdot\Delta\theta+\frac{l}{r}\cos^2\theta
-\sin\theta\cdot\Delta\theta+\frac{l}{r}\cos^2\theta\right]\\
&=\frac{p}{4\pi\epsilon_0r^2}\left[\frac{2l}{r}\cos^2\theta-2\sin\theta\cdot\frac{l\sin\theta}{2r}\right]\\
&=\frac{ql^2}{4\pi\epsilon_0r^3}(3\cos^2\theta-1)
=\frac{Q}{8\pi\epsilon_0r^3}(3\cos^2\theta-1),
\end{align*}
where quantities of second order and higher have been neglected and use has been made of $\Delta\theta=\dfrac{l\sin\theta}{2r}$. The electric field strength is
\begin{align*}
E_r&=-\frac{\partial\Phi}{\partial r}=\frac{3ql^2}{4\pi\epsilon_0r^4}(3\cos^2\theta-1)
=\frac{3}{2}\frac{Q}{4\pi\epsilon_0r^4}(3\cos^2\theta-1),\\
E_\theta&=-\frac{1}{r}\frac{\partial\Phi}{\partial\theta}
=-\frac{1}{r}\frac{ql^2}{4\pi\epsilon_0r^3}(-6\cos\theta\sin\theta)
=\frac{6ql^2\sin\theta\cos\theta}{4\pi\epsilon_0r^4}
=\frac{3Q}{4\pi\epsilon_0r^4}\sin\theta\cos\theta.
\end{align*}
In the above expressions for the potential and the electric field, $Q=2ql^2$ is called the quadrupole moment of the electric quadrupole.

\begin{figbox}{3.02}
\begin{tikzpicture}[scale=0.95,>=latex]
  % point P
  \fill (6.2,3.0) circle (1.2pt);
  \node[right] at (6.25,3.05) {$P$};
  % quadrupole charges on the vertical axis
  \coordinate (A) at (0,0.75);
  \coordinate (O) at (0,0);
  \coordinate (B) at (0,-0.75);
  \draw[thickline] (A)--(B);
  \fill (A) circle (2.2pt);
  \fill (O) circle (2.2pt);
  \fill (B) circle (2.2pt);
  \node[left] at (0,0.75) {$+q$};
  \node[left] at (0,0) {$-q$};
  \node[left] at (0,-0.75) {$+q$};
  % dipole moment arrows
  \draw[force] (0,0.9)--(0,1.75) node[right] {$p_+$};
  \draw[force] (0,-0.9)--(0,-1.75) node[right] {$p_-$};
  % rays from P
  \draw[thickline] (6.2,3.0)--(A);
  \draw[thickline] (6.2,3.0)--(O);
  \draw[thickline] (6.2,3.0)--(B);
  \node[above] at (3.2,2.05) {$r_+$};
  \node[above] at (3.2,1.62) {$r$};
  \node[below right] at (3.2,1.05) {$r_-$};
  % angle marks at P
  \draw[thickline] (5.05,2.82) arc (178.5:185.4:1.2);
  \node at (4.85,2.72) {$\Delta\theta$};
  \draw[thickline] (5.05,2.68) arc (185.4:192.3:1.2);
  \node at (4.85,2.45) {$\Delta\theta$};
\end{tikzpicture}
\end{figbox}

\begin{commentary}
Part (1) of this problem is solved by the energy method, which is much more convenient than solving it directly by the composition of forces. If part (2) were solved by superposition of the fields of point charges, second-order small quantities $l^2/r^3$ would have to be retained in the approximation of the distances, and the calculation would be more involved. The concepts of the electric (di)pole and quadrupole involved in this problem have wide applications in electricity.
\end{commentary}
\end{solution}

\begin{solution}{3.03}
(1) Since $q'<q$, the lines of force issuing from $+q$ cannot all terminate on $-q'$; a part of them will extend to infinity. Let the lines of force with $\alpha>\alpha_0$ extend to infinity; then the lines of force issuing from $+q$ and lying inside the cone of half-angle $\alpha_0$ will all terminate on $-q'$. Since the number of lines of force is proportional to the electric flux, by Gauss's law, and noting that the solid angle contained in a cone of half-angle $\alpha_0$ is $2\pi(1-\cos\alpha_0)$, we have
\[
\frac{2\pi(1-\cos\alpha_0)}{4\pi}\,q=q'.
\]
Solving this gives
\[
\alpha_0=\arccos\left(1-2\frac{q'}{q}\right),
\]
or
\[
\alpha_0=2\arcsin\sqrt{\frac{q'}{q}}.
\]
Therefore the angle $\beta$ should be discussed in the two cases $\alpha<\alpha_0$ and $\alpha>\alpha_0$.

When $\alpha<\alpha_0$, in analogy with the discussion above, when the line of force terminates on $-q'$ the angle $\beta$ it makes with the line $AB$ satisfies
\[
\frac{2\pi(1-\cos\alpha)}{4\pi}\,q=\frac{2\pi(1-\cos\beta)}{4\pi}\,q'.
\]
From this we obtain
\[
\cos\beta=1-\frac{q}{q'}(1-\cos\alpha),
\]
hence
\[
\beta=\arccos\left[1-\frac{q}{q'}(1-\cos\alpha)\right],
\]
or
\[
\beta=2\arcsin\left(\sin\frac{\alpha}{2}\sqrt{\frac{q}{q'}}\right).
\]

When $\alpha>\alpha_0$, the lines of force extend to infinity. At this time the lines of force issuing from $+q$ and lying inside the cone of half-angle $(\pi-\alpha_0)$ will ultimately be distributed uniformly within the solid angle $4\pi$, the total number of lines of force being proportional to the charge $(q-q')$. That is to say, the lines of force issuing from $+q$ and lying within the solid angle $[4\pi-2\pi(1-\cos\alpha)]$ are ultimately as if emitted from a point charge of magnitude $(q-q')$, and lie within the solid angle $[4\pi-2\pi(1-\cos\beta)]$. Hence
\[
\frac{4\pi-2\pi(1-\cos\alpha)}{4\pi}\,q=\frac{4\pi-2\pi(1-\cos\beta)}{4\pi}(q-q'),
\]
that is,
\[
q\left[1-\frac{1}{2}(1-\cos\alpha)\right]=(q-q')\left[1-\frac{1}{2}(1-\cos\beta)\right].
\]
Solving this gives
\[
\cos\beta=\frac{q}{q-q'}(1+\cos\alpha)-1,
\]
or
\[
\cos\frac{\beta}{2}=\cos\frac{\alpha}{2}\sqrt{\frac{q}{q-q'}},
\]
hence
\[
\beta=\arccos\left[\frac{q}{q-q'}(1+\cos\alpha)-1\right],
\]
or
\[
\beta=2\arccos\left(\cos\frac{\alpha}{2}\sqrt{\frac{q}{q-q'}}\right).
\]

(2) For a line of force terminating on $-q'$, the point of intersection is just at $-q'$. For a line of force going to infinity, in order to find the point of intersection of its final tangent with the $x$-axis, consider the electric field at a distant point $P$. The electric field $\vec{E}_P$ at $P$ is the vector sum of the field $\vec{E}_q$ produced by the charge $q$ and the field $\vec{E}_{-q'}$ produced by $-q'$.

From the figure it is not difficult to see that the extension of $E_P$ must intersect the $x$-axis to the left of the point $A$. Let the point of intersection be $C$, and set $CA=x_1$, $CB=x_2$; write $AP=r_1$, $BP=r_2$, $\angle CPA=\theta$, $\angle CPB=\varphi$, $\angle PCA=\psi$. Applying the sine rule in $\triangle PCA$ and in $\triangle PCB$ respectively, we have
\[
\frac{x_1}{\sin\theta}=\frac{r_1}{\sin\psi},
\qquad
\frac{x_2}{\sin\varphi}=\frac{r_2}{\sin\psi}.
\]
Since $P$ is very far away, $r_1\approx r_2$, so from the above two equations we obtain
\[
\frac{x_1}{x_2}=\frac{\sin\theta}{\sin\varphi}.
\]
Moreover, from the vector diagram $\vec{E}_q+\vec{E}_{-q'}=\vec{E}_P$, and from $E_q\propto q/r_1^2$, $E_{-q'}\propto q'/r_2^2$, we have
\[
\frac{q'/r_2^2}{\sin\theta}=\frac{q/r_1^2}{\sin\varphi}.
\]
Since also $r_1\approx r_2$, the above equation becomes
\[
\frac{\sin\theta}{\sin\varphi}=\frac{q'}{q}.
\]
Comparing this with the previous equation, we obtain
\[
\frac{x_1}{x_2}=\frac{q'}{q}.
\]
If we put $AB=l$, it is not difficult to obtain from the above equation
\[
x_1=\frac{q'}{q-q'}l,
\qquad
x_2=\frac{q}{q-q'}l.
\]

\begin{figbox}{3.03}
\begin{tikzpicture}[scale=0.95,>=latex]
  % x-axis
  \draw[axis] (-1.2,0)--(8.6,0) node[right] {$x$};
  \node[below] at (-0.9,0) {$C$};
  \node[below left] at (2.15,0) {$A$};
  \node[below right] at (3.35,0) {$B$};
  % charges
  \fill (2.15,0) circle (1.6pt) node[above] {$q$};
  \fill (3.35,0) circle (1.6pt) node[above right] {$-q'$};
  % point P and rays
  \coordinate (P) at (7.4,3.15);
  \fill (P) circle (1.2pt) node[above right] {$P$};
  \draw[thickline] (-0.9,0)--(P);
  \draw[dashedline] (2.15,0)--(P);
  \draw[dashedline] (3.35,0)--(P);
  % field vectors at P
  \draw[force] (P)--(7.95,4.15) node[above right] {$E_q$};
  \draw[force] (P)--(8.35,3.32) node[right] {$E_{-q'}$};
  \draw[force] (P)--(8.02,3.78) node[above right] {$E_P$};
  % labels of distances and angles
  \node[below] at (0.6,0.05) {$x_1$};
  \node[below] at (2.75,0.02) {$x_2$};
  \node[above left] at (4.75,1.68) {$r_1$};
  \node[below right] at (5.4,1.45) {$r_2$};
  \draw[thickline] (200.8:1.05) arc (200.8:211:1.05);
  \node at (206:1.25) {$\theta$};
  \draw[thickline] (211:0.78) arc (211:217.9:0.78);
  \node at (214.5:0.98) {$\varphi$};
  \draw[thickline] (-0.3,0) arc (0:20.8:0.6);
  \node at (-0.12,0.16) {$\psi$};
\end{tikzpicture}
\end{figbox}

\begin{commentary}
(1) The key to solving this problem is to know that the number of lines of force is proportional to the electric flux and that the electric flux is proportional to the charge (this is Gauss's law).

(2) Comparing the expressions for $x_1$ and $x_2$ with the expression for the centre-of-mass coordinate, it is not difficult to find that, provided one takes into account that a charge has not only magnitude but also sign, these two expressions are no different from the expression for the $x$-coordinate of the centre of mass. Hence it may be said that, viewed from afar, the lines of force of a system of charges appear to issue from the ``centre of gravity'' of the system of charges. This conclusion is applicable to any system of charges. As with the centre of mass, the position of the ``centre of gravity'' $C$ of a system of charges is determined by
\[
\vec{r}_C=\frac{\sum q_i\vec{r}_i}{\sum q_i},
\]
where $\vec{r}_i$ is the position vector of the point charge $q_i$. Note that $q_i$ may be positive or negative. The above expression can also be extended to continuously distributed systems of charges, provided that the sum is replaced by an integral. This is precisely the physical basis of the assertion that any static system of charges of finite extent can be regarded as a point charge at a sufficiently great distance.
\end{commentary}
\end{solution}

\begin{solution}{3.04}
First consider the potential produced at the point $P$ by the ring carrying $+q$. Take the vertical plane containing $OP$ and the $z$-axis as the $xz$-plane and set up coordinates as shown in the figure. Consider the potential produced at $P$ by the charge $\eta a\,\dd\varphi$ of the arc element $\varphi\to\varphi+\dd\varphi$ on the ring, where $\eta=q/2\pi a$ is the linear charge density. The distance $AP$ from the position $A$ of this charge element to the point $P$ is not difficult to obtain (see the figure):
\begin{align*}
\overline{AP}^2&=\overline{PP'}^2+\overline{AP'}^2\\
&=(r\cos\theta)^2+a^2+(r\sin\theta)^2+2a\cdot r\sin\theta\cos\varphi\\
&=r^2+a^2+2ra\sin\theta\cos\varphi.
\end{align*}
The potential produced by this charge element at the point $P$ (to accuracy $a^2/r^2$) is
\begin{align*}
\dd\Phi_+&=\frac{1}{4\pi\epsilon_0}\frac{\eta a\,\dd\varphi}{\sqrt{r^2+a^2+2ra\sin\theta\cos\varphi}}\\
&=\frac{\eta a}{4\pi\epsilon_0r}\frac{\dd\varphi}{\sqrt{1+\dfrac{a^2}{r^2}+2\dfrac{a}{r}\sin\theta\cos\varphi}}\\
&\approx\frac{\eta a\,\dd\varphi}{4\pi\epsilon_0r}
\left(1-\frac{1}{2}\frac{a^2}{r^2}-\frac{a}{r}\sin\theta\cos\varphi
+\frac{3}{2}\frac{a^2}{r^2}\sin^2\theta\cos^2\varphi\right).
\end{align*}
The potential produced by the whole ring carrying $+q$ at the point $P$ is the integral of the above expression:
\begin{align*}
\Phi_+&=\int\dd\Phi_+
=\frac{2\pi\eta a}{4\pi\epsilon_0r}\int_0^\pi
\left(1-\frac{1}{2}\frac{a^2}{r^2}-\frac{a}{r}\sin\theta\cos\varphi
+\frac{3}{2}\frac{a^2}{r^2}\sin^2\theta\cos^2\varphi\right)\dd\varphi\\
&=\frac{2\pi\eta a}{4\pi\epsilon_0r}
\left[\left(1-\frac{1}{2}\frac{a^2}{r^2}\right)+\frac{3}{4}\frac{a^2}{r^2}\sin^2\theta\right]\pi\\
&=\frac{q}{4\pi\epsilon_0r}\left[1-\frac{1}{2}\frac{a^2}{r^2}
\left(1-\frac{3}{2}\sin^2\theta\right)\right].
\end{align*}
Similarly, the potential produced at the point $P$ by the ring carrying $-q$ is
\[
\Phi_-=-\frac{q}{4\pi\epsilon_0r}\left[1-\frac{1}{2}\frac{b^2}{r^2}
\left(1-\frac{3}{2}\sin^2\theta\right)\right].
\]
The total potential at the point $P$ is
\begin{align*}
\Phi&=\Phi_++\Phi_-
=\frac{q}{4\pi\epsilon_0r}\left[\frac{1}{2r^2}(b^2-a^2)
-\frac{3}{4r^2}\sin^2\theta\,(b^2-a^2)\right]\\
&=\frac{q(b^2-a^2)}{4\pi\epsilon_0\cdot 2r^3}
\left(1-\frac{3}{2}\sin^2\theta\right)\\
&=\frac{q(b^2-a^2)}{4\pi\epsilon_0\cdot 4r^3}(3\cos^2\theta-1).
\end{align*}
The electric field strength is
\begin{align*}
E_r&=-\frac{\partial\Phi}{\partial r}
=\frac{3q(b^2-a^2)}{16\pi\epsilon_0r^4}(3\cos^2\theta-1),\\
E_\theta&=-\frac{1}{r}\frac{\partial\Phi}{\partial\theta}
=\frac{6q(b^2-a^2)}{16\pi\epsilon_0r^4}\cos\theta\sin\theta.
\end{align*}

\begin{figbox}{3.04}
\begin{tikzpicture}[scale=0.95,>=latex]
  % ring in perspective
  \draw[thickline] (0,0) ellipse (1.6 and 0.55);
  % axes
  \draw[axis] (0,0)--(-3.1,0) node[left] {$x$};
  \draw[axis] (0,0)--(1.35,-1.15) node[below right] {$y$};
  \draw[axis] (0,0)--(0,3.3) node[above] {$z$};
  \node[below right] at (0.06,-0.06) {$O$};
  % point P and its projection
  \coordinate (P) at (-3.3,2.3);
  \coordinate (Q) at (-3.3,0);
  \draw[dashedline] (P)--(Q);
  \node[below] at (-3.3,0) {$P'$};
  \fill (P) circle (1.2pt);
  \node[above left] at (P) {$P$};
  \draw[thickline] (0,0)--(P);
  \node[above right] at (-1.9,1.35) {$r$};
  % angle theta at O
  \draw[thickline] (0,1.25) arc (90:145:1.25);
  \node at (-0.45,1.42) {$\theta$};
  % charge element A
  \coordinate (A) at (-1.52,0.17);
  \fill (A) circle (1.6pt);
  \node[below left] at (A) {$A$};
  \draw[thickline] (A)--(P);
  \draw[thickline] (0,0)--(A);
  \node[above] at (-0.8,0.28) {$a$};
  % angle phi and element dphi
  \draw[thickline] (-0.95,0) arc (180:174:0.95);
  \node at (-0.98,0.42) {$\varphi$};
  \draw[thickline] (-1.15,0.35) arc (162:170:1.18);
  \node at (-1.38,0.50) {$\dd\varphi$};
\end{tikzpicture}
\end{figbox}

\begin{commentary}
Comparing with Problem 3.02(2), it is seen that the electric field of two concentric rings carrying equal and opposite charges is similar to that of an electric quadrupole, whose quadrupole moment is $Q=\dfrac{q}{2}(b^2-a^2)$.
\end{commentary}
\end{solution}

\begin{solution}{3.05}
(1) Let the sphere have radius $R$ and total charge $Q$; then the potential at a point on the spherical surface is
\[
\varphi_s=\frac{1}{4\pi\epsilon_0}\frac{Q}{R}.
\]
The electric field at any point inside the sphere can be found from Gauss's law. For a point at distance $r$ from the centre,
\[
E(r)=\frac{1}{4\pi\epsilon_0}\frac{r^2}{R^3}\frac{Q}{r^2}
=\frac{1}{4\pi\epsilon_0}\frac{Q}{R^3}r.
\]
Hence the potential at the centre of the sphere is
\begin{align*}
\varphi_0&=\int_0^R E\,\dd r+\varphi_s
=\frac{1}{4\pi\epsilon_0}\frac{Q}{R^3}\int_0^R r\,\dd r+\varphi_s\\
&=\frac{1}{4\pi\epsilon_0}\frac{Q}{2R}+\frac{1}{4\pi\epsilon_0}\frac{Q}{R}
=\frac{3}{2}\cdot\frac{1}{4\pi\epsilon_0}\frac{Q}{R}.
\end{align*}
Therefore the required ratio of the potentials is
\[
\varphi_s/\varphi_0=\frac{2}{3}.
\]

(2) Let the cube have edge length $l$ and total charge $Q$, with charge density $\rho$. Imagine cutting the cube into eight small cubes of edge $l/2$, as shown in the figure; then the potential at the centre $O$ is equal to the sum of the potentials at the vertices of the eight small cubes:
\[
\varphi_O^{l}=8\varphi^{l/2}.
\]
But the potential at a vertex is obviously proportional to the total charge and inversely proportional to the edge length, hence
\[
\varphi\propto\frac{Q}{l}=\frac{\rho l^3}{l}=\rho l^2.
\]
With $\rho$ unchanged, the potential is therefore proportional to the square of the edge length, so that $\varphi^{l/2}=\dfrac{1}{4}\varphi^{l}$. Substituting into the original equation,
\[
\varphi_O^{l}=8\varphi^{l/2}=8\cdot\frac{1}{4}\varphi^{l}=2\varphi^{l}.
\]
Hence the required ratio of the potentials is
\[
\varphi^{l}/\varphi_O^{l}=\frac{1}{2}.
\]

\begin{figbox}{3.05}
\begin{tikzpicture}[scale=0.95,>=latex]
  \def\s{2.2}
  \def\dx{0.85}
  \def\dy{0.7}
  % front face
  \draw[thickline] (0,0) rectangle (\s,\s);
  % back face
  \draw[thickline] (\dx,\dy)--(\dx+\s,\dy)--(\dx+\s,\dy+\s)--(\dx,\dy+\s)--cycle;
  % connecting edges
  \draw[thickline] (0,0)--(\dx,\dy);
  \draw[thickline] (\s,0)--(\dx+\s,\dy);
  \draw[thickline] (\s,\s)--(\dx+\s,\dy+\s);
  \draw[thickline] (0,\s)--(\dx,\dy+\s);
  % division lines on the front face
  \draw[thickline] (\s/2,0)--(\s/2,\s);
  \draw[thickline] (0,\s/2)--(\s,\s/2);
  % division lines on the top face
  \draw[thickline] (\s/2,\s)--(\dx+\s/2,\dy+\s);
  \draw[thickline] (\dx,\dy+\s/2)--(\dx+\s,\dy+\s/2);
  % division lines on the right face
  \draw[thickline] (\s,\s/2)--(\dx+\s,\dy+\s/2);
  \draw[thickline] (\s/2,0)--(\dx+\s/2,\dy);
  % centre O and dashed construction lines
  \coordinate (O) at (\s/2+\dx/2,\s/2+\dy/2);
  \fill (O) circle (1.2pt);
  \node[right] at (O) {$O$};
  \draw[dashedline] (O)--(\s/2,\s/2);
  \draw[dashedline] (O)--(\s/2+\dx,\s/2+\dy);
\end{tikzpicture}
\end{figbox}
\end{solution}

\begin{solution}{3.06}
The magnitude of the electric field at any point outside the spherical shell can be obtained from the electric field at the surface of the shell:
\[
E=\frac{1}{4\pi\epsilon_0}\frac{q}{a^2}.
\]
To find the charge density $\rho(r)$ at $r$, take two concentric spherical surfaces of radii $r$ and $r+\Delta r$ as Gaussian surfaces. By Gauss's law,
\[
E\cdot 4\pi(r+\Delta r)^2-E\cdot 4\pi r^2=\frac{1}{\epsilon_0}\rho\cdot 4\pi r^2\Delta r.
\]
From this we obtain
\[
\rho=\rho(r)=\frac{2\epsilon_0E}{r}.
\]
Substituting the expression for $E$, we obtain
\[
\rho(r)=\frac{q}{2\pi a^2r}.
\]
\end{solution}

\begin{solution}{3.07}
For simplicity, first assume that the positive ion layer is at rest and that the negative ion layer is displaced to the right by $x$, with $x\ll l$ (see the figure). To find the electrostatic force exerted on the negative ion layer by the positive ion layer, divide it into two parts: part I, which coincides with the positive ion layer, and part II, which does not coincide. Part I is acted on only by the thin positive ion layer of thickness $x$ that does not overlap the negative ion layer, while part II is acted on by the whole positive ion layer. A positive ion layer of thickness $x$ produces a uniform electric field on both sides, directed perpendicular to the plane and away from the thin layer on either side; its magnitude is easily found from Gauss's law:
\[
E_1=\frac{ne}{2\epsilon_0}x.
\]
Hence the force per unit area on the negative ions of part I is
\[
F_{\mathrm{I}}=-ne(l-x)\frac{ne}{2\epsilon_0}x\approx-\frac{n^2e^2}{2\epsilon_0}lx.
\]
Here we have neglected the second-order small quantity $x^2$. The whole positive ion layer also produces a uniform electric field on both sides, and its strength can likewise be found from Gauss's law:
\[
E_2=\frac{ne}{2\epsilon_0}l.
\]
Thus the force per unit area on the negative ions of part II is
\[
F_{\mathrm{II}}=-nex\cdot\frac{ne}{2\epsilon_0}l=-\frac{n^2e^2}{2\epsilon_0}lx.
\]
The resultant force per unit area on the negative ion layer is
\[
F=F_{\mathrm{I}}+F_{\mathrm{II}}=-\frac{n^2e^2}{\epsilon_0}lx.
\]
The force is proportional to the displacement $x$ and opposite to $x$ in direction; it is a simple harmonic force, so the negative ion layer performs simple harmonic motion with an angular frequency equal to the square root of the ratio of the force to the displacement divided by the mass:
\[
\omega'=\sqrt{\frac{n^2e^2l}{\epsilon_0m_-\cdot nl}}
=\sqrt{\frac{ne^2}{\epsilon_0m_-}}.
\]
Taking the motion of the positive ions into account, it is only necessary to replace $m_-$ in the above expression by the reduced mass $m_+m_-/(m_++m_-)$ to obtain the actual angular frequency
\[
\omega=\sqrt{\frac{ne^2(m_++m_-)}{\epsilon_0m_+m_-}}.
\]

\begin{figbox}{3.07}
\begin{tikzpicture}[scale=0.9,>=latex]
  % positive ion layer
  \draw[thickline] (0,0) rectangle (1.9,3.8);
  \foreach \y in {0.3,0.75,...,3.6} {
    \foreach \xx in {0.28,0.7,1.12,1.54} {
      \node[font=\scriptsize] at (\xx,\y) {$+$};
    }
  }
  % negative ion layer (displaced)
  \draw[thickline] (0.4,0) rectangle (2.35,3.8);
  \foreach \y in {0.3,0.75,...,3.6} {
    \foreach \xx in {0.62,1.04,1.46,1.88,2.2} {
      \node[font=\scriptsize] at (\xx,\y) {$-$};
    }
  }
  % dimensions
  \draw[dim] (0,4.15)--(1.9,4.15) node[midway,above] {$l$};
  \draw[dim] (1.9,4.15)--(2.35,4.15) node[midway,above] {$x$};
  \draw[dashedline] (1.9,3.8)--(1.9,4.3);
  \draw[dashedline] (2.35,3.8)--(2.35,4.3);
  \draw[dashedline] (0,3.8)--(0,4.3);
  % part labels
  \node[below] at (1.15,-0.15) {I};
  \node[below] at (2.12,-0.15) {II};
\end{tikzpicture}
\end{figbox}
\end{solution}

\begin{solution}{3.08}
(1) Now the potential at a distance $r$ from the point charge $q$ (still taking infinity as the zero of potential) is
\[
\varphi=\int_r^\infty\frac{1}{4\pi\epsilon_0}\frac{q}{r^{2+\delta}}\,\dd r
=\frac{q}{4\pi\epsilon_0(1+\delta)}\frac{1}{r^{1+\delta}}.
\]
Now consider the potential at a point $P$ at distance $r$ from the centre of the sphere. Take the line joining the centre $O$ of the sphere to $P$ as axis and construct cones of half-angle $\theta$ and $\theta+\dd\theta$, which cut out a ring band on the spherical surface. This band is at the same distance from $P$, denoted $r'$, so the potential produced at $P$ by the charge on this band is
\[
\dd\varphi=\frac{1}{4\pi\epsilon_0}\frac{\sigma\cdot 2\pi R\sin\theta\,R\,\dd\theta}{(1+\delta)r'^{(1+\delta)}}.
\]

\begin{figbox}{3.08}
\begin{tikzpicture}[scale=0.95,>=latex]
  % spherical shell
  \draw[thickline] (0,0) circle (1.6);
  \node[below left] at (-0.1,-0.12) {$O$};
  % point P and axis
  \coordinate (P) at (5.4,0);
  \fill (P) circle (1.2pt);
  \node[right] at (5.45,0) {$P$};
  \draw[dashedline] (0,0)--(5.4,0);
  \node[below] at (2.9,0) {$r$};
  % cone cutting the ring band
  \coordinate (A1) at (28:1.6);
  \coordinate (A2) at (36:1.6);
  \draw[thickline] (0,0)--(A1);
  \draw[thickline] (0,0)--(A2);
  \draw[line width=1.6pt] (A1) arc (28:36:1.6);
  \draw[thickline] (P)--(32:1.6);
  \node[above right] at (3.35,1.05) {$r'$};
  % angles at O
  \draw[thickline] (1.0,0) arc (0:28:1.0);
  \node at (1.12,0.30) {$\theta$};
  \draw[thickline] (28:1.22) arc (28:36:1.22);
  \node at (0.98,0.80) {$\dd\theta$};
\end{tikzpicture}
\end{figbox}

Since $r'^2=R^2+r^2-2Rr\cos\theta$, this may be written as
\[
\dd\varphi=\frac{\sigma R^2}{2\epsilon_0(1+\delta)}
\frac{\sin\theta\,\dd\theta}{(R^2+r^2-2Rr\cos\theta)^{(1+\delta)/2}}.
\]
The potential at the point $P$ is the integral of the above expression over $\theta$:
\begin{align*}
\varphi&=\frac{\sigma R^2}{2\epsilon_0(1+\delta)}\int_0^\pi
\frac{\sin\theta\,\dd\theta}{(R^2+r^2-2Rr\cos\theta)^{(1+\delta)/2}}\\
&=\frac{\sigma R}{2\epsilon_0(1-\delta^2)r}\bigl[(R+r)^{1-\delta}-|R-r|^{1-\delta}\bigr].
\end{align*}
When $r<R$ (the point $P$ is inside the sphere), $|R-r|=R-r$; when $r>R$ (the point $P$ is outside the sphere), $|R-r|=r-R$. Hence
\[
\varphi=
\begin{cases}
\dfrac{\sigma R}{2\epsilon_0(1-\delta^2)r}\bigl[(R+r)^{1-\delta}-(R-r)^{1-\delta}\bigr], & r<R\ (\text{inside the sphere}),\\[10pt]
\dfrac{\sigma R}{2\epsilon_0(1-\delta^2)r}\bigl[(R+r)^{1-\delta}-(r-R)^{1-\delta}\bigr], & r>R\ (\text{outside the sphere}).
\end{cases}
\]
(2) Since $\delta\ll 1$,
\begin{align*}
(R+r)^{1-\delta}&\approx(R+r)-\delta(R+r)\ln(R+r),\\
|R-r|^{1-\delta}&\approx|R-r|-\delta|R-r|\ln|R-r|,\\
1-\delta^2&\approx 1.
\end{align*}
The potential may then be written as
\[
\varphi\approx
\begin{cases}
\dfrac{\sigma R}{2\epsilon_0r}\bigl\{2r+\delta\bigl[(R-r)\ln(R-r)-(R+r)\ln(R+r)\bigr]\bigr\}, & r<R,\\[10pt]
\dfrac{\sigma R}{2\epsilon_0r}\bigl\{2R+\delta\bigl[(r-R)\ln(r-R)-(R+r)\ln(R+r)\bigr]\bigr\}, & r>R,\\[10pt]
\dfrac{\sigma}{2\epsilon_0}\bigl[2R-\delta\cdot 2R\ln 2R\bigr], & r=R.
\end{cases}
\]
Applying this result to the two spherical shells in the problem, and denoting their surface charge densities by $\sigma_a$ and $\sigma_b$ ($\sigma_b\ll\sigma_a$), we have, retaining first-order small quantities,
\begin{align*}
\varphi_a&=\frac{\sigma_a}{2\epsilon_0}\bigl[2a-\delta\cdot 2a\ln 2a\bigr]
+\frac{b\sigma_b}{2\epsilon_0a}\bigl\{2b+\delta\bigl[(a-b)\ln(a-b)-(a+b)\ln(a+b)\bigr]\bigr\}\\
&\approx\frac{\sigma_aa}{\epsilon_0}-\delta\frac{\sigma_aa}{\epsilon_0}\ln 2a+\frac{\sigma_bb^2}{\epsilon_0a},\\
\varphi_b&=\frac{\sigma_b}{2\epsilon_0}\bigl[2b-\delta\cdot 2b\ln 2b\bigr]
+\frac{a\sigma_a}{2\epsilon_0b}\bigl\{2b+\delta\bigl[(a-b)\ln(a-b)-(a+b)\ln(a+b)\bigr]\bigr\}\\
&\approx\frac{\sigma_bb}{\epsilon_0}+\frac{\sigma_aa}{\epsilon_0}
+\delta\frac{\sigma_aa}{2\epsilon_0b}\bigl[(a-b)\ln(a-b)-(a+b)\ln(a+b)\bigr].
\end{align*}
Putting $\varphi_a=\varphi_b=U$, we obtain
\[
\frac{b\sigma_b}{\epsilon_0}\left(1-\frac{b}{a}\right)
=-\delta\left\{\frac{\sigma_aa}{\epsilon_0}\ln 2a
+\frac{\sigma_aa}{2\epsilon_0b}\bigl[(a-b)\ln(a-b)-(a+b)\ln(a+b)\bigr]\right\}.
\]
But $\dfrac{\sigma_aa}{\epsilon_0}\approx U$; by the statement of the problem, $\dfrac{b\sigma_b}{\epsilon_0}\left(1-\dfrac{b}{a}\right)\le u$. Substituting into the above equation gives
\[
u\ge\frac{1}{2}\delta U\left\{\frac{1}{b}\bigl[(a+b)\ln(a+b)-(a-b)\ln(a-b)\bigr]-2\ln 2a\right\}
\]
\[
=\frac{1}{2}\delta U\left(\ln\frac{a^2-b^2}{4a^2}+\frac{a}{b}\ln\frac{a+b}{a-b}\right).
\]
Therefore
\[
\delta\le\frac{2u}{U\left(\ln\dfrac{a^2-b^2}{4a^2}+\dfrac{a}{b}\ln\dfrac{a+b}{a-b}\right)}.
\]

\begin{commentary}
The inverse-square law of the electrostatic force is of profound significance, and experimental tests of the inverse-square law have always been taken seriously by the scientific community. This problem is the theoretical basis of the experimental test method used by Maxwell in 1872. The experimental result gave a value of $u$ equal to zero; later, from the sensitivity of the apparatus and the relevant data, it was estimated that $\delta<5\times10^{-5}$. The result obtained by more precise experimental methods in modern times (the 1970s) is $2\times10^{-16}$.
\end{commentary}
\end{solution}

\begin{solution}{3.09}
Although the geometrical structure of this apparatus is very symmetric, the uncertainty of such factors as the distribution of ions in the air and cosmic rays makes the potentials of the aluminium cylinders $A_1$ and $A_2$ slightly different. For example, if the potential of $A_1$ is higher than that of $A_2$, then by electrostatic induction the water drops above $A_1$ become negatively charged and the water drops above $A_2$ become positively charged. After the charged water drops fall into the aluminium cups below, the cup $B_1$ becomes negatively charged. Since $B_1$ and $A_2$ are connected by a wire, the potential of $A_2$ is further lowered, and in the same way the potential of $A_1$ is further raised; this in turn makes the drops above $A_1$ carry still more negative charge and the drops above $A_2$ carry still more positive charge. Continuing in this way, the potential of the aluminium cup $B_2$ becomes higher and higher and that of $B_1$ lower and lower, until finally a potential difference of several thousand volts can be produced between the two aluminium cups. Of course, because of the uncertainty of the various factors, at the beginning of the next experiment the potential of $A_2$ may be higher than that of $A_1$, and finally the potential of $B_1$ will be several thousand volts higher than that of $B_2$. But it does not change that the potential difference between $A_1$ and $A_2$ produced by accidental factors is amplified by the positive-feedback effect described above.
\end{solution}

\begin{solution}{3.10}
(1) \textit{Capacitive type.}\quad From the formula $u=q/C$, in order to convert the sound vibration (which causes a change of capacitance) into an output electric signal, $q$ should be kept essentially constant; hence the circuit shown in Fig.~3.10(a) may be used, with $R\gg\dfrac{1}{C\omega}$. For this circuit we have the following equations:
\begin{align*}
u&=\frac{q}{C}+\frac{\dd q}{\dd t}R\\
&=\frac{q}{\epsilon_0S}(d+x)+\frac{\dd q}{\dd t}R\\
&=u_1(t)+u_2(t).
\end{align*}
When $R\gg\dfrac{1}{C\omega}$, $q$ is essentially constant and $u_2(t)$ is not easy to find, but $u_1(t)$ can be found:
\begin{align*}
u_1(t)&=\frac{q}{\epsilon_0S}(d+x)=\frac{q}{\epsilon_0S}d+\frac{q}{\epsilon_0S}x=u_{10}+\Delta u_1(t),\\
\Delta u_1(t)&=\frac{q}{\epsilon_0S}x=\frac{q}{\epsilon_0S/d}\cdot\frac{x}{d}=\frac{q}{C}\frac{x}{d}=u\frac{x}{d}\\
&=u\cdot\frac{1}{d}x_0\sin\omega t.
\end{align*}
This is the output signal. If $R$ is very small $\left(R\ll\dfrac{1}{C\omega}\right)$, $u_1(t)$ is not easy to find (because $q$ varies greatly), but $u_2(t)$ can be found:
\begin{align*}
u_2(t)&=\frac{\dd q}{\dd t}R=R\cdot\frac{\dd(uC)}{\dd t}=Ru\frac{\dd C}{\dd t}
=uR\frac{\epsilon_0S}{d^2}\frac{\dd x}{\dd t}\\
&\approx uRC\frac{1}{d}\omega x=uRC\omega\frac{x}{d}\ll u\frac{x}{d}.
\end{align*}
It is thus seen that the output signal is much smaller than that for $R\gg\dfrac{1}{C\omega}$.

\begin{figbox}{3.10(a)}
\begin{tikzpicture}[scale=0.92,>=latex]
  % top-left branch with capacitor
  \draw[thickline] (0,0)--(0,3.0)--(1.35,3.0);
  \draw[thickline] (1.35,2.45)--(1.35,3.55);
  \draw[thickline] (2.05,2.45)--(2.05,3.55);
  \draw[thickline] (2.05,3.0)--(6.4,3.0);
  % labels of capacitor
  \node[above] at (1.7,3.6) {$C$};
  \draw[dim] (1.35,2.72)--(2.05,2.72) node[midway,below] {$d$};
  \draw[force] (2.05,3.75)--(2.35,3.75) node[right] {$x_0\sin\omega t$};
  % right branch with resistor
  \draw[thickline] (6.4,3.0)--(6.4,2.05);
  \draw[thickline] (6.15,2.05) rectangle (6.65,1.25);
  \node[right] at (6.75,1.65) {$R$};
  \draw[thickline] (6.4,1.25)--(6.4,0);
  % bottom branch
  \draw[thickline] (6.4,0)--(4.05,0);
  \draw[thickline] (3.35,0)--(0,0);
  % source u
  \draw[thickline] (3.95,0.32)--(3.95,-0.32);
  \draw[thickline] (3.45,0.18)--(3.45,-0.18);
  \node[below] at (3.7,-0.42) {$u$};
  % output terminal A
  \fill (6.4,3.0) circle (1.5pt);
  \node[above right] at (6.45,3.05) {$A$ (output)};
  % node B and earth
  \fill (0,0) circle (1.5pt);
  \node[left] at (-0.05,0.1) {$B$};
  \draw[thickline] (-0.35,0)--(-0.35,-0.3);
  \draw[thickline] (-0.55,-0.3)--(-0.15,-0.3);
  \draw[thickline] (-0.47,-0.42)--(-0.23,-0.42);
  \draw[thickline] (-0.4,-0.54)--(-0.3,-0.54);
\end{tikzpicture}
\end{figbox}

(2) \textit{Electret type.}\quad Since there is a constant applied charge, no external power supply is needed. The circuit shown in Fig.~3.10(b) may be used, with $R\gg\dfrac{1}{C\omega}$ (for the same reason as above). For this circuit, in the static state,
\begin{align}
\frac{\sigma_1}{\epsilon_r\epsilon_0}d_1&=\frac{\sigma_2}{\epsilon_0}d \tag{1}\\
\sigma_1+\sigma_2&=\sigma \tag{2}
\end{align}
From (1) and (2) we obtain
\[
\sigma_2=\frac{\sigma d_1}{\epsilon_rd+d_1}.
\]
Since $R$ is very large, this charge remains essentially unchanged when $d$ undergoes a small change. Hence the output signal is
\[
\Delta u(t)=\frac{\sigma_2}{\epsilon_0}x
=\frac{d_1\sigma}{\epsilon_0(\epsilon_rd+d_1)}x_0\sin\omega t.
\]

\begin{figbox}{3.10(b)}
\begin{tikzpicture}[scale=0.92,>=latex]
  % left branch
  \draw[thickline] (0,0)--(0,3.0)--(1.0,3.0);
  \draw[thickline] (1.0,2.45)--(1.0,3.55);
  % dielectric (electret) layer
  \fill[pattern=north east lines,pattern color=black!35] (1.05,2.5) rectangle (1.55,3.5);
  \draw[thickline] (1.05,2.5) rectangle (1.55,3.5);
  \node[above] at (1.3,3.55) {$\sigma_1$};
  \draw[dim] (1.05,2.7)--(1.55,2.7) node[midway,below] {$d_1$};
  % movable plate
  \draw[thickline] (2.3,2.45)--(2.3,3.55);
  \node[above] at (2.3,3.55) {$\sigma_2$};
  \draw[dim] (1.55,3.2)--(2.3,3.2) node[midway,above] {$d$};
  \draw[force] (2.3,3.85)--(2.6,3.85) node[right] {$x_0\sin\omega t$};
  \draw[thickline] (2.3,3.0)--(6.4,3.0);
  % right branch with resistor
  \draw[thickline] (6.4,3.0)--(6.4,2.05);
  \draw[thickline] (6.15,2.05) rectangle (6.65,1.25);
  \node[right] at (6.75,1.55) {$R\gg\dfrac{1}{C\omega}$};
  \draw[thickline] (6.4,1.25)--(6.4,0);
  % bottom branch
  \draw[thickline] (6.4,0)--(0,0);
  % output terminal
  \fill (6.4,3.0) circle (1.5pt);
  \node[above right] at (6.45,3.05) {$A$ (output)};
  % node B and earth
  \fill (0,0) circle (1.5pt);
  \node[left] at (-0.05,0.1) {$B$};
  \draw[thickline] (-0.35,0)--(-0.35,-0.3);
  \draw[thickline] (-0.55,-0.3)--(-0.15,-0.3);
  \draw[thickline] (-0.47,-0.42)--(-0.23,-0.42);
  \draw[thickline] (-0.4,-0.54)--(-0.3,-0.54);
\end{tikzpicture}
\end{figbox}

\begin{commentary}
The key to both parts of this problem is to keep the charge on the moving plate essentially unchanged, so that the field strength near the moving plate is also essentially unchanged; in this way the displacement of the moving plate gives rise to a voltage change to the greatest extent.
\end{commentary}
\end{solution}

\begin{solution}{3.11}
(1) The electric field in the region where the point charge $q$ is located (the first quadrant) can be regarded as the superposition of the fields of the point charge $q$ and of three other image charges $-q$, $-q$, $q$, as shown in the figure, while in the remaining regions (the second, third and fourth quadrants) the electric field is zero. Therefore the work done by the external force in moving the point charge $q$ from infinity to its present position is equal to one quarter of the work done by the external force in moving these four charges from infinity to their present positions (this is not difficult to understand if one imagines that during the process of moving, the positions of the four charges always remain symmetric), and the work done by the external force in moving the four charges from infinity to their present positions is equal to the interaction energy of the four charges. Hence the required work $W_{\mathrm{in}}$ is
\begin{align*}
W_{\mathrm{in}}&=\frac{1}{4}\left[4\left(\frac{1}{4\pi\epsilon_0}\right)\frac{-q^2}{2d}
+2\left(\frac{1}{4\pi\epsilon_0}\right)\frac{q^2}{2\sqrt{2}d}\right]\\
&=\frac{1}{4\pi\epsilon_0}\frac{q^2}{d}\left(\frac{1}{4\sqrt{2}}-\frac{1}{2}\right).
\end{align*}
This work is negative, which is obvious.

(2) When the point charge $q$ is moved, the force acting on it is equivalent to the force exerted on it by the three image charges; hence the required work $W_{\mathrm{out}}$ is the negative of the interaction energy of this point charge with the three image charges in the present configuration:
\begin{align*}
W_{\mathrm{out}}&=-\left[2\cdot\frac{1}{4\pi\epsilon_0}\frac{-q^2}{2d}
+\frac{1}{4\pi\epsilon_0}\frac{q^2}{2\sqrt{2}d}\right]\\
&=\frac{1}{4\pi\epsilon_0}\frac{q^2}{d}\left(1-\frac{1}{2\sqrt{2}}\right).
\end{align*}
Since the charge is always attracted during the displacement, this work is positive.

(3) In effect, $W_{\mathrm{in}}$ is equivalent to establishing the system of induced charges, that is, the system of charges on the insulating plate, and then moving the point charge $q$ from infinity to its present position, while $W_{\mathrm{out}}$ is the work done by the external force in moving the point charge $q$ from its present position to infinity in the electric field of the charge system on the insulating plate. Hence the work done by the external force in establishing the charge system on the insulating plate, that is, the electrostatic energy $W$ of the charge system on the insulating plate, is equal to the sum of $W_{\mathrm{in}}$ and $W_{\mathrm{out}}$:
\begin{align*}
W&=W_{\mathrm{in}}+W_{\mathrm{out}}
=\frac{1}{4\pi\epsilon_0}\frac{q^2}{d}
\left(\frac{1}{4\sqrt{2}}-\frac{1}{2}+1-\frac{1}{2\sqrt{2}}\right)\\
&=\frac{1}{4\pi\epsilon_0}\frac{q^2}{d}
\left(\frac{1}{2}-\frac{1}{4\sqrt{2}}\right).
\end{align*}
This energy is obviously positive.

\begin{figbox}{3.11}
\begin{tikzpicture}[scale=0.95,>=latex]
  % conducting planes
  \draw[dashedline] (-3.0,0)--(3.4,0);
  \draw[dashedline] (0,-3.0)--(0,3.0);
  \node[below right] at (0.05,-0.05) {$O$};
  % charges
  \fill (1.5,1.5) circle (1.7pt);
  \node[above right] at (1.55,1.55) {$q$};
  \fill (-1.5,1.5) circle (1.7pt);
  \node[above left] at (-1.55,1.55) {$-q$};
  \fill (-1.5,-1.5) circle (1.7pt);
  \node[below left] at (-1.55,-1.55) {$q$};
  \fill (1.5,-1.5) circle (1.7pt);
  \node[below right] at (1.55,-1.55) {$-q$};
  % dimensions
  \draw[dim] (0,1.5)--(1.5,1.5) node[midway,above] {$d$};
  \draw[dim] (1.5,0)--(1.5,1.5) node[midway,right] {$d$};
\end{tikzpicture}
\end{figbox}

\begin{commentary}
The method of images is one of the important methods for solving electrostatic problems (especially those involving conductors). This problem is an example of the method of images.
\end{commentary}
\end{solution}

\begin{solution}{3.12}
The point charge $q$ will induce on the inner surface of the spherical shell a charge of total amount $-q$ distributed non-uniformly, as shown in the figure, while on the outer surface of the shell a charge $q$ will be distributed uniformly. Since the charge distributed uniformly on the outer surface produces no electric field in its interior, the point charge $q$ inside the shell is acted on only by the force of the electric field produced by the induced charge on the inner surface of the shell, and the induced charge on the inner surface of the shell is equivalent to that on a grounded shell. But the electric field produced inside the sphere by the induced charge on the inner surface can be replaced by the field of an image charge outside the sphere.

To find the magnitude and position of this image charge $q'$, it suffices to examine the potentials at the two points $A$ and $B$ where the line joining $q$ to the centre $O$ of the shell intersects the inner surface of the shell. By symmetry, $q'$ must lie on this line; let its distance from the centre $O$ be $r'$. From the vanishing of the potentials at the points $A$ and $B$ we obtain the following equations:
\begin{align*}
\frac{q}{a-r}+\frac{q'}{r'-a}&=0,\\
\frac{q}{a+r}+\frac{q'}{r'+a}&=0.
\end{align*}
Solving these, we obtain
\begin{align*}
r'&=\frac{a^2}{r},\\
q'&=-\frac{a}{r}q.
\end{align*}
Hence the required force is
\[
F=\frac{1}{4\pi\varepsilon_0}\frac{qq'}{(r'-r)^2}
=-\frac{q^2}{4\pi\varepsilon_0}\frac{ar}{(a^2-r^2)^2}.
\]
When $r\to0$ this force is zero, and when $r\to a$ this force tends to infinity.

\begin{figbox}{3.12}
\begin{tikzpicture}[scale=0.95,>=latex]
  \coordinate (O) at (0,0);
  \coordinate (Q) at (1.4,0);
  \coordinate (Qp) at (3.46,0);
  % spherical shell
  \draw[thickline] (O) circle (2.2);
  % field lines of q (dashed) ending on the inner surface
  \foreach \a in {0,30,...,330}{%
    \draw[dashedline] (Q) -- (\a:2.2);}
  % line through O, A and B
  \draw[dashedline] (-2.5,0)--(3.95,0);
  \fill (2.2,0) circle (1.0pt) node[above right] {$A$};
  \fill (-2.2,0) circle (1.0pt) node[above left] {$B$};
  \node[below left] at (O) {$O$};
  \fill (Q) circle (1.2pt) node[above] {$q$};
  \fill (Qp) circle (1.2pt) node[above right] {$q'$};
  % dimensions
  \draw[dim] (0,1.15)--(1.4,1.15) node[midway,above] {$r$};
  \draw[dim] (0,-1.15)--(3.46,-1.15) node[midway,below] {$r'$};
\end{tikzpicture}
\end{figbox}

\begin{commentary}
This problem is another example of the method of images. It is not difficult to prove that the potential produced by $q$ and $q'$ over the whole spherical shell is zero.
\end{commentary}
\end{solution}

\begin{solution}{3.13}
(1) Let the capacitance between two adjacent conducting plates be $C$. Then when the positive terminal of the source is connected to the $N$-th plate, the source is, as it were, connected across the $(N-1)$ capacitors in series, so the $N$-th plate will carry the charge
\[
q_N=U\cdot\frac{C}{N-1}.
\]
At the same time, the left surface of the $(N-1)$-th plate carries the charge $-q_N$ and its right surface carries the charge $+q_N$; the charge distribution on each of the remaining conducting plates is similar to that on the $(N-1)$-th plate. When the positive terminal of the battery is next connected to the $(N-1)$-th plate, the charge on the left surface of this plate, being attracted by the positive charge on the $N$-th plate, no longer changes, while its right surface becomes the positive surface of $(N-2)$ capacitors in series and carries the charge
\[
q'_{N-1}=U\cdot\frac{C}{N-2}.
\]
Hence the total charge on the $(N-1)$-th plate is
\[
q_{N-1}=-U\frac{C}{N-1}+U\frac{C}{N-2}
=UC\left(\frac{1}{N-2}-\frac{1}{N-1}\right).
\]
By the same argument, the charge on the $k$-th plate is
\[
q_k=UC\left(\frac{1}{k-1}-\frac{1}{k}\right),
\]
and evidently
\[
q_1=-UC.
\]
Therefore
\[
|q_k|:|q_1|=\left(\frac{1}{k-1}-\frac{1}{k}\right):1.
\]

(2) From the above analysis it is known that the potential difference between the $N$-th plate and the $(N-1)$-th plate is $U/(N-1)$, that between the $(N-1)$-th plate and the $(N-2)$-th plate is $U/(N-2)$, and so on; hence the potential of the $k$-th plate is
\[
U_k=U\left(\frac{1}{k-1}+\frac{1}{k-2}+\cdots+\frac{1}{2}+1\right).
\]

\begin{commentary}
This problem can also be solved without the idea of capacitors in series, by regarding each charged conducting plate as an infinite charged plane and applying Gauss's law; but this is far more complicated than the capacitance method.
\end{commentary}
\end{solution}

\begin{solution}{3.14}
(1) When the radius is $r$, the resistance of the atmosphere is
\[
R=\int_r^\infty\rho\,\frac{\dd r}{4\pi r^2}=\frac{\rho}{4\pi r}.
\]
By Ohm's law,
\begin{equation}
-\frac{\rho}{4\pi r}\frac{\dd q}{\dd t}=V=\text{constant}, \tag{1}
\end{equation}
where $q$ is the charge carried by the balloon. The quantities $q$ and $r$ are connected by the following relation:
\begin{equation}
\frac{1}{4\pi\varepsilon_0}\frac{q}{r}=V. \tag{2}
\end{equation}
From this we obtain
\begin{equation}
q=4\pi\varepsilon_0V\cdot r. \tag{3}
\end{equation}
Differentiating,
\begin{equation}
\frac{\dd q}{\dd t}=4\pi\varepsilon_0V\frac{\dd r}{\dd t}. \tag{4}
\end{equation}
Substituting into equation (1) gives
\[
-\frac{\rho}{4\pi r}\cdot4\pi\varepsilon_0V\frac{\dd r}{\dd t}=V,
\]
that is,
\[
\frac{\dd r}{r}=-\frac{1}{\rho\varepsilon_0}\dd t.
\]
Integrating both sides, we obtain
\[
\ln r\Big|_a^r=-\frac{1}{\rho\varepsilon_0}t,
\]
hence
\[
\therefore\quad r=a\ee^{-\frac{1}{\rho\varepsilon_0}t}.
\]

(2) The Joule heat consumed in the whole leakage process is
\[
Q=\int_0^\infty\left(-\frac{\dd q}{\dd t}\right)^2R\,\dd t
=\int_0^\infty\left(-\frac{\dd q}{\dd t}\right)\cdot V\,\dd t=Vq_0,
\]
where $q_0$ is the charge on the balloon at $t=0$; from (3), $q_0=4\pi\varepsilon_0aV$. The original electrostatic energy of the balloon is
\[
W=\frac{1}{2}q_0V=\frac{1}{2}Q.
\]
The Joule heat consumed is greater than the original electrostatic energy because, during the gradual decrease of the radius of the balloon, external forces (the atmospheric pressure and the elastic force of the balloon) do work, so that the electrostatic energy is continually replenished.
\end{solution}

\begin{solution}{3.15}
Since the outer shell is grounded and the inner shell is connected by a wire to a conducting sphere of radius $c$ (conductor 3), the spherical capacitor formed by the inner and outer shells may be regarded as connected in parallel with the spherical conductor 3 (which acts as a capacitor). Let the charge on conductor 3 be $q$; then the charge on the inner shell is $Q_a=(Q-q)$, the two charges being proportional to their respective capacitances. The capacitance $C_{ab}$ of the spherical capacitor and the capacitance $C_c$ of conductor 3 are respectively
\[
C_{ab}=\frac{4\pi\varepsilon_0}{\dfrac{1}{a}-\dfrac{1}{b}},\qquad
C_c=4\pi\varepsilon_0c.
\]
Hence
\[
q=\frac{c}{c+\dfrac{1}{\dfrac{1}{a}-\dfrac{1}{b}}}\,Q
=\frac{c\left(\dfrac{1}{a}-\dfrac{1}{b}\right)}{c\left(\dfrac{1}{a}-\dfrac{1}{b}\right)+1}\,Q,
\]
\[
Q_a=Q-q=\frac{1}{c\left(\dfrac{1}{a}-\dfrac{1}{b}\right)+1}\,Q.
\]
The inner surface of the outer shell will carry the charge $-Q_a$. Owing to the induction of the charge on conductor 3, the outer surface of the outer shell will also carry a small amount of charge. Let this charge be $q'$; since the potential of the outer shell is zero,
\[
\frac{q'}{b}+\frac{q}{r}=0,
\]
and hence
\[
q'=-\frac{b}{r}q.
\]
Since the inner shell and the inner surface of the outer shell carry equal and opposite charges, they produce no electric field outside themselves; therefore the electrostatic force acting on conductor 3 is in fact the force between $q$ and $q'$:
\[
F=\frac{1}{4\pi\varepsilon_0}\frac{qq'}{r^2}
=\frac{1}{4\pi\varepsilon_0r^2}\left(-\frac{b}{r}\right)q^2
=-\frac{1}{4\pi\varepsilon_0r^3}
\frac{bc^2\left(\dfrac{1}{a}-\dfrac{1}{b}\right)^2}
{\left[c\left(\dfrac{1}{a}-\dfrac{1}{b}\right)+1\right]^2}Q^2.
\]
The negative sign indicates that the force is attractive.
\end{solution}

\begin{solution}{3.16}
To facilitate the solution, we first examine the behaviour of an electrically neutral conducting sphere in a uniform external field. The conducting sphere may be regarded as the superposition of two spheres carrying equal and opposite volume charges. Let the radius of the sphere be $R$ and the volume charge density be $\pm\rho$. When an external field $\vec{E}_0$ is present, the two charged spheres undergo a small displacement in the direction of the field, so as to set up inside the sphere a reverse field that cancels the external field, ultimately making the resultant field inside the sphere zero.

Let the displacement of the two spheres be $l$. Consider now the electric field at an arbitrary point $P$ inside the sphere (see the figure). Let the distance from $P$ to the centre $O_+$ of the positively charged sphere be $r_+$, and the distance from $P$ to the centre $O_-$ of the negatively charged sphere be $r_-$. By Gauss's law, the field produced by the positively charged sphere at $P$ is
\[
\vec{E}_+=\frac{1}{4\pi\varepsilon_0}\frac{\frac{4}{3}\pi r_+^3\rho}{r_+^3}\,\vec{r}_+
=\frac{\rho}{3\varepsilon_0}\vec{r}_+.
\]
Similarly, the field produced by the negatively charged sphere at $P$ is
\[
\vec{E}_-=-\frac{\rho}{3\varepsilon_0}\vec{r}_-.
\]
Here $\vec{r}_+$ and $\vec{r}_-$ are the vectors from $O_+$ and $O_-$ respectively to the point $P$; hence the resultant field is
\[
\vec{E}=\vec{E}_++\vec{E}_-
=\frac{\rho}{3\varepsilon_0}\left(\vec{r}_+-\vec{r}_-\right)
=-\frac{\rho}{3\varepsilon_0}\vec{l},
\]
where $\vec{l}$ is the vector from $O_-$ to $O_+$. It is seen that the field inside the sphere is uniform. It is only necessary to set $\vec{E}=-\vec{E}_0$ to obtain the actual value of the displacement:
\[
\frac{\rho}{3\varepsilon_0}l=E_0,
\]
and hence
\[
l=\frac{3\varepsilon_0E_0}{\rho}.
\]
Since $\rho$ is very large for a conductor, $l$ is a small quantity; therefore, as a result of the displacement of the two charged spheres, the greater part of the overlapping region inside the sphere remains electrically neutral, and a thin layer of charge appears only in the two thin non-overlapping regions at the sides, which may be regarded as a surface charge. The surface charge density $\sigma$ is a function of the angle $\theta$ (this is in fact the induced charge on the conductor):
\[
\sigma=\rho d=\rho l\cos\theta.
\]
Substituting the expression for $l$, we obtain
\begin{equation}
\sigma=\sigma_0\cos\theta=3\varepsilon_0E_0\cos\theta. \tag{1}
\end{equation}
When $\theta<90\degree$, $\sigma>0$; when $\theta>90\degree$, $\sigma<0$; when $\theta=90\degree$, $\sigma=0$.

\begin{figbox}{3.16}
\begin{tikzpicture}[scale=1.0,>=latex]
  \draw[thickline] (0,0) circle (1.6);
  % induced surface charges
  \foreach \a in {-55,-18,18,55}{\node at (\a:1.80) {$+$};}
  \foreach \a in {125,162,198,235}{\node at (\a:1.80) {$-$};}
  % centres of the two charged spheres
  \coordinate (Op) at (0.22,0);
  \coordinate (Om) at (-0.22,0);
  \fill (Op) circle (1.1pt) node[below right] {$O_+$};
  \fill (Om) circle (1.1pt) node[below left] {$O_-$};
  % point P and the vectors r_+, r_-
  \coordinate (P) at (0.58,0.85);
  \fill (P) circle (1.1pt) node[above right] {$P$};
  \draw[force] (Op)--(P) node[pos=0.55,right] {$\vec{r}_+$};
  \draw[force] (Om)--(P) node[pos=0.55,left] {$\vec{r}_-$};
  % external field
  \draw[force] (1.05,2.05)--++(1.05,0) node[right] {$\vec{E}_0$};
\end{tikzpicture}
\end{figbox}

Now we return to this problem. Imagine a conducting sphere of radius $R$ placed in a uniform electric field $\vec{E}_0$ directed vertically upwards; then, from the above discussion, induced charges will appear on the surface of the conductor, their surface density being given by expression (1), in which the angle $\theta$ should now be understood as the angle with the vertical direction. At this time the horizontal plane through the centre of the sphere must be an equipotential surface of the conducting sphere; if a conducting flat plate is placed in this plane, it does not affect the distribution of the electric field in the upper part of the conducting sphere. Therefore, in the situation of the present problem, the charge distribution on the hemisphere is still determined by (1), and induced charges will appear on the conducting plate; their effect above the conducting plate is just like that produced by the charges on the original lower hemisphere (the charges on the lower hemisphere may be regarded as the image charges of the charges on the upper hemisphere), while below the conducting plate the electric field is zero everywhere. Thus the electric force acting on the hemisphere is the force on the surface charge of the hemisphere. This resultant force must be directed upwards, and its magnitude $F$ can be found by integrating the vertical component of the normal force acting on each surface element of charge. We know that the force per unit area on the surface of a charged conductor is equal to $\dfrac{1}{2\varepsilon_0}\sigma^2$; hence
\begin{align*}
F&=\int_0^{\pi/2}\frac{1}{2\varepsilon_0}\sigma^2\cdot2\pi R\sin\theta\,R\,\dd\theta\cdot\cos\theta\\
&=\frac{\pi R^2\sigma_0^2}{\varepsilon_0}\int_0^{\pi/2}\sin\theta\cos^3\theta\,\dd\theta\\
&=-\frac{\pi R^2\sigma_0^2}{\varepsilon_0}\int_0^{\pi/2}\cos^3\theta\,\dd(\cos\theta)\\
&=\frac{\pi R^2\sigma_0^2}{4\varepsilon_0}.
\end{align*}
Substituting $\sigma_0=3\varepsilon_0E_0$, we obtain
\begin{equation}
F=\frac{9}{4}\pi R^2\varepsilon_0E_0^2. \tag{2}
\end{equation}
Setting $F=mg=\dfrac{2}{3}\pi R^3\rho g$, we obtain
\[
E_0=\frac{2}{3}\sqrt{\frac{2R\rho g}{3\varepsilon_0}}.
\]
This is the minimum value of the electric field that can lift the hemisphere off the flat plate.
\end{solution}

\begin{solution}{3.17}
(1) Imagine that the cylindrical wire carries positive charge, with a linear charge density $+\eta$. By the method of images, the electric field above the ground is excited jointly by this wire and by a cylindrical wire with linear charge density $-\eta$ located at the symmetric position below the ground.

When $r\ll d$, the charge on the wire may be regarded as uniformly distributed, and the electric field in space may be regarded as excited by two uniformly charged straight lines lying on the axes of the wires, with linear charge densities $+\eta$ and $-\eta$ respectively.

Take the ground as the $yz$-plane, and let the two charged straight lines run along the $z$ direction, located at $x=\pm a$ respectively, as shown in the figure (in this part $a=d$). Take the $y$-axis (actually the $yz$-plane) as the zero-potential surface. The electric field at an arbitrary point $P(x,y)$ is
\begin{align*}
\vec{E}_P&=\vec{E}_++\vec{E}_-\\
&=\frac{\eta}{2\pi\varepsilon_0}\left(\frac{1}{r_+}\hat{r}_+-\frac{1}{r_-}\hat{r}_-\right).
\end{align*}
The potential at $P$ is the integral of the electric field from $P$ to $O$ (the zero point of potential). We now want the potential of the thin wire; it is most convenient to integrate along the $x$-axis, taking the starting point of the integration at $A$, the intersection of the surface of the thin wire with the $x$-axis, and setting $a=d$. The potential of the wire is
\[
\varphi_A=\frac{\eta}{2\pi\varepsilon_0}\int_r^d\frac{\dd r}{r}
-\frac{\eta}{2\pi\varepsilon_0}\int_{2d-r}^d\frac{\dd r}{r}
\approx\frac{\eta}{2\pi\varepsilon_0}\ln\frac{2d}{r}.
\]
The required capacitance is
\[
C^*=\frac{\eta}{\varphi}=\frac{2\pi\varepsilon_0}{\ln\dfrac{2d}{r}}.
\]

\begin{figbox}{3.17}
\begin{tikzpicture}[scale=1.0,>=latex]
  \draw[axis] (-3.0,0)--(3.4,0) node[right] {$x$};
  \draw[axis] (0,-2.0)--(0,2.5) node[above] {$y$};
  \node[below left] at (0,0) {$O$};
  % the two wires (cross-sections)
  \draw[thickline] (-1.5,0) circle (0.12);
  \draw[thickline] (1.5,0) circle (0.12);
  \node[below left] at (-1.6,-0.16) {$-\eta$};
  \node[below right] at (1.6,-0.16) {$+\eta$};
  % dimensions a
  \draw[dim] (-1.5,-0.62)--(0,-0.62) node[midway,below] {$a$};
  \draw[dim] (0,-0.62)--(1.5,-0.62) node[midway,below] {$a$};
  % point P
  \coordinate (P) at (2.45,1.75);
  \fill (P) circle (1.1pt) node[right] {$P(x,y)$};
  \draw[thickline] (-1.5,0)--(P) node[pos=0.62,below right] {$r_-$};
  \draw[thickline] (1.5,0)--(P) node[pos=0.45,right] {$r_+$};
  % point A
  \fill (1.37,0) circle (1.0pt);
  \node[above left] at (1.36,0.10) {$A$};
\end{tikzpicture}
\end{figbox}

(2) When $r$ is of the same order of magnitude as $d$, the charge on the wire is no longer distributed uniformly over the cylindrical surface; being attracted by the charge of opposite sign on the image wire, it shifts towards the side of the image. The electric field in space can therefore no longer be regarded as the superposition of the fields of two lines of opposite sign placed along the axes. It can, however, still be regarded as the superposition of the fields of two lines of opposite sign parallel to the axis of the wire, the positions of both lines being shifted from the axes towards each other, while the linear charge densities remain unchanged. This is because the equipotential surfaces of two parallel straight lines carrying charges of opposite sign are a family of non-coaxial cylindrical surfaces.

For this purpose, consider the equipotential surfaces of the two charged straight lines in the figure. Let the distance of each of the two charged lines from the $y$-axis be $a$ ($a\ne d$). According to the discussion above, the potential at an arbitrary point $P$ is
\begin{align*}
\varphi_P&=\int_P^O\vec{E}\cdot\dd\vec{r}
=\int_P^O\vec{E}_+\cdot\dd\vec{r}+\int_P^O\vec{E}_-\cdot\dd\vec{r}\\
&=\frac{\eta}{2\pi\varepsilon_0}\int_{r_+}^a\frac{\dd r}{r}
-\frac{\eta}{2\pi\varepsilon_0}\int_{r_-}^a\frac{\dd r}{r}\\
&=\frac{\eta}{2\pi\varepsilon_0}\ln\frac{a}{r_+}
-\frac{\eta}{2\pi\varepsilon_0}\ln\frac{a}{r_-}
=\frac{\eta}{2\pi\varepsilon_0}\ln\frac{r_-}{r_+}.
\end{align*}
Thus the points whose potential is $\varphi$ satisfy
\[
\frac{r_-}{r_+}=\ee^{2\pi\varepsilon_0\varphi/\eta}.
\]
Substituting the coordinates of the point $P(x,y)$ into the above expression and squaring, we obtain
\[
\frac{(x+a)^2+y^2}{(x-a)^2+y^2}=\ee^{4\pi\varepsilon_0\varphi/\eta}.
\]
Setting $\ee^{4\pi\varepsilon_0\varphi/\eta}=k^2$, the equation of the equipotential surface of potential $\varphi$ is
\[
(x+a)^2+y^2=k^2(x-a)^2+k^2y^2.
\]
Simplified, this gives
\[
\left(x-\frac{k^2+1}{k^2-1}a\right)^2+y^2=\frac{4k^2}{(k^2-1)^2}a^2.
\]
The above is the equation of a circle, whose centre and radius both depend on $k^2$. Let this circle be the cross-section of the cylindrical wire in the problem; then
\[
\left(\frac{k^2+1}{k^2-1}\right)^2a^2=d^2,\qquad
\frac{4k^2}{(k^2-1)^2}a^2=r^2.
\]
From these we solve for $k^2$ (noting that $k>1$):
\[
k^2=\left[\frac{d}{r}+\sqrt{\left(\frac{d}{r}\right)^2-1}\right]^2.
\]
From the definition $k^2=\ee^{4\pi\varepsilon_0\varphi/\eta}$ given above we then obtain
\[
C^*=\frac{\eta}{\varphi}=\frac{4\pi\varepsilon_0}{\ln k^2}
=\frac{2\pi\varepsilon_0}{\ln k}
=\frac{2\pi\varepsilon_0}{\ln\left[\dfrac{d}{r}+\sqrt{\left(\dfrac{d}{r}\right)^2-1}\right]}.
\]
When $r\ll d$, the result agrees with that of part (1).
\end{solution}

\begin{solution}{3.18}
(1) The electric field inside a uniformly charged spherical surface of radius $a$ is zero. This can be explained by Gauss's law, or proved directly: for an arbitrary point $P$ inside the spherical surface, draw a pair of small opposite cones with vertex at $P$; they cut out two surface elements on the sphere (see the figure of Solution 3.01). From the geometrical relations and the inverse-square law of the electric force it follows that the fields produced at $P$ by the charges on these two surface elements cancel each other. The action of the other surface elements is similar. Hence the electric field at $P$ is zero, that is, the electric field at every point inside the spherical surface is zero. Imagine that the spherical surface is made of an insulating material, and now let it be charged and then be deformed into a spheroid (keeping its volume unchanged). Because the charges repel one another, the charge will redistribute so that the spheroidal surface becomes an equipotential surface, and the electric field inside the spheroid is zero. By the condition of electrostatic equilibrium, this is just the charge distribution on a spheroidal conductor. It is only necessary to find the potential produced by this charged spheroid at its centre, and the capacitance can then be found.

Let the surface charge density on the original spherical surface of radius $a$ be $\sigma_0$. When the charge on the ring on this spherical surface between the cone angles $\theta$ and $\theta+\dd\theta$ is stretched onto the spheroid, its amount remains unchanged, but its distance to the centre $O$ increases to $OA$, as shown in the figure. The potential it produces at $O$ is
\begin{align*}
\dd\varphi&=\frac{1}{4\pi\varepsilon_0}\frac{\sigma_0\,2\pi a\sin\theta\,a\,\dd\theta}{OA}\\
&=\frac{1}{4\pi\varepsilon_0}\frac{\sigma_0\cdot2\pi a\sin\theta\,a\,\dd\theta}
{\sqrt{a^2\sin^2\theta+\left(a\cos\theta\cdot\dfrac{b}{a}\right)^2}}\\
&=\frac{\sigma_0\,2\pi a^2}{4\pi\varepsilon_0}
\frac{\sin\theta\,\dd\theta}{\sqrt{a^2\sin^2\theta+b^2\cos^2\theta}}\\
&=-\frac{q}{8\pi\varepsilon_0a}
\frac{\dd(\cos\theta)}{\sqrt{1+\dfrac{b^2-a^2}{a^2}\cos^2\theta}},
\end{align*}
where $q=4\pi a^2\sigma_0$ is the total charge on the spherical surface.

\begin{figbox}{3.18}
\begin{tikzpicture}[scale=0.95,>=latex]
  % z-axis
  \draw[axis] (-2.5,0)--(3.0,0) node[right] {$z$};
  % original sphere and the spheroid
  \draw[dashedline] (0,0) circle (1.55);
  \draw[thickline] (0,0) ellipse (2.05 and 1.55);
  \node[below left] at (0,0) {$O$};
  % ring element and its image A
  \coordinate (Ae) at ({2.05*cos(65)},{1.55*sin(65)});
  \draw[dashedline] (0,0)--(Ae);
  \draw[dashedline] (0,0)--({2.05*cos(55)},{1.55*sin(55)});
  \fill (Ae) circle (1.1pt) node[above right] {$A$};
  \draw[thickline] (55:1.25) arc (55:65:1.25);
  \node at (60:0.98) {$\dd\theta$};
\end{tikzpicture}
\end{figbox}

(i) When $b>a$,
\[
\dd\varphi=-\frac{q}{8\pi\varepsilon_0a}\cdot\frac{1}{c}\cdot
\frac{\dd(c\cos\theta)}{\sqrt{1+c^2\cos^2\theta}},
\]
where $c^2=\dfrac{b^2-a^2}{a^2}$; hence the potential at $O$ is
\begin{align*}
\varphi&=-\frac{q}{8\pi\varepsilon_0ac}\int_0^\pi
\frac{\dd(c\cos\theta)}{\sqrt{1+c^2\cos^2\theta}}
=-\frac{q}{8\pi\varepsilon_0}\frac{1}{\sqrt{b^2-a^2}}
\int_c^{-c}\frac{\dd x}{\sqrt{1+x^2}}\\
&=\frac{q}{8\pi\varepsilon_0\sqrt{b^2-a^2}}
\ln\left(x+\sqrt{1+x^2}\right)\Big|_{-c}^{c}\\
&=\frac{q}{8\pi\varepsilon_0\sqrt{b^2-a^2}}
\ln\frac{b+\sqrt{b^2-a^2}}{b-\sqrt{b^2-a^2}}.
\end{align*}
Hence the capacitance is
\[
C=\frac{q}{\varphi}
=\frac{8\pi\varepsilon_0\sqrt{b^2-a^2}}
{\ln\dfrac{b+\sqrt{b^2-a^2}}{b-\sqrt{b^2-a^2}}}
\qquad(b>a).
\]

(ii) When $b<a$,
\[
\dd\varphi=-\frac{q}{8\pi\varepsilon_0a}\cdot\frac{1}{c}\cdot
\frac{\dd(c\cos\theta)}{\sqrt{1-c^2\cos^2\theta}},
\]
where $c^2=\dfrac{a^2-b^2}{a^2}$, and the potential at $O$ is
\[
\varphi=-\frac{q}{8\pi\varepsilon_0ac}\int_0^\pi
\frac{\dd(c\cos\theta)}{\sqrt{1-c^2\cos^2\theta}}
=-\frac{q}{8\pi\varepsilon_0}\frac{1}{\sqrt{a^2-b^2}}
\int_c^{-c}\frac{\dd x}{\sqrt{1-x^2}}
\]
\[
=-\frac{q}{8\pi\varepsilon_0\sqrt{a^2-b^2}}\arcsin x\Big|_c^{-c}
=\frac{q}{4\pi\varepsilon_0\sqrt{a^2-b^2}}\arcsin\frac{\sqrt{a^2-b^2}}{a}.
\]
Therefore
\[
\therefore\quad C=\frac{q}{\varphi}
=\frac{4\pi\varepsilon_0\sqrt{a^2-b^2}}{\arcsin\dfrac{\sqrt{a^2-b^2}}{a}}
\qquad(b<a).
\]

(2) When $b\gg a$ (needle-like),
\[
\ln\frac{b+\sqrt{b^2-a^2}}{b-\sqrt{b^2-a^2}}
=\ln\frac{1+\sqrt{1-\dfrac{a^2}{b^2}}}{1-\sqrt{1-\dfrac{a^2}{b^2}}}
\approx2\ln\frac{2b}{a},
\]
\[
C=\frac{8\pi\varepsilon_0\sqrt{b^2-a^2}}
{\ln\dfrac{b+\sqrt{b^2-a^2}}{b-\sqrt{b^2-a^2}}}
\approx\frac{4\pi\varepsilon_0b}{\ln\dfrac{2b}{a}}
\qquad(b\gg a).
\]
When $b\ll a$ (disk-like), $\arcsin\dfrac{\sqrt{a^2-b^2}}{a}\approx\dfrac{\pi}{2}$, and
\[
C=\frac{4\pi\varepsilon_0\sqrt{a^2-b^2}}{\arcsin\dfrac{\sqrt{a^2-b^2}}{a}}
\approx8\varepsilon_0a
\qquad(b\ll a).
\]

\begin{commentary}
This problem skilfully uses the method of deformation to determine the charge distribution on a spheroidal conductor under electrostatic equilibrium, so that its capacitance problem is readily solved.

For a disk-like conductor, because the charge on it is not uniformly distributed, the capacitance obtained differs from the result ($2\pi\varepsilon_0a$) obtained by assuming a uniform charge distribution. For a needle-like conductor, because the charge on it is essentially uniformly distributed, the two results are the same. The reader may compare and check them.
\end{commentary}
\end{solution}

\begin{solution}{3.19}
(1) First find the electric field inside the cylindrical surface. Let the linear charge density on the wire be $\lambda$; then
\[
E=\frac{1}{2\pi\varepsilon_0}\frac{\lambda}{r}.
\]
Also,
\[
V_0=\int_{r_0}^R E\,\dd r=\frac{\lambda}{2\pi\varepsilon_0}\int_{r_0}^R\frac{\dd r}{r}
=\frac{\lambda}{2\pi\varepsilon_0}\ln\frac{R}{r_0},
\]
\[
\therefore\quad E=\frac{V_0}{\ln\dfrac{R}{r_0}}\frac{1}{r}.
\]
Consider the motion of the ions. The velocity of an ion is
\[
\frac{\dd r}{\dd t}=wE=\frac{V_0w}{\ln\dfrac{R}{r_0}}\frac{1}{r}.
\]
Then
\[
r\frac{\dd r}{\dd t}=\frac{V_0w}{\ln\dfrac{R}{r_0}}.
\]
Integrating,
\[
\int_{r_0}^r r\,\dd r=\int_0^t\frac{V_0w}{\ln\dfrac{R}{r_0}}\,\dd t,
\]
that is,
\[
\frac{1}{2}\left(r^2-r_0^2\right)=\frac{V_0w}{\ln\dfrac{R}{r_0}}t,
\]
\[
\therefore\quad r^2=r_0^2+2\frac{V_0w}{\ln\dfrac{R}{r_0}}t,
\]
or
\[
r=\sqrt{r_0^2+2\frac{V_0w}{\ln(R/r_0)}t}.
\]

(2) Between $r_0$ and the radial position $r_i$ of the positive ions there is a reverse electric field. Let the counter have length $L$; then the linear charge density of the ions is $\lambda_i=\dfrac{Q_i}{L}$, and the potential difference of this reverse field is
\[
\Delta V=\frac{Q_i}{L\cdot2\pi\varepsilon_0}\ln\frac{r_i}{r_0}.
\]
This potential difference must be cancelled by the potential difference produced by the charge $Q$ flowing into the wire (at the same time a charge $-Q$ flows into the outer wall of the cylinder); this cancelling potential difference is
\[
\Delta V'=\frac{Q}{L\cdot2\pi\varepsilon_0}\ln\frac{R}{r_0}.
\]
Setting $\Delta V=\Delta V'$, we obtain
\[
Q=Q_i\frac{\ln\dfrac{r_i}{r_0}}{\ln(R/r_0)}.
\]
Substituting $r_i=\sqrt{r_0^2+2\dfrac{V_0w}{\ln(R/r_0)}t}$ into the above expression, we obtain
\[
Q=\frac{Q_i}{2\ln(R/r_0)}\ln\left(1+\frac{2V_0w}{r_0^2\ln(R/r_0)}t\right).
\]
This action lasts only until $r_i=R$, that is, it applies only for
\[
t\le\frac{R^2-r_0^2}{2\dfrac{V_0w}{\ln(R/r_0)}}=t_1.
\]
After $t\ge t_1$, $Q=Q_i$.
\end{solution}

\begin{solution}{3.20}
Imagine the point charge to be distributed on a small conducting sphere, and further imagine thin conducting spherical shells of radii $a$ and $b$ to exist at the inner and outer surfaces of the dielectric respectively. Since the conducting shells are very thin, the charges of opposite sign on them lie very close together, and their presence may be considered not to affect the original electric field. In this way the whole system may be regarded as three capacitors connected in series: the capacitance $C_1$ between the innermost small conducting sphere and the conducting shell at the inner surface of the dielectric, the capacitance $\varepsilon_rC_2$ of the spherical dielectric of inner and outer radii $a$ and $b$, and the capacitance $C_3$ between the conducting shell of radius $b$ and infinity. When the point charge is moved to infinity, at the place where the point charge was there still exists the capacitance between the small conducting sphere of the point charge and infinity; this capacitance can still be regarded as the series combination of the three vacuum capacitances $C_1$, $C_2$ and $C_3$. In this way the electrostatic energies of the two states differ only in the stored energy of the vacuum capacitance $C_2$ and of the dielectric capacitance $\varepsilon_rC_2$, and the work $A$ done by the external force is just equal to this difference of electrostatic energies:
\[
A=\frac{q^2}{2}\left(\frac{1}{C_2}-\frac{1}{\varepsilon_rC_2}\right)
=\frac{q^2}{2C_2}\left(1-\frac{1}{\varepsilon_r}\right).
\]
Also,
\[
C_2=\frac{4\pi\varepsilon_0}{\dfrac{1}{a}-\dfrac{1}{b}}.
\]
Substituting into the above expression, we obtain
\[
A=\frac{q^2}{8\pi\varepsilon_0}\left(\frac{1}{a}-\frac{1}{b}\right)
\left(1-\frac{1}{\varepsilon_r}\right).
\]

\begin{commentary}
Finding the work through the change of energy is a method commonly used in electrostatic problems. But attention must be paid to whether there are other processes doing work in the course of the process (heat-transfer problems are seldom involved in electrostatics), especially work done by a source (in this problem there is no work done by any source).

This problem can also be solved by directly calculating, by integration, the difference between the energies of the electrostatic field in the two cases.
\end{commentary}
\end{solution}

\begin{solution}{3.21}
Suppose that, when the dielectric particles are absent, a uniform electric field $E_0$ already exists between the plates. When dielectric particles are present, since the particles are sparse, the external field polarizing a given dielectric particle may be taken to be $E_0$. Under the action of $E_0$ the spherical dielectric particle is uniformly polarized; let its polarization be $P$. At the same time equivalent polarization charges are formed on the surface of the dielectric sphere, with a surface density $\sigma'=P\cos\theta$ (see the figure). The field produced by these polarization charges is also uniform, opposite in direction to $P$, and of magnitude $P/(3\varepsilon_0)$ (see Solution 3.16); and the dielectric is polarized precisely under the action of the resultant of $E_0$ and the field of these polarization charges. Hence
\[
P=(\varepsilon_r-1)\varepsilon_0\left[E_0-\frac{P}{3\varepsilon_0}\right].
\]
That is,
\[
\left(1+\frac{\varepsilon_r-1}{3}\right)P=(\varepsilon_r-1)\varepsilon_0E_0.
\]
From this we obtain
\[
P=\frac{3\varepsilon_0(\varepsilon_r-1)}{\varepsilon_r+2}E_0.
\]
The electric moment of each particle is
\[
p=\frac{4}{3}\pi a^3P
=\frac{4\pi\varepsilon_0a^3(\varepsilon_r-1)}{\varepsilon_r+2}E_0.
\]
The presence of the polarized particles is equivalent to a uniform polarization of the dielectric between the plates, with an equivalent polarization
\[
P'=np=4\pi\varepsilon_0a^3n\frac{\varepsilon_r-1}{\varepsilon_r+2}E_0.
\]
Since the particles are sparse, $P'$ must be very small; the reverse field produced by the polarization charges corresponding to $P'$ may be neglected, and the field polarizing this equivalent dielectric may be taken to be the external field $E_0$. Thus
\[
P'=(\varepsilon_r'-1)\varepsilon_0E_0.
\]
Comparing this with the above expression, we obtain the equivalent relative permittivity
\[
\varepsilon_r'=1+\frac{\varepsilon_r-1}{\varepsilon_r+2}\cdot4\pi a^3n.
\]

\begin{figbox}{3.21}
\begin{tikzpicture}[scale=1.0,>=latex]
  % uniform external field lines
  \foreach \y in {-0.9,-0.45,0.45,0.9}{%
    \draw[force] (-2.6,\y)--(-1.55,\y);
    \draw[force] (-1.15,\y)--(0.85,\y);
    \draw[force] (1.55,\y)--(2.6,\y);}
  \draw[force] (-2.6,0)--(-1.55,0);
  \draw[force] (1.55,0)--(2.6,0);
  % dielectric sphere
  \draw[thickline] (0,0) circle (1.55);
  % polarization charges on the surface
  \foreach \a in {-50,-25,0,25,50}{\node at (\a:1.75) {$+$};}
  \foreach \a in {130,155,180,205,230}{\node at (\a:1.75) {$-$};}
  % polarization vector
  \draw[force] (0,0)--(0.95,0) node[below right] {$\vec{P}$};
  % surface charge density label
  \draw[thickline] (45:1.55)--++(45:0.35);
  \node[right] at (1.20,1.28) {$\sigma'=P\cos\theta$};
  % external field label
  \node[above] at (-2.0,1.05) {$\vec{E}_0$};
\end{tikzpicture}
\end{figbox}
\end{solution}

\begin{solution}{3.22}
First, around this point charge there will necessarily gather polarization charge $q''$ of sign opposite to $q$, so that $q+q''=\dfrac{q}{\varepsilon_r}$; hence it suffices to examine the behaviour of the point charge $q/\varepsilon_r$.

By symmetry, it suffices to consider the surface polarization charge at a distance $r$ from the foot $O$ of the perpendicular from the point charge to the surface; let its density be $\sigma'(r)$. Above and below the surface the electric field is the superposition of the field $\vec{E}_1$ produced by the point charge $q/\varepsilon_r$ and the field $\vec{E}'$ produced by $\sigma'$. At the surface, in the normal direction, it follows from Gauss's law that the normal component of the field $\vec{E}'$ produced by $\sigma'$ is $\sigma'/(2\varepsilon_0)$, directed away from the surface upwards and downwards; hence
\[
E_{n,\text{up}}=\frac{1}{4\pi\varepsilon_0}\frac{q/\varepsilon_r\,h}{(h^2+r^2)^{3/2}}+\frac{\sigma'}{2\varepsilon_0},
\]
\[
E_{n,\text{down}}=\frac{1}{4\pi\varepsilon_0}\frac{q/\varepsilon_r\,h}{(h^2+r^2)^{3/2}}-\frac{\sigma'}{2\varepsilon_0}.
\]
From the continuity of the normal component of the electric displacement, $D_n$,
\[
\varepsilon_0E_{n,\text{up}}=\varepsilon_r\varepsilon_0E_{n,\text{down}}.
\]
That is,
\[
\frac{1}{4\pi\varepsilon_r}\frac{qh}{(h^2+r^2)^{3/2}}+\frac{\sigma'}{2}
=\frac{1}{4\pi}\frac{qh}{(h^2+r^2)^{3/2}}-\frac{\varepsilon_r\sigma'}{2}.
\]
Solving this, we obtain
\[
\sigma'=\frac{\varepsilon_r-1}{(\varepsilon_r+1)\varepsilon_r}\cdot
\frac{qh}{2\pi(h^2+r^2)^{3/2}}.
\]
The total polarization charge is
\[
q'=\int_0^\infty\sigma'\cdot2\pi r\,\dd r
=\frac{\varepsilon_r-1}{\varepsilon_r(\varepsilon_r+1)}
\int_0^\infty\frac{qhr\,\dd r}{(r^2+h^2)^{3/2}}
=\frac{\varepsilon_r-1}{\varepsilon_r(\varepsilon_r+1)}q.
\]
From the expression for $\sigma'(r)$ it is seen that the field produced by the polarization charges below the surface is the same as if there were a point charge $q'$ at a height $h$ above the surface; hence the electrostatic force acting on $q$ is the repulsion between $q$ and $q'$:
\[
F=\frac{1}{4\pi\varepsilon_0}\frac{q^2(\varepsilon_r-1)}{\varepsilon_r(\varepsilon_r+1)\cdot4h^2}.
\]

\begin{figbox}{3.22}
\begin{tikzpicture}[scale=0.95,>=latex]
  % interface
  \draw[thickline] (-3.2,0)--(3.2,0);
  \node[below left] at (-2.6,0) {$\varepsilon_r$};
  % vertical line through q', O and q
  \draw[dashedline] (0,1.55)--(0,-1.55);
  \fill (0,1.55) circle (1.2pt) node[above] {$q'$};
  \fill (0,-1.55) circle (1.2pt) node[below] {$q$};
  \node[above left] at (0,0) {$O$};
  % dimensions h
  \draw[dim] (0,0)--(0,1.55) node[midway,left] {$h$};
  \draw[dim] (0,-1.55)--(0,0) node[midway,left] {$h$};
  % observation point on the surface
  \coordinate (S) at (2.0,0);
  \fill (S) circle (1.0pt);
  \draw[dim] (0,-0.45)--(2.0,-0.45) node[midway,below] {$r$};
  \draw[dashedline] (0,-1.55)--(S);
  % fields at the surface point
  \draw[force] (S)--++(34:1.05) node[above right] {$\vec{E}_1$};
  \draw[force] (S)--++(-34:1.05) node[below right] {$\vec{E}'$};
\end{tikzpicture}
\end{figbox}

\begin{commentary}
The field above the surface is just the field of a point charge equal to $(q/\varepsilon_r+q')$. The charge $q'$ is in essence the ``image charge'' in the case of a dielectric.
\end{commentary}
\end{solution}

\begin{solution}{3.23}
(1) Let the charge per unit length on the metal rod be $\lambda$. Then from Gauss's law,
\[
E=\frac{\lambda}{2\pi\varepsilon_0 r},
\]
so
\[
U=\int_a^b E\,\dd r=\frac{\lambda}{2\pi\varepsilon_0}\int_a^b\frac{\dd r}{r}
=\frac{\lambda}{2\pi\varepsilon_0}\ln\frac{b}{a},
\]
and therefore
\[
E=\frac{U}{r\ln\dfrac{b}{a}}.
\]

(2) At the place where the dust lies the electric field may be regarded as uniform, and the spherical dielectric (the dust) will be uniformly polarised, the direction of the polarisation being the same as that of the field; denote the polarisation by $P$. Then
\[
P=(\varepsilon_r-1)\varepsilon_0E',
\]
where
\[
E'=E-\frac{P}{3\varepsilon_0}.
\]
(see Solution 3.21) Substituting into the above expression gives
\[
P=(\varepsilon_r-1)\varepsilon_0\left(E-\frac{P}{3\varepsilon_0}\right)
=(\varepsilon_r-1)\varepsilon_0E-\frac{\varepsilon_r-1}{3}P,
\]
from which we obtain
\[
P=\frac{3(\varepsilon_r-1)}{\varepsilon_r+2}\varepsilon_0E
=\frac{3(\varepsilon_r-1)}{\varepsilon_r+2}\varepsilon_0\frac{U}{r\ln\dfrac{b}{a}}.
\]
Here the expression for $E$ obtained in part (1) has already been substituted. The dust is equivalent to a dipole, and its dipole moment is
\[
p=\frac{4}{3}\pi c^3P
=\frac{4}{3}\pi c^3\frac{3(\varepsilon_r-1)}{\varepsilon_r+2}\varepsilon_0
\frac{U}{\ln\dfrac{b}{a}}\cdot\frac{1}{r}.
\]
The force acting on the dust is
\[
F=p\frac{\partial E}{\partial r}
=-\frac{4\pi\varepsilon_0c^3(\varepsilon_r-1)U^2}
{(\varepsilon_r+2)\left(\ln\dfrac{b}{a}\right)^2}\frac{1}{r^3}.
\]
The sign ``$-$'' shows that the direction of the force points towards the axis.

(3) The dust must move from the wall of the cylinder towards the metal rod, and its acceleration is
\[
a=\frac{\dd v}{\dd t}=\frac{F}{m}=-\frac{A}{r^3}
\qquad\left(\text{where }A=\frac{4\pi\varepsilon_0c^3(\varepsilon_r-1)U^2}
{m(\varepsilon_r+2)\left(\ln\dfrac{b}{a}\right)^2}\right).
\]
Multiplying both sides of the above expression by $\dd r$ and integrating, and noting that $\dfrac{\dd r}{\dd t}=v$, we have
\[
\int_0^v v\,\dd v=\int_b^r-\frac{A}{r^3}\,\dd r,
\]
that is,
\[
\frac{1}{2}v^2=\frac{A}{2}\left.\frac{1}{r^2}\right|_b^r
=\frac{A}{2}\left(\frac{1}{r^2}-\frac{1}{b^2}\right),
\]
and hence
\[
v=-\sqrt{A\left(\frac{1}{r^2}-\frac{1}{b^2}\right)}
=-\frac{\sqrt{A(b^2-r^2)}}{br}.
\]
Then
\[
\dd t=\frac{\dd r}{v}=-\frac{br\,\dd r}{\sqrt{A(b^2-r^2)}}.
\]
Integrating, we obtain
\begin{align*}
t=\int\dd t&=-\frac{b}{\sqrt{A}}\int_b^a\frac{r\,\dd r}{\sqrt{b^2-r^2}}
=\frac{b}{2\sqrt{A}}\int\frac{\dd(b^2-r^2)}{\sqrt{b^2-r^2}}\\
&=\frac{b}{2\sqrt{A}}\left.2\sqrt{b^2-r^2}\right|_b^a
=\frac{b}{\sqrt{A}}\sqrt{b^2-a^2}.
\end{align*}
\end{solution}

\begin{solution}{3.24}
(1) Polarisation charges will be formed on the wall of the cavity, and the superposition of the field produced by these polarisation charges inside the sphere with the original field $E_0$ is just the field in the cavity; outside the sphere the result of superposing the field produced by the polarisation charges with the original field will make the polarisation of the medium no longer uniform. Inspired by the result that in a uniform field the field inside a dielectric sphere is still uniform, the field inside the cavity should also be uniform. Hence we may suppose the surface density of the polarisation charge to be
\[
\sigma_p=-\sigma_0\cos\theta.
\]
It is only necessary to find $\sigma_0$, and then the field in the cavity can be found, because the polarisation charge with this distribution produces a uniform field inside the cavity:
\[
E'=\frac{\sigma_0}{3\varepsilon_0},
\]
whose direction is the same as that of $E_0$ (see Fig.~1). Hence the field in the cavity is
\[
\vec{E}=\vec{E}_0+\frac{\sigma_0}{3\varepsilon_0}\hat{\imath}.
\]

\begin{figbox}{3.24(a)}
\begin{tikzpicture}[scale=1.0,>=latex]
  % cavity sphere
  \draw[thickline] (0,0) circle (1.45);
  % external field arrows
  \draw[force] (-2.75,1.95)--(-1.55,1.95) node[right] {$\vec{E}_0$};
  \draw[force] (2.95,1.95)--(4.15,1.95) node[right] {$\vec{E}_0$};
  % polarisation charges
  \foreach \a in {130,150,170,190,210,230}{\node at (\a:1.22) {$+$};}
  \foreach \a in {-50,-30,-10,10,30,50}{\node at (\a:1.22) {$-$};}
  % axes
  \draw[axis] (0,0)--(2.15,0) node[right] {$x$};
  \node[below left] at (0,0) {$O$};
  % radius at angle theta
  \draw[thickline] (0,0)--(35:1.45);
  \draw[thickline] (0.45,0) arc (0:35:0.45);
  \node at (0.62,0.22) {$\theta$};
\end{tikzpicture}
\end{figbox}

In order to find $\sigma_0$ it must be related to the polarisation, and therefore the field produced by this polarisation charge outside the cavity should first be found. This can be done by using the boundary conditions. Denote the region inside the cavity by 1 and that outside by 2; then at $\theta$ we must have
\[
E_{2n}-E_{1n}=\frac{\sigma_P}{\varepsilon_0}=-\frac{\sigma_0}{\varepsilon_0}\cos\theta.
\]
But
\[
E_{1n}=\left(E_0+\frac{\sigma_0}{3\varepsilon_0}\right)\cos\theta,
\]
\[
\therefore\ E_{2n}=E_{1n}-\frac{\sigma_0}{\varepsilon_0}\cos\theta
=\left(E_0-\frac{2}{3}\frac{\sigma_0}{\varepsilon_0}\right)\cos\theta,
\]
\[
\therefore\ P_{2n}=(\varepsilon_r-1)\varepsilon_0E_{2n}
=(\varepsilon_r-1)\varepsilon_0E_0\cos\theta-\frac{2}{3}(\varepsilon_r-1)\sigma_0\cos\theta.
\]
However, from the relation between the polarisation charge and the polarisation, $P_{2n}=-\sigma'=\sigma_0\cos\theta$; substituting this into the above expression gives
\[
\sigma_0=\frac{3(\varepsilon_r-1)}{2\varepsilon_r+1}\varepsilon_0E_0.
\]
Substituting into the expression for $E$, we obtain
\[
E=E_0+\frac{\sigma_0}{3\varepsilon_0}
=E_0+\frac{\varepsilon_r-1}{2\varepsilon_r+1}E_0
=\frac{3\varepsilon_r}{2\varepsilon_r+1}E_0.
\]
Written in vector form,
\begin{equation}
\vec{E}=\frac{3\varepsilon_r}{2\varepsilon_r+1}\vec{E}_0. \tag{1}
\end{equation}

(2) Putting a charged conducting sphere into the cavity may be imagined as first putting an uncharged conducting sphere into the cavity and then placing charge on the conducting sphere. When the uncharged conducting sphere is put into the cavity, induced charges will appear on the conducting sphere, and the field of these induced charges will in turn change the polarisation of the medium, so that finally the field produced inside the sphere by the total charge (free charge plus polarisation charge) formed on the spherical surface should exactly cancel the original external field, making the field inside the sphere zero. Hence we may suppose the surface density of the total charge to be
\[
\sigma_t=\sigma_{t0}\cos\theta
\]

\begin{figbox}{3.24(b)}
\begin{tikzpicture}[scale=1.0,>=latex]
  \draw[thickline] (0,0) circle (1.45);
  \draw[axis] (0,0)--(2.35,0) node[right] {$x$};
  \node[below left] at (0,0) {$O$};
  \draw[thickline] (0,0)--(35:1.45);
  \draw[thickline] (0.45,0) arc (0:35:0.45);
  \node at (0.62,0.22) {$\theta$};
  \node at (35:1.05) {$\sigma_t$};
  \draw[force] (35:1.45)--(35:2.15) node[right] {$\vec{E}_1$};
  \node[right] at (1.05,-0.12) {$\sigma_{t0}$};
\end{tikzpicture}
\end{figbox}

(Fig.~2) where $\sigma_{t0}$ must satisfy
\[
\frac{\sigma_{t0}}{3\varepsilon_0}=E_0
\qquad(\text{that is, }\sigma_{t0}=3\varepsilon_0E_0),
\]
\[
\therefore\ \sigma_t=3\varepsilon_0E_0\cos\theta.
\]
By the boundary conditions, the field at $\theta$ on the outer surface of the conducting sphere (inside the medium) is perpendicular to the surface, and its value is
\[
E_1=\frac{\sigma_t}{\varepsilon_0}=3E_0\cos\theta.
\]
When a charge $Q$ is then placed on the conducting sphere, $Q$ must be distributed uniformly over the surface of the sphere (the field inside the sphere still being zero), and the polarisation charges caused by $Q$ will also be distributed uniformly over the surface of the sphere; at this time the total field at $\theta$ is
\begin{equation}
E=E_1+E_2=3E_0\cos\theta+\frac{Q}{4\pi\varepsilon_0\varepsilon_rR^2} \tag{2}
\end{equation}
and its direction is perpendicular to the spherical surface, where $E_2$ is the field excited after $Q$ is added. The force density on the conductor (the force per unit area) is equal to the energy density of the field, and by symmetry the resultant force must be along the direction of $E_0$ (that is, the $x$ direction); the component along the $x$ direction of the force on the belt from $\theta$ to $\theta+\dd\theta$ is
\[
\dd F_x=\frac{1}{2}\varepsilon_r\varepsilon_0E^2\cdot2\pi R\sin\theta\,R\,\dd\theta\cdot\cos\theta
=\varepsilon_r\varepsilon_0\pi R^2E^2\sin\theta\cos\theta\,\dd\theta.
\]
The resultant force is $F=\displaystyle\int\dd F_x=\int_0^\pi\varepsilon_r\varepsilon_0\pi R^2E^2\sin\theta\cos\theta\,\dd\theta$, and
\begin{align*}
E^2&=\left(3E_0\cos\theta+\frac{Q}{4\pi\varepsilon_0\varepsilon_rR^2}\right)^2\\
&=9E_0^2\cos^2\theta+\left(\frac{Q}{4\pi\varepsilon_0\varepsilon_rR^2}\right)^2
+6E_0\cos\theta\frac{Q}{4\pi\varepsilon_0\varepsilon_rR^2}.
\end{align*}
Only the last term contributes to the integral, so
\begin{align*}
\therefore\ F&=\int_0^\pi\varepsilon_r\varepsilon_0\pi R^2\cdot
\frac{6E_0Q}{4\pi\varepsilon_0\varepsilon_rR^2}\cos^2\theta\sin\theta\,\dd\theta\\
&=\frac{3}{2}E_0Q\int_0^\pi(-\cos^2\theta)\,\dd(\cos\theta)\\
&=\frac{3}{2}E_0Q\left.\frac{\cos^3\theta}{3}\right|_\pi^0\\
&=QE_0.
\end{align*}
\end{solution}

\begin{solution}{3.25}
The spherical cavity in which $q'$ lies and the medium near it are in the electric field excited by the charge $q$ and the surrounding medium. Since $q$ and $q'$ are very far apart, this field may be regarded as a uniform field, denoted by $E_q$. If this field did not exist, $q'$ would experience no force in the cavity, because the polarisation charges produced by the polarisation of the medium caused by the field excited by $q'$ are distributed uniformly on the surface of the spherical cavity and exert no force on $q'$. But $q'$ is also not simply acted on by $E_q$, because the existence of $E_q$ makes the medium around $q'$ produce an additional polarisation, so that the field in the cavity is no longer equal to $E_q$. The problem comes down to finding the field inside a spherical cavity in a uniform electric field $E_q$. According to Solution 3.24, the magnitude of this field is
\[
E=\frac{3\varepsilon_r}{2\varepsilon_r+1}E_q.
\]
(see equation (1) of Solution 3.24) Neglecting the influence of $q'$ and the polarisation charges around it on the polarisation of the medium near $q$ (this is a second-order effect), $E_q$ is evidently just the field of the point charge $q$ partially shielded by the medium, that is,
\[
E_q=\frac{1}{4\pi\varepsilon_0}\frac{q}{\varepsilon_rr^2}.
\]
The required force is
\[
F'=q'E=q'\frac{3\varepsilon_r}{2\varepsilon_r+1}\cdot\frac{1}{4\pi\varepsilon_0}\frac{q}{\varepsilon_rr^2}.
\]
\[
=\frac{1}{4\pi\varepsilon_0}\frac{3qq'}{(2\varepsilon_r+1)r^2}.
\]
\end{solution}

\begin{solution}{3.26}
From the hint it is known that the distribution of the polarisation charge on the surface of the dielectric sphere is the same as in the case of uniform polarisation. Let the polarisation charge be $\sigma=\sigma_0\cos\theta$, the axis of the polar coordinates being the same as the direction of the dipole $p$ at the centre of the sphere. Then the field outside the sphere is produced jointly by the dipole moment formed by the surface polarisation charge of the sphere and the equivalent dipole moment $p/\varepsilon_r$, which is the dipole at the centre of the sphere shielded by the medium; the total dipole moment is
\[
p_t=\frac{p}{\varepsilon_r}+\frac{4\pi}{3}a^3\sigma_0.
\]
(Note that the surface density of the polarisation charge on the surface of a uniformly polarised sphere with polarisation $P$ is $\sigma=P\cos\theta$.) The field inside the sphere is the superposition of the field produced by the equivalent dipole $p/\varepsilon_r$ at the centre of the sphere and the uniform field produced by the surface polarisation charge. The problem comes down to finding $\sigma_0$, which can be found from the boundary conditions.

Let the region inside the sphere be region 2 and the region outside the sphere be region 1; then at the surface $(a,\theta)$,
\begin{align*}
E_{1r}&=\frac{1}{4\pi\varepsilon_0}\frac{(2p/\varepsilon_r)\cos\theta}{a^3}
+\frac{1}{4\pi\varepsilon_0}\frac{2\cdot\dfrac{4}{3}\pi a^3\sigma_0\cos\theta}{a^3}\\
&=\frac{1}{4\pi\varepsilon_0}\frac{2p\cos\theta}{\varepsilon_ra^3}
+\frac{2\sigma_0}{3\varepsilon_0}\cos\theta,\\[2pt]
E_{2r}&=\frac{1}{4\pi\varepsilon_0}\frac{(2p/\varepsilon_r)\cos\theta}{a^3}
-\frac{\sigma_0}{3\varepsilon_0}\cos\theta\\
&=\frac{1}{4\pi\varepsilon_0}\frac{2p\cos\theta}{\varepsilon_ra^3}
-\frac{\sigma_0}{3\varepsilon_0}\cos\theta.
\end{align*}
From the boundary condition $\varepsilon_rE_{2r}=E_{1r}$ we obtain
\[
\frac{2p\cos\theta}{4\pi\varepsilon_0a^3}-\frac{\varepsilon_r\sigma_0}{3\varepsilon_0}\cos\theta
=\frac{2p\cos\theta}{4\pi\varepsilon_0\varepsilon_ra^3}+\frac{2\sigma_0}{3\varepsilon_0}\cos\theta.
\]
From this we obtain
\[
\sigma_0=\frac{3(\varepsilon_r-1)}{\varepsilon_r(\varepsilon_r+2)}\frac{p}{2\pi a^3}.
\]
Substituting into the expression for $p_t$ gives
\begin{align*}
p_t&=\frac{p}{\varepsilon_r}+\frac{4\pi}{3}a^3\cdot
\frac{3(\varepsilon_r-1)}{\varepsilon_r(\varepsilon_r+2)}\frac{p}{2\pi a^3}\\
&=\left[\frac{1}{\varepsilon_r}+\frac{2(\varepsilon_r-1)}{\varepsilon_r(\varepsilon_r+2)}\right]p\\
&=\frac{3}{\varepsilon_r+2}p.
\end{align*}
The potential outside the sphere is just the potential of $p_t$. Taking the potential at infinity as zero, we have
\[
\varphi_1=\frac{1}{4\pi\varepsilon_0}\frac{p_t\cos\theta}{r^2}
=\frac{1}{4\pi\varepsilon_0}\frac{3}{\varepsilon_r+2}\frac{p\cos\theta}{r^2}
\qquad(r\ge a).
\]
The potential inside the sphere is the superposition of the potential of the uniform field $\dfrac{\sigma_0}{3\varepsilon_0}$ and the potential of the dipole $p/\varepsilon_r$:
\begin{align*}
\varphi_2&=\frac{(p/\varepsilon_r)\cos\theta}{4\pi\varepsilon_0r^2}
+\frac{\sigma_0}{3\varepsilon_0}r\cos\theta\\
&=\frac{p\cos\theta}{4\pi\varepsilon_0\varepsilon_r}
\left[\frac{1}{r^2}+\frac{2(\varepsilon_r-1)}{(\varepsilon_r+2)a^3}r\right]
\qquad(r\le a).
\end{align*}
\end{solution}

\begin{solution}{3.27}
The current in the resistive material is the superposition of the inflowing current and the outflowing current. When only the inflowing current exists, the current density in the medium is
\[
j_1=\frac{I}{2\pi r^2}.
\]
From the differential form of Ohm's law, the field is
\[
E_1=\rho j_1=\frac{\rho I}{2\pi r^2},
\]
where $r$ is the distance from the point under consideration to the point at which the current enters. Compared with the field of a point charge, $E=\dfrac{1}{4\pi\varepsilon_0}\dfrac{q}{r^2}$, $E_1$ is equivalent to the field of a point charge with $q=2\varepsilon_0\rho I$; hence the potential difference between the positions of the two ends of the voltmeter is
\[
U_1=\frac{q}{4\pi\varepsilon_0}\left(\frac{1}{d}-\frac{1}{2d}\right)
=\frac{q}{4\pi\varepsilon_0}\cdot\frac{1}{2d}.
\]
Substituting $q=2\varepsilon_0\rho I$, we obtain
\[
U_1=\frac{\rho I}{2\pi}\cdot\frac{1}{2d}.
\]
Similarly, when only the outflowing current exists the situation is analogous, and moreover on the line joining the points at which the current enters and leaves (that is, the straight line on which the four probe tips lie) the directions of the current are the same, and the corresponding potential differences are also the same:
\[
U_2=\frac{\rho I}{2\pi}\cdot\frac{1}{2d}.
\]
When the inflowing and the outflowing currents exist simultaneously, the reading of the voltmeter is the sum of $U_1$ and $U_2$:
\[
U=U_1+U_2=\frac{\rho I}{2\pi d}.
\]
From this we obtain
\[
\rho=\frac{2\pi dU}{I}.
\]
\end{solution}

\begin{solution}{3.28}
Let the required resistance be $xR$. For convenience of calculation in what follows, we may take $R$ as the unit of resistance; then the required resistance is set equal to $x$. Since resistance is proportional to length, the resistance between two opposite corners of the small network left after the two outermost squares are removed should be $x/2$, so the original network may be simplified to Fig.~(1). (The values of the resistances in the lower half of the figure are the same as those in the upper half, and are not marked.)

\begin{figbox}{3.28}
\begin{tikzpicture}[scale=0.86,>=latex]
% ---------- panel (1) ----------
\begin{scope}[xshift=0cm]
  \coordinate (A) at (0,2);   \coordinate (B) at (2.8,2);
  \coordinate (T) at (1.4,3.4); \coordinate (D) at (1.4,0.6);
  \coordinate (P) at (0.7,2); \coordinate (Q) at (2.1,2);
  \coordinate (U) at (0.7,2.7); \coordinate (V) at (2.1,2.7);
  \coordinate (W) at (0.7,1.3); \coordinate (X) at (2.1,1.3);
  \draw[thickline] (A)--(T)--(B)--(D)--cycle;
  \draw[thickline] (A)--(P); \draw[thickline] (Q)--(B);
  \draw[thickline] (T)--(U); \draw[thickline] (T)--(V);
  \draw[thickline] (U)--(V); \draw[thickline] (U)--(P); \draw[thickline] (V)--(Q);
  \draw[thickline] (P)--(V);
  \draw[thickline] (D)--(W); \draw[thickline] (D)--(X); \draw[thickline] (W)--(X);
  \draw[thickline] (P)--(0.95,2); \draw[thickline] (1.85,2)--(Q);
  \draw[thickline] (0.95,1.74) rectangle (1.85,2.26);
  \node[above left] at (0.66,2.72) {$\tfrac12$};
  \node[above right] at (2.14,2.72) {$\tfrac12$};
  \node[above] at (1.4,2.72) {$\tfrac12$};
  \node[left] at (0.66,2.35) {$\tfrac12$};
  \node[right] at (2.14,2.35) {$\tfrac12$};
  \node[left] at (0.66,1.98) {$\tfrac12$};
  \node[right] at (2.14,1.98) {$\tfrac12$};
  \node at (1.09,2.36) {$\tfrac{1}{\sqrt2}$};
  \node at (1.4,2.0) {$\tfrac{x}{2}$};
  \node[left] at (0.66,1.55) {$\tfrac{2}{\sqrt2}$};
  \node[right] at (2.14,1.55) {$\tfrac{2}{\sqrt2}$};
  \node[below left] at (0,1.96) {$A$};
  \node[below right] at (2.8,1.96) {$B$};
  \node at (1.4,0.42) {(1)};
\end{scope}
% ---------- panel (2) ----------
\begin{scope}[xshift=4.6cm]
  \coordinate (A2) at (0,2);  \coordinate (B2) at (2.8,2);
  \coordinate (P2) at (0.7,2); \coordinate (Q2) at (2.1,2);
  \coordinate (U2) at (1.0,3.0); \coordinate (V2) at (1.8,3.0);
  \draw[thickline] (U2)--(V2);
  \draw[thickline] (U2)--(A2); \draw[thickline] (V2)--(B2);
  \draw[thickline] (U2)--(P2); \draw[thickline] (V2)--(Q2);
  \draw[thickline] (A2)--(P2); \draw[thickline] (Q2)--(B2);
  \draw[thickline] (P2)--(0.95,2); \draw[thickline] (1.85,2)--(Q2);
  \draw[thickline] (0.95,1.74) rectangle (1.85,2.26);
  \node[above] at (1.4,3.02) {$\tfrac14$};
  \node[above left] at (0.94,2.5) {$\tfrac14$};
  \node[above right] at (1.86,2.5) {$\tfrac14$};
  \node[left] at (0.74,2.4) {$\tfrac14$};
  \node[right] at (2.06,2.4) {$\tfrac14$};
  \node[left] at (0.66,1.98) {$\tfrac14$};
  \node[right] at (2.14,1.98) {$\tfrac14$};
  \node at (1.4,2.0) {$\tfrac{x}{2}$};
  \node[left] at (0.66,1.66) {$\tfrac{2}{\sqrt2}$};
  \node[right] at (2.14,1.66) {$\tfrac{2}{\sqrt2}$};
  \node[left] at (-0.06,1.98) {$A$};
  \node[right] at (2.86,1.98) {$B$};
  \node at (1.4,1.42) {(2)};
\end{scope}
% ---------- panel (3) ----------
\begin{scope}[xshift=9.2cm]
  \coordinate (A3) at (0,2);  \coordinate (B3) at (2.8,2);
  \coordinate (P3) at (0.7,2); \coordinate (Q3) at (2.1,2);
  \coordinate (U3) at (1.0,3.0); \coordinate (V3) at (1.8,3.0);
  \draw[thickline] (U3)--(V3);
  \draw[thickline] (U3)--(A3); \draw[thickline] (V3)--(B3);
  \draw[thickline] (U3)--(P3); \draw[thickline] (V3)--(Q3);
  \draw[thickline] (A3)--(P3); \draw[thickline] (Q3)--(B3);
  \draw[thickline] (P3)--(0.95,2); \draw[thickline] (1.85,2)--(Q3);
  \draw[thickline] (0.95,1.74) rectangle (1.85,2.26);
  \node[above] at (1.4,3.02) {$\tfrac{1}{2(1+\sqrt2)}$};
  \node[left] at (0.74,2.45) {$\tfrac14$};
  \node[right] at (2.06,2.45) {$\tfrac14$};
  \node[left] at (0.66,1.98) {$\tfrac14$};
  \node[right] at (2.14,1.98) {$\tfrac14$};
  \node at (1.4,2.0) {$\tfrac{x}{2}$};
  \node[left] at (0.66,1.66) {$\tfrac{2}{\sqrt2}$};
  \node[right] at (2.14,1.66) {$\tfrac{2}{\sqrt2}$};
  \node[left] at (-0.06,1.98) {$A$};
  \node[right] at (2.86,1.98) {$B$};
  \node at (1.4,1.42) {(3)};
\end{scope}
\end{tikzpicture}
\end{figbox}

Next fold the network along the line $AB$; coincident resistances are then in parallel, so their values are halved, giving Fig.~(2). Then combine the two topmost resistances in parallel, giving Fig.~(3).

At this point we obtain the equation
\[
x=\frac{1}{2}+\left[2(1+\sqrt2)+\frac{2\sqrt2}{\sqrt2x+1}\right]^{-1}.
\]
Simplifying,
\[
2x-1=\frac{1}{\sqrt2+1+\dfrac{\sqrt2}{\sqrt2x+1}},
\]
that is,
\[
(2x-1)\left[(2+\sqrt2)x+2\sqrt2+1\right]=\sqrt2x+1.
\]
Finally we obtain the quadratic equation for $x$:
\[
x^2+(\sqrt2-1)x-\frac{1}{\sqrt2}=0.
\]
Solving it, we obtain
\[
x=\frac{1}{2}\left(\sqrt3-\sqrt2+1\right)\approx0.659.
\]
\[
\therefore\ R_{AB}=\frac{1}{2}\left(\sqrt3-\sqrt2+1\right)R\approx0.659R.
\]
\end{solution}

\begin{solution}{3.29}
\textit{Method 1. Kirchhoff's equation method.}

By symmetry we may assume the currents in the branches as shown in Fig.~1, and write down the following equations.

Node equations:
\begin{align*}
I_2+I_4-I_1&=0,\\
I_3-I_4-I_5&=0.
\end{align*}
Loop equations:
\begin{align*}
I_2-I_3-I_4&=0,\\
I_3+I_5-I_4-I_1&=0.
\end{align*}

Solving the above four equations, we obtain
\[
I_2=\frac{5}{6}I_1,\qquad I_3=\frac{2}{3}I_1,\qquad
I_4=\frac{1}{6}I_1,\qquad I_5=\frac{1}{2}I_1.
\]

\begin{figbox}{3.29(a)}
\begin{tikzpicture}[scale=1.0,>=latex,
  car/.style={postaction={decorate,decoration={markings,
    mark=at position 0.5 with {\arrow{>}}}}}]
  \coordinate (L1) at (0,2.2);   % A
  \coordinate (L2) at (2.0,2.2);
  \coordinate (L3) at (0,0.6);   % D
  \coordinate (L4) at (2.0,0.6);
  \coordinate (R1) at (0.9,3.0); % C
  \coordinate (R2) at (2.9,3.0); % B
  \coordinate (R3) at (0.9,1.4); % F
  \coordinate (R4) at (2.9,1.4); % E
  \draw[thickline] (L1)--(L2)--(L4)--(L3)--cycle;
  \draw[thickline] (R1)--(R2)--(R4)--(R3)--cycle;
  \draw[thickline] (L1)--(R1); \draw[thickline] (L2)--(R2);
  \draw[thickline] (L3)--(R3); \draw[thickline] (L4)--(R4);
  \draw[thickline,car] (L1)--(R1);   % I2
  \draw[thickline,car] (R1)--(R2);   % I4
  \draw[thickline,car] (R1)--(R3);   % I5
  \draw[thickline,car] (L1)--(L3);   % I1
  \draw[thickline,car] (R3)--(R4);   % I3
  \draw[thickline,car] (L3)--(R3);   % I2
  \node[left] at (L1) {$A$};
  \node[above] at (R1) {$C$};
  \node[right] at (R2) {$B$};
  \node[left] at (L3) {$D$};
  \node[below right] at (R3) {$F$};
  \node[right] at (R4) {$E$};
  \node[above left] at (0.45,2.62) {$I_2$};
  \node[above] at (1.9,3.04) {$I_4$};
  \node[left] at (0.86,2.2) {$I_5$};
  \node[left] at (-0.02,1.4) {$I_1$};
  \node[above] at (1.9,1.44) {$I_3$};
  \node[below] at (0.45,0.86) {$I_2$};
  \draw[force] (-0.95,2.95)--(L1) node[midway,above left] {$I$};
  \draw[force] (R4)--(3.55,0.95) node[midway,below right] {$I$};
\end{tikzpicture}
\end{figbox}

Also,
\[
U_{AE}=I_1+I_2=\frac{11}{6}I_1,\qquad
I=I_1+I_2+I_3=\frac{15}{6}I_1,
\]
hence
\[
R_{AE}=\frac{U_{AE}}{I}=\frac{11}{15}\ \Omega.
\]

\textit{Method 2. Method of superposition of currents.}

Imagine a current $I$ flowing in at $A$ and out at $D$; the branch currents obtained are as shown in Fig.~2(1). Also imagine a current flowing in at $D$ and out at $E$; the branch currents obtained are as shown in Fig.~2(2). Superposing the above two cases, we obtain the case in which the current flows in at $A$ and out at $E$. Fig.~2(3) shows the currents in two of the branches after superposition (the remaining branches are not drawn). From this we can find
\[
R_{AE}:U_{AE}=\frac{2}{5}I+\frac{5}{15}I=\frac{11}{15}I.
\]
Therefore
\[
R_{AE}=U_{AE}/I=\frac{11}{15}\ \Omega.
\]

\begin{figbox}{3.29(b)}
\begin{tikzpicture}[scale=0.78,>=latex]
% ---- panel (1): current in at A, out at D ----
\begin{scope}
  \draw[thickline] (0,1.5)--(1.3,1.5)--(1.3,0.3)--(0,0.3)--cycle;
  \draw[thickline] (0.55,1.9)--(1.85,1.9)--(1.85,0.7)--(0.55,0.7)--cycle;
  \draw[thickline] (0,1.5)--(0.55,1.9); \draw[thickline] (1.3,1.5)--(1.85,1.9);
  \draw[thickline] (0,0.3)--(0.55,0.7); \draw[thickline] (1.3,0.3)--(1.85,0.7);
  \node[left] at (0,1.5) {$A$};
  \node[above] at (0.55,1.9) {$C$};
  \node[right] at (1.85,1.9) {$B$};
  \node[left] at (0,0.3) {$D$};
  \node[below right] at (0.55,0.7) {$F$};
  \node[right] at (1.85,0.7) {$E$};
  \node[above] at (0.28,1.72) {$\tfrac15I$};
  \node[left] at (-0.02,0.9) {$\tfrac{2}{5}I$};
  \node[below] at (0.9,0.26) {$\tfrac15I$};
  \draw[force] (-0.85,2.0)--(0,1.5) node[midway,above left] {$I$};
  \draw[force] (0,0.3)--(-0.7,-0.1) node[midway,below left] {$I$};
  \node at (0.65,-0.35) {(1)};
\end{scope}
% ---- panel (2): current in at D, out at E ----
\begin{scope}[xshift=3.9cm]
  \draw[thickline] (0,1.5)--(1.3,1.5)--(1.3,0.3)--(0,0.3)--cycle;
  \draw[thickline] (0.55,1.9)--(1.85,1.9)--(1.85,0.7)--(0.55,0.7)--cycle;
  \draw[thickline] (0,1.5)--(0.55,1.9); \draw[thickline] (1.3,1.5)--(1.85,1.9);
  \draw[thickline] (0,0.3)--(0.55,0.7); \draw[thickline] (1.3,0.3)--(1.85,0.7);
  \node[left] at (0,1.5) {$A$};
  \node[above] at (0.55,1.9) {$C$};
  \node[right] at (1.85,1.9) {$B$};
  \node[left] at (0,0.3) {$D$};
  \node[below right] at (0.55,0.7) {$F$};
  \node[right] at (1.85,0.7) {$E$};
  \node[above] at (0.3,1.72) {$\tfrac{1}{15}I$};
  \node[right] at (1.34,1.2) {$\tfrac15I$};
  \node[left] at (-0.02,0.9) {$\tfrac15I$};
  \node at (0.78,1.15) {$\tfrac{2}{15}I$};
  \node at (0.95,0.52) {$\tfrac{4}{15}I$};
  \node[below] at (1.0,0.24) {$\tfrac{8}{15}I$};
  \draw[force] (-0.72,-0.08)--(0,0.3) node[midway,below left] {$I$};
  \draw[force] (1.85,0.7)--(2.4,0.25) node[midway,below right] {$I$};
  \node at (0.65,-0.35) {(2)};
\end{scope}
% ---- panel (3): superposition ----
\begin{scope}[xshift=7.6cm]
  \draw[thickline] (0,1.5)--(1.3,1.5)--(1.3,0.3)--(0,0.3)--cycle;
  \draw[thickline] (0.55,1.9)--(1.85,1.9)--(1.85,0.7)--(0.55,0.7)--cycle;
  \draw[thickline] (0,1.5)--(0.55,1.9); \draw[thickline] (1.3,1.5)--(1.85,1.9);
  \draw[thickline] (0,0.3)--(0.55,0.7); \draw[thickline] (1.3,0.3)--(1.85,0.7);
  \node[left] at (0,1.5) {$A$};
  \node[above] at (0.55,1.9) {$C$};
  \node[right] at (1.85,1.9) {$B$};
  \node[left] at (0,0.3) {$D$};
  \node[below right] at (0.55,0.7) {$F$};
  \node[right] at (1.85,0.7) {$E$};
  \node[left] at (-0.02,0.9) {$\tfrac25I$};
  \node[below] at (1.15,0.24) {$\tfrac{5}{15}I$};
  \draw[force] (-0.85,2.0)--(0,1.5) node[midway,above left] {$I$};
  \draw[force] (1.85,0.7)--(2.4,0.25) node[midway,below right] {$I$};
  \node at (0.65,-0.35) {(3)};
\end{scope}
\end{tikzpicture}
\end{figbox}

\textit{Method 3. Star-delta transformation method.}

The actual transformation process is shown in Fig.~3. Finally we obtain

\begin{figbox}{3.29(c)}
\begin{tikzpicture}[scale=0.72,>=latex]
% ---- (1) the original cube ----
\begin{scope}
  \draw[thickline] (0,1.5)--(1.3,1.5)--(1.3,0.3)--(0,0.3)--cycle;
  \draw[thickline] (0.55,1.9)--(1.85,1.9)--(1.85,0.7)--(0.55,0.7)--cycle;
  \draw[thickline] (0,1.5)--(0.55,1.9); \draw[thickline] (1.3,1.5)--(1.85,1.9);
  \draw[thickline] (0,0.3)--(0.55,0.7); \draw[thickline] (1.3,0.3)--(1.85,0.7);
  \node[left] at (0,1.5) {$A$};
  \node[right] at (1.85,0.7) {$E$};
  \node at (0.65,-0.3) {(1)};
\end{scope}
\draw[force] (2.1,0.9)--(2.6,0.9);
\node at (2.35,1.12) {\small star--delta};
% ---- (2) ----
\begin{scope}[xshift=3.4cm]
  \draw[thickline] (0,1.5)--(0.55,1.9)--(1.85,1.9)--(1.3,1.5)--cycle;
  \draw[thickline] (0,1.5)--(0,0.3); \draw[thickline] (1.3,1.5)--(1.3,0.3);
  \draw[thickline] (0,0.3)--(0.55,0.7)--(1.85,0.7)--(1.3,0.3);
  \draw[thickline] (0.3,1.5)--(0.3,0.5); \draw[thickline] (1.0,1.5)--(1.0,0.5);
  \node[left] at (0,1.5) {$A$};
  \node[right] at (1.3,0.3) {$E$};
  \node[above] at (0.28,1.72) {$\tfrac13$};
  \node[above] at (1.0,1.72) {$\tfrac13$};
  \node[above] at (1.58,1.72) {$\tfrac13$};
  \node[left] at (-0.02,1.05) {$1$};
  \node[left] at (-0.02,0.6) {$1$};
  \node[below] at (0.28,0.48) {$\tfrac13$};
  \node[below] at (1.0,0.48) {$\tfrac13$};
  \node[below] at (1.58,0.48) {$\tfrac13$};
  \node[right] at (1.34,0.9) {$1$};
  \node at (0.65,-0.3) {(2)};
\end{scope}
\draw[force] (5.15,0.9)--(5.6,0.9);
\node at (5.37,1.12) {\small simplify};
% ---- (3) bridge ----
\begin{scope}[xshift=6.4cm]
  \coordinate (a) at (0,0.9); \coordinate (e) at (2.0,0.9);
  \coordinate (t) at (1.0,1.8); \coordinate (b) at (1.0,0);
  \draw[thickline] (a)--(t)--(e); \draw[thickline] (a)--(b)--(e);
  \draw[thickline] (t)--(b);
  \node[left] at (a) {$A$};
  \node[right] at (e) {$E$};
  \node[above left] at (0.5,1.35) {$\tfrac13$};
  \node[above right] at (1.5,1.35) {$\tfrac43$};
  \node[below left] at (0.5,0.45) {$\tfrac43$};
  \node[below right] at (1.5,0.45) {$\tfrac13$};
  \node[right] at (1.03,1.0) {$\tfrac53$};
  \node at (1.0,-0.3) {(3)};
\end{scope}
\draw[force] (8.6,0.9)--(9.05,0.9);
\node at (8.82,1.12) {\small star--delta};
% ---- (4) ----
\begin{scope}[xshift=9.9cm]
  \coordinate (a) at (0,0.9); \coordinate (e) at (1.8,0.9);
  \coordinate (t) at (0.9,1.7); \coordinate (b) at (0.9,0.1);
  \draw[thickline] (a)--(t)--(e); \draw[thickline] (a)--(b)--(e);
  \draw[thickline] (t)--(b);
  \node[left] at (a) {$A$};
  \node[right] at (e) {$E$};
  \node[above left] at (0.45,1.28) {$\tfrac{2}{15}$};
  \node[above right] at (1.35,1.28) {$\tfrac43$};
  \node[below left] at (0.45,0.52) {$\tfrac23$};
  \node[below right] at (1.35,0.52) {$\tfrac13$};
  \node[right] at (0.93,1.0) {$\tfrac16$};
  \node at (0.9,-0.3) {(4)};
\end{scope}
\draw[force] (11.95,0.9)--(12.4,0.9);
\node at (12.17,1.12) {\small simplify};
% ---- (5) ----
\begin{scope}[xshift=13.2cm]
  \draw[thickline] (0,0.9)--(0.6,0.9);
  \draw[thickline] (0.6,0.75) rectangle (1.2,1.05);
  \node at (0.9,0.9) {$\tfrac{2}{15}$};
  \draw[thickline] (1.2,0.9)--(1.7,0.9);
  \draw[thickline] (1.7,0.75) rectangle (2.3,1.05);
  \node at (2.0,0.9) {$\tfrac32$};
  \draw[thickline] (2.3,0.9)--(2.7,0.9);
  \draw[thickline] (2.7,0.8) ellipse (0.18 and 0.1);
  \node[above] at (2.7,0.98) {$1$};
  \draw[thickline] (2.88,0.9)--(3.2,0.9);
  \node[left] at (0,0.9) {$A$};
  \node[right] at (3.2,0.9) {$E$};
  \node at (1.6,-0.3) {(5)};
\end{scope}
\end{tikzpicture}
\end{figbox}

\[
R_{AE}=\frac{2}{15}+\frac{\dfrac{3}{2}\cdot1}{\dfrac{3}{2}+1}=\frac{11}{15}\ \Omega.
\]
\end{solution}

\begin{solution}{3.30}
According to the statement of the problem, the circuit shown in the figure may be drawn. The current flowing through the galvanometer may be found by Th\'evenin's theorem (note that quantities of order $\delta^2$ and higher are neglected in the calculation). Imagine that what is connected across the galvanometer is an equivalent source, whose electromotive force is
\[
\varepsilon_{\text{eff}}=(U_A-U_B)_{A,\,B\ \text{open}}.
\]
\begin{align*}
&=U\frac{R_2}{R_1+R_2}-U\frac{R_2(1-\delta)}{R_1+R_2-R_2\delta}\\
&=U\frac{R_2}{R_1+R_2}-\frac{UR_2}{R_1+R_2}\frac{(1-\delta)}{1-\delta\dfrac{R_2}{R_1+R_2}}\\
&\approx U\frac{R_2}{R_1+R_2}-U\frac{R_2}{R_1+R_2}(1-\delta)
\left(1+\delta\frac{R_2}{R_1+R_2}\right)\\
&\approx\delta U\frac{R_1R_2}{(R_1+R_2)^2}.
\end{align*}

\begin{figbox}{3.30}
\begin{tikzpicture}[scale=0.9,>=latex]
  \coordinate (A) at (2.6,2.2);
  \coordinate (B) at (2.6,0.6);
  \coordinate (L) at (0.6,1.4);
  \coordinate (R) at (4.6,1.4);
  \draw[thickline] (A)--(L)--(B)--(R)--(A);
  \draw[thickline] (A)--(B);
  % resistors
  \draw[thickline] (1.35,1.98) rectangle (1.95,1.62);
  \node at (1.65,1.8) {$R_1$};
  \draw[thickline] (3.25,1.98) rectangle (3.85,1.62);
  \node at (3.55,1.8) {$R_1$};
  \draw[thickline] (1.4,0.85) rectangle (1.8,1.15);
  \node at (1.6,1.0) {$R_2$};
  \draw[thickline] (3.95,1.05) rectangle (4.35,0.75);
  \node at (4.15,0.9) {$R_2(1-\delta)$};
  % galvanometer
  \draw[thickline] (2.6,1.65) circle (0.17);
  \node at (2.6,1.65) {$G$};
  \node[right] at (2.72,1.9) {$I_G$};
  \node[above] at (A) {$A$};
  \node[below] at (B) {$B$};
  % battery and current
  \draw[thickline] (L)--(0.6,-0.5)--(4.6,-0.5)--(R);
  \draw[thickline] (2.45,-0.62)--(2.45,-0.38);
  \draw[thickline] (2.62,-0.55)--(2.62,-0.45);
  \node[below] at (2.6,-0.62) {$U$};
  \draw[force] (0.6,-0.05)--(0.6,0.55) node[midway,left] {$I$};
\end{tikzpicture}
\end{figbox}

Since $\varepsilon_{\text{eff}}$ is already a small quantity, the internal resistance of the equivalent source may be replaced by its value when the bridge is balanced:
\[
r_{\text{eff}}=2\frac{R_1R_2}{R_1+R_2},
\]
\[
\therefore\ I_G=\frac{\varepsilon_{\text{eff}}}{r_{\text{eff}}}
=\frac{\delta U\dfrac{R_1R_2}{(R_1+R_2)^2}}{2\dfrac{R_1R_2}{R_1+R_2}}
=\frac{\delta U}{2(R_1+R_2)}.
\]
Since $I_G$ is a first-order small quantity, the current flowing through the source may be replaced by its value when the bridge is balanced:
\[
I=\frac{U}{\dfrac{1}{2}(R_1+R_2)}=\frac{2U}{R_1+R_2},
\]
\[
\therefore\ I_G/I=\frac{\delta U}{2(R_1+R_2)}\Big/\frac{2U}{R_1+R_2}=\frac{\delta}{4}.
\]
\end{solution}

\begin{solution}{3.31}
(1) According to Th\'evenin's theorem, looking to the left from $A_n$, the network may be replaced by an equivalent source with electromotive force $V_n/2$ and internal resistance $R$, as shown in Fig.~1. Connecting this equivalent source with the next section, it may again be replaced by an equivalent source with electromotive force $\left(\dfrac{V_{n-1}}{2}+\dfrac{V_n}{4}\right)$ and internal resistance $R$. This is also the equivalent source seen looking to the left from $A_{n-1}$. And so on: when $A$ and $B$ are taken as the output terminals, the network may still be replaced by an equivalent source, whose electromotive force $V$ and internal resistance $r$ are respectively
\[
V=\frac{V_1}{2}+\frac{V_2}{2^2}+\frac{V_3}{2^3}+\cdots+\frac{V_n}{2^n},
\qquad r=R.
\]

\begin{figbox}{3.31(a)}
\begin{tikzpicture}[scale=0.85,>=latex]
% ---------- left: one section ----------
\begin{scope}
  \draw[thickline] (0,2.1)--(2.4,2.1);
  \draw[thickline] (0.35,2.25) rectangle (1.15,1.95);
  \node at (0.75,2.1) {$2R$};
  \node[right] at (2.45,2.1) {$A_n$};
  \draw[thickline] (1.9,2.1)--(1.9,0.9);
  \draw[thickline] (1.75,1.7) rectangle (2.05,1.1);
  \node[right] at (2.08,1.4) {$2R$};
  \draw[thickline] (1.9,0.9)--(1.9,0.45);
  \draw[thickline] (1.72,0.45)--(2.08,0.45);
  \draw[thickline] (1.78,0.32)--(2.02,0.32);
  \node[right] at (2.14,0.38) {$V_n$};
  \draw[thickline] (1.9,0.32)--(1.9,0)--(0,0)--(0,2.1);
  \draw[force] (0,-0.02)--(0,0.55) node[midway,left] {$I$};
  \node at (1.2,-0.35) {(1)};
\end{scope}
\draw[force] (3.1,1.1)--(3.7,1.1);
% ---------- right: equivalent ----------
\begin{scope}[xshift=4.4cm]
  \draw[thickline] (1.2,2.1)--(2.3,2.1);
  \node[right] at (2.32,2.1) {$A_n$};
  \draw[thickline] (1.2,2.1)--(1.2,1.5);
  \draw[thickline] (1.05,1.5) rectangle (1.35,0.95);
  \node[right] at (1.38,1.22) {$R$};
  \draw[thickline] (1.2,0.95)--(1.2,0.55);
  \draw[thickline] (1.02,0.55)--(1.38,0.55);
  \draw[thickline] (1.08,0.42)--(1.32,0.42);
  \node[right] at (1.44,0.48) {$\tfrac{V_n}{2}$};
  \draw[thickline] (1.2,0.42)--(1.2,0)--(0,0)--(0,2.1)--(1.2,2.1);
  \node at (1.2,-0.35) {(2)};
\end{scope}
\end{tikzpicture}
\end{figbox}

as shown in Fig.~2.

\begin{figbox}{3.31(b)}
\begin{tikzpicture}[scale=0.62,>=latex]
% ---------- panel 1 ----------
\begin{scope}
  \draw[thickline] (0,2.4)--(0.5,2.4);
  \draw[thickline] (0.5,2.55) rectangle (1.2,2.25);
  \node at (0.85,2.4) {$2R$};
  \draw[thickline] (1.2,2.4)--(1.9,2.4);
  \node[above] at (1.9,2.42) {$A_n$};
  \draw[thickline] (1.9,2.4)--(2.5,2.4);
  \draw[thickline] (2.5,2.55) rectangle (3.1,2.25);
  \node at (2.8,2.4) {$R$};
  \draw[thickline] (3.1,2.4)--(3.7,2.4);
  \node[above] at (3.7,2.42) {$A_{n-1}$};
  % shunts
  \draw[thickline] (1.9,2.4)--(1.9,1.55);
  \draw[thickline] (1.75,1.95) rectangle (2.05,1.55);
  \node[right] at (2.08,1.75) {$2R$};
  \draw[thickline] (1.9,1.55)--(1.9,1.15);
  \draw[thickline] (1.72,1.15)--(2.08,1.15);
  \draw[thickline] (1.78,1.02)--(2.02,1.02);
  \node[right] at (2.14,1.08) {$V_n$};
  \draw[thickline] (1.9,1.02)--(1.9,0.6);
  \draw[thickline] (3.7,2.4)--(3.7,1.55);
  \draw[thickline] (3.55,1.95) rectangle (3.85,1.55);
  \node[right] at (3.88,1.75) {$2R$};
  \draw[thickline] (3.7,1.55)--(3.7,1.15);
  \draw[thickline] (3.52,1.15)--(3.88,1.15);
  \draw[thickline] (3.58,1.02)--(3.82,1.02);
  \node[right] at (3.94,1.08) {$V_{n-1}$};
  \draw[thickline] (3.7,1.02)--(3.7,0.6);
  \draw[thickline] (0,2.4)--(0,0.6)--(3.7,0.6);
  \node at (1.8,0.15) {(1)};
\end{scope}
\draw[force] (4.35,1.4)--(4.9,1.4);
% ---------- panel 2 ----------
\begin{scope}[xshift=5.3cm]
  \draw[thickline] (0,2.4)--(0.55,2.4);
  \node[above] at (0.55,2.42) {$A_n$};
  \draw[thickline] (0.55,2.4)--(1.15,2.4);
  \draw[thickline] (1.15,2.55) rectangle (1.75,2.25);
  \node at (1.45,2.4) {$R$};
  \draw[thickline] (1.75,2.4)--(2.35,2.4);
  \node[above] at (2.35,2.42) {$A_{n-1}$};
  \draw[thickline] (0.55,2.4)--(0.55,1.55);
  \draw[thickline] (0.40,1.95) rectangle (0.70,1.55);
  \node[right] at (0.73,1.75) {$R$};
  \draw[thickline] (0.55,1.55)--(0.55,1.15);
  \draw[thickline] (0.37,1.15)--(0.73,1.15);
  \draw[thickline] (0.43,1.02)--(0.67,1.02);
  \node[left] at (0.33,1.08) {$\tfrac{V_n}{2}$};
  \draw[thickline] (0.55,1.02)--(0.55,0.6);
  \draw[thickline] (2.35,2.4)--(2.35,1.55);
  \draw[thickline] (2.20,1.95) rectangle (2.50,1.55);
  \node[right] at (2.53,1.75) {$2R$};
  \draw[thickline] (2.35,1.55)--(2.35,1.15);
  \draw[thickline] (2.17,1.15)--(2.53,1.15);
  \draw[thickline] (2.23,1.02)--(2.47,1.02);
  \node[right] at (2.59,1.08) {$V_{n-1}$};
  \draw[thickline] (2.35,1.02)--(2.35,0.6);
  \draw[thickline] (0,2.4)--(0,0.6)--(2.35,0.6);
  \node at (1.2,0.15) {(2)};
\end{scope}
\draw[force] (8.0,1.4)--(8.55,1.4);
% ---------- panel 3 ----------
\begin{scope}[xshift=8.95cm]
  \draw[thickline] (0,2.4)--(1.0,2.4);
  \node[above] at (1.0,2.42) {$A_{n-1}$};
  \draw[thickline] (0,2.4)--(0,1.55);
  \draw[thickline] (-0.15,1.95) rectangle (0.15,1.55);
  \node[right] at (0.18,1.75) {$R$};
  \draw[thickline] (0,1.55)--(0,1.15);
  \draw[thickline] (-0.18,1.15)--(0.18,1.15);
  \draw[thickline] (-0.12,1.02)--(0.12,1.02);
  \node[right] at (0.24,1.08) {$\tfrac{V_{n-1}}{2}+\tfrac{V_n}{4}$};
  \draw[thickline] (0,1.02)--(0,0.6)--(1.0,0.6);
  \node at (0.5,0.15) {(3)};
\end{scope}
\draw[force] (10.6,1.4)--(11.15,1.4);
\node at (10.87,1.68) {$\cdots$};
% ---------- panel 4 ----------
\begin{scope}[xshift=12.0cm]
  \draw[thickline] (0,2.4)--(1.2,2.4);
  \node[above] at (0,2.42) {$A_1$};
  \node[right] at (1.22,2.4) {$A$};
  \draw[thickline] (0,2.4)--(0,1.55);
  \draw[thickline] (-0.15,1.95) rectangle (0.15,1.55);
  \node[right] at (0.18,1.75) {$R$};
  \draw[thickline] (0,1.55)--(0,1.15);
  \draw[thickline] (-0.18,1.15)--(0.18,1.15);
  \draw[thickline] (-0.12,1.02)--(0.12,1.02);
  \node[right] at (0.24,1.02) {$\tfrac{V_1}{2}+\tfrac{V_2}{4}+\cdots+\tfrac{V_n}{2^n}$};
  \draw[thickline] (0,1.02)--(0,0.6)--(1.2,0.6);
  \node[right] at (1.22,0.6) {$B$};
  \node at (0.6,0.15) {(4)};
\end{scope}
\end{tikzpicture}
\end{figbox}

(2) At this time, looking to the left from $A_i$ gives an equivalent source, and looking to the right from $A_{i-1}$ also gives an equivalent source; their electromotive forces are $V_{\text{L}}$ and $V_{\text{R}}$ respectively, as shown in Fig.~3. Therefore which of the potentials of the two points $A_i$ and $A_{i-1}$ is the higher depends on which of $V_{\text{L}}$ and $V_{\text{R}}$ is the larger. From the results of part (1) it is known that when $i\le6$ we must have $V_{\text{L}}>V_{\text{R}}$, because the electromotive force of any battery on the left is higher than that of a battery on the right. Hence the point of highest potential must be among the points $A_6,A_7,\cdots,A_{10}$. If we can find the point $A_i$ with the smallest ordinal number $i$ such that $V_{\text{L}}<V_{\text{R}}$, then $A_{i-1}$ at this time must be the point of highest potential.

\begin{figbox}{3.31(c)}
\begin{tikzpicture}[scale=0.62,>=latex]
% ---------- left: the whole ladder to the left of A_i ----------
\begin{scope}
  \draw[dashedline] (-0.35,0.15) rectangle (7.6,3.1);
  \draw[thickline] (0.1,2.6)--(0.6,2.6);
  \draw[thickline] (0.6,2.75) rectangle (1.2,2.45);
  \node at (0.9,2.6) {$2R$};
  \draw[thickline] (1.2,2.6)--(1.8,2.6);
  \node[above] at (1.8,2.62) {$A_n$};
  \draw[thickline] (1.8,2.6)--(2.3,2.6);
  \draw[thickline] (2.3,2.75) rectangle (2.9,2.45);
  \node at (2.6,2.6) {$R$};
  \draw[thickline] (2.9,2.6)--(3.4,2.6);
  \node[above] at (3.4,2.62) {$A_{n-1}$};
  \draw[thickline] (4.1,2.6)--(4.6,2.6);
  \node[above] at (4.6,2.62) {$A_i$};
  \draw[dashedline] (3.55,2.6)--(4.0,2.6);
  % shunts
  \foreach \x/\lab in {1.8/{$V_n$},3.4/{$V_{n-1}$},4.6/{$V_i$}}{
    \draw[thickline] (\x,2.6)--(\x,1.85);
    \draw[thickline] (\x-0.15,1.85) rectangle (\x+0.15,1.45);
    \draw[thickline] (\x,1.45)--(\x,1.15);
    \draw[thickline] (\x-0.18,1.15)--(\x+0.18,1.15);
    \draw[thickline] (\x-0.12,1.02)--(\x+0.12,1.02);
    \node[below] at (\x,0.98) {\lab};
    \draw[thickline] (\x,1.02)--(\x,0.55);
  }
  \node[right] at (1.95,1.65) {$2R$};
  \node[right] at (3.55,1.65) {$2R$};
  \node[right] at (4.75,1.65) {$2R$};
  \draw[thickline] (0.1,2.6)--(0.1,0.55)--(4.6,0.55);
\end{scope}
% ---------- middle: one section ----------
\begin{scope}[xshift=8.3cm]
  \draw[dashedline] (-0.2,0.15) rectangle (3.5,3.1);
  \draw[thickline] (0.1,2.6)--(0.6,2.6);
  \node[above] at (0.6,2.62) {$A_{i-1}$};
  \draw[thickline] (0.6,2.6)--(1.1,2.6);
  \draw[thickline] (1.1,2.75) rectangle (1.7,2.45);
  \node at (1.4,2.6) {$R$};
  \draw[thickline] (1.7,2.6)--(2.3,2.6);
  \node[above] at (2.3,2.62) {$A_i$};
  \draw[thickline] (2.3,2.6)--(2.3,1.85);
  \draw[thickline] (2.15,1.85) rectangle (2.45,1.45);
  \node[right] at (2.5,1.65) {$2R$};
  \draw[thickline] (2.3,1.45)--(2.3,1.15);
  \draw[thickline] (2.12,1.15)--(2.48,1.15);
  \draw[thickline] (2.18,1.02)--(2.42,1.02);
  \node[below] at (2.3,0.98) {$V_i$};
  \draw[thickline] (2.3,1.02)--(2.3,0.55);
  \draw[thickline] (0.1,2.6)--(0.1,0.55)--(2.3,0.55);
\end{scope}
\draw[force] (11.9,1.5)--(12.4,1.5);
% ---------- right: the equivalent circuit ----------
\begin{scope}[xshift=12.9cm]
  \draw[thickline] (0.1,2.6)--(0.9,2.6);
  \node[above] at (0.1,2.62) {$A_i$};
  \draw[thickline] (0.9,2.75) rectangle (1.5,2.45);
  \node at (1.2,2.6) {$R$};
  \draw[thickline] (1.5,2.6)--(2.3,2.6);
  \node[above] at (2.3,2.62) {$A_{i-1}$};
  \draw[thickline] (0.1,2.6)--(0.1,1.8);
  \draw[thickline] (-0.05,1.8) rectangle (0.25,1.4);
  \node[left] at (-0.12,1.6) {$R$};
  \draw[thickline] (0.1,1.4)--(0.1,1.1);
  \draw[thickline] (-0.08,1.1)--(0.28,1.1);
  \draw[thickline] (-0.02,0.97)--(0.22,0.97);
  \node[below] at (0.1,0.93) {$V_{\text{L}}$};
  \draw[thickline] (0.1,0.97)--(0.1,0.5);
  \draw[thickline] (2.3,2.6)--(2.3,1.8);
  \draw[thickline] (2.15,1.8) rectangle (2.45,1.4);
  \node[right] at (2.52,1.6) {$R$};
  \draw[thickline] (2.3,1.4)--(2.3,1.1);
  \draw[thickline] (2.12,1.1)--(2.48,1.1);
  \draw[thickline] (2.18,0.97)--(2.42,0.97);
  \node[below] at (2.3,0.93) {$V_{\text{R}}$};
  \draw[thickline] (2.3,0.97)--(2.3,0.5);
  \draw[thickline] (0.1,0.5)--(2.3,0.5);
\end{scope}
\end{tikzpicture}
\end{figbox}

Let us first look at $A_8$, the middle point of $6\sim10$. At this time
\[
V_{\text{L}}=\frac{18}{2}+\frac{19}{4}+\frac{20}{8},\qquad
V_{\text{R}}=\frac{17}{2}+\frac{16}{4}+\frac{15}{8}+\frac{14}{16}
+\frac{13}{32}+\frac{12}{64}+\frac{11}{128},
\]
then
\[
V_{\text{L}}-V_{\text{R}}=\frac{1}{2}+\frac{3}{4}+\frac{5}{8}
-\left(\frac{14}{16}+\frac{13}{32}+\frac{12}{64}+\frac{11}{128}\right)
=\frac{41}{128}>0.
\]
Now look at $A_9$. At this time
\[
V_{\text{L}}=\frac{19}{2}+\frac{20}{4},\qquad
V_{\text{R}}=\frac{18}{2}+\frac{17}{4}+\frac{16}{8}+\frac{15}{16}
+\frac{14}{32}+\frac{13}{64}+\frac{12}{128}+\frac{11}{256},
\]
\[
V_{\text{L}}-V_{\text{R}}=\frac{1}{2}+\frac{3}{4}
-\left(\frac{16}{8}+\frac{15}{16}+\frac{14}{32}+\frac{13}{64}
+\frac{12}{128}+\frac{11}{256}\right)=-\frac{631}{256}<0.
\]
It can be seen that the potential of $A_8$ is the highest. Its potential $U_8$ may be found from the corresponding equivalent circuit (the right-hand part of Fig.~3):
\begin{align*}
U_8&=\frac{V_{\text{L}}-V_{\text{R}}}{3R}\cdot R+V_{\text{R}}
=\frac{1}{3}V_{\text{L}}+\frac{2}{3}V_{\text{R}}\\
&=\frac{1}{3}\left(\frac{19}{2}+\frac{20}{4}\right)
+\frac{2}{3}\left(\frac{18}{2}+\frac{17}{4}+\frac{16}{8}+\cdots
+\frac{12}{128}+\frac{11}{256}\right)
=\frac{6199}{384}=16.14\ (\text{V}).
\end{align*}
\end{solution}

\begin{solution}{3.32}
(1) It can be determined. This can be judged from whether the direction of the current in the ring is the same as or opposite to the direction of the velocity of the ring: if they are the same, the ring carries positive charge; otherwise it carries negative charge. And the direction of the current can be determined from its magnetic effect.

(2) During the deceleration of the ring, the charge carriers will be acted on by an inertial force, and this inertial force may be regarded as a force originating from a non-electrostatic source. Let the linear velocity of the ring be $v$, and let the field strength of the field corresponding to this force originating from a non-electrostatic source be $E$. Let the mass and the charge of a charge carrier be $m$ and $e$ respectively; then we have
\[
-m\frac{\dd v}{\dd t}=eE,
\]
\[
E=-\frac{m}{e}\frac{\dd v}{\dd t}.
\]
The equivalent electromotive force is
\[
\mathcal{E}=E\cdot2\pi b=-\frac{m}{e}\frac{\dd v}{\dd t}\cdot2\pi b.
\]
By Ohm's law,
\[
\mathcal{E}=IR,
\]
that is,
\[
-\frac{m}{e}\frac{\dd v}{\dd t}\cdot2\pi b=IR.
\]
Integrating both sides with respect to $t$,
\[
-\frac{m}{e}\cdot2\pi b\int_{v_0}^{0}\dd v=R\int_0^\infty I\,\dd t=Rq,
\]
that is,
\[
\frac{m}{e}\cdot2\pi bv_0=Rq,
\]
\[
\therefore\ \frac{e}{m}=\frac{2\pi bv_0}{Rq}.
\]
\end{solution}

\begin{solution}{3.33}
Since the bubble does not conduct electricity, the current will accumulate charge on the surface of the bubble, as shown in the figure. Let the surface charge density be
\[
\sigma=-\sigma_0\cos\theta
\]
(only in this way can $E_r=0$ hold on the side of the surface facing the wire). Then the electric field inside the bubble is
\[
\vec{E}_{\text{in}}=\left(E_0+\frac{\sigma_0}{3\varepsilon_0}\right)\hat{\imath}
\]
(see Solution 3.24). Here $E_0$ is the electric field in the wire when the bubble is absent, that is, $E_0=\dfrac{1}{g}j_0$, where $g$ is the conductivity. Outside the bubble the field produced by the surface charges is the field of a dipole. The dipole moment is
\[
\vec{p}=\frac{4}{3}\pi a^3\sigma_0(-\hat{\imath})
\]
At the outer boundary of the bubble $(a,\theta)$ (which lies inside the wire) we must have $E_r=0$, that is,
\[
E_0\cos\theta+\frac{2p\cos\theta}{4\pi\varepsilon_0a^3}=0.
\]
Substituting the expression for $p$, we obtain
\[
E_0\cos\theta-\frac{2\cos\theta}{4\pi\varepsilon_0a^3}\cdot\frac{4\pi}{3}a^3\sigma_0=0,
\]
that is,
\[
E_0\cos\theta-\frac{2}{3\varepsilon_0}\sigma_0\cos\theta=0.
\]
Hence
\[
\sigma_0=\frac{3}{2}\varepsilon_0E_0,\qquad
\vec{p}=-2\pi\varepsilon_0a^3E_0=-\frac{2\pi\varepsilon_0a^3}{g}\vec{j}_0 .
\]

\begin{figbox}{3.33}
\begin{tikzpicture}[scale=1.0,>=latex]
  % bubble
  \draw[thickline] (0,0) circle (1.5);
  % accumulated surface charge
  \node at (-1.00,0.62) {$+$};
  \node at (-1.36,0.00) {$+$};
  \node at (-1.00,-0.62) {$+$};
  \node at (1.00,0.62) {$-$};
  \node at (1.36,0.00) {$-$};
  \node at (1.00,-0.62) {$-$};
  % current density in the wire
  \draw[force] (-2.75,1.15)--(-1.75,1.15) node[midway,above] {$\vec{j}_0$};
  % x axis
  \draw[axis] (0,0)--(2.45,0) node[right] {$x$};
  % radius vector to the field point P
  \draw[force] (0,0)--(47:1.5) node[above right] {$P$};
  \node at (0.50,0.68) {$r$};
  \draw[thickline] (0.75,0) arc (0:47:0.75);
  \node at (0.92,0.30) {$\theta$};
\end{tikzpicture}
\end{figbox}

(1) The electric field inside the bubble is
\[
\vec{E}_{\text{in}}=\left(E_0+\frac{\sigma_0}{3\varepsilon_0}\right)\hat{\imath}
=\frac{3}{2}\frac{1}{g}\vec{j}_0 .
\]

(2) The field at the point $P$ is the superposition of $E_0$ and the field of the dipole $p$:
\begin{align}
E_r &= E_0\cos\theta+\frac{2p\cos\theta}{4\pi\varepsilon_0r^3}
=E_0\cos\theta-\frac{a^3}{r^3}\cos\theta\,E_0\\
E_\theta &= -E_0\sin\theta+\frac{p\sin\theta}{4\pi\varepsilon_0r^3}
=-E_0\sin\theta-\frac{a^3}{2r^3}\sin\theta\,E_0
\end{align}
Noting that $E_0=\dfrac{1}{g}j_0$,
\begin{align}
j_r=gE_r &= j_0\cos\theta-j_0\left(\frac{a^3}{r^3}\right)\cos\theta
=j_0\left(1-\frac{a^3}{r^3}\right)\cos\theta\\
j_\theta=gE_\theta &= -j_0\left(1+\frac{a^3}{2r^3}\right)\sin\theta
\end{align}

(3) When $\theta=90\degree$, $j_r=0$ and $j_\theta=-j_0\left(1+\dfrac{a^3}{2r^3}\right)$, so
\[
I=\int_a^b -j_\theta\cdot2\pi r\dd r
=\int_a^b j_0\left(1+\frac{a^3}{2r^3}\right)2\pi r\dd r
\]
\[
=j_0\pi\left(b^2-a^2\right)+\int_a^b\pi j_0a^3\frac{\dd r}{r^2}
\]
\[
=j_0\pi b^2-j_0\pi a^2\frac{a}{b}.
\]
\end{solution}

\begin{solution}{3.34}
(1) The current (strength) may be regarded as steady, and the current density $j$ may be regarded as depending only on $r$; hence
\[
4\pi jr^2=I\ (\text{a constant})=j_0\cdot4\pi R^2=\sigma_0E_0\cdot4\pi R^2.
\]
\[
\therefore\ I=\sigma_0E_0\cdot4\pi R^2=3\times10^{-14}\times100\times4\pi\times\left(6.4\times10^6\right)^2=1.54\times10^3\ (\text{A})
\]
\[
j=\frac{R^2}{r^2}j_0=\frac{\sigma_0E_0R^2}{r^2}=\sigma E.
\]
\[
\therefore\ E(r)=\frac{j}{\sigma}=\frac{\sigma_0E_0R^2}{\sigma r^2}
=\frac{\sigma_0E_0R^2}{r^2\left[\sigma_0+a(r-R)^2\right]},
\]
the direction of $\vec{E}$ being towards the centre of the Earth.

(2) Take the thin shell of radii $r\sim r+\Delta r$. By Gauss's law,
\[
E(r+\Delta r)\cdot4\pi(r+\Delta r)^2-E(r)\cdot4\pi r^2=-4\pi r^2\Delta r\,\rho/\varepsilon_0,
\]
\[
\therefore\ \rho=-\frac{\varepsilon_0}{4\pi r^2}\frac{\dd}{\dd r}\left[E(r)\cdot4\pi r^2\right]
=\frac{2\varepsilon_0\sigma_0E_0R^2a(r-R)}{r^2\left[\sigma_0+a(r-R)^2\right]^2}.
\]
From $E_0=-\Sigma_0/\varepsilon_0$ we obtain
\[
\Sigma_0=-\varepsilon_0E_0=-\frac{100}{4\pi\times9\times10^9}=-8.85\times10^{-10}\ \left(\text{C/m}^2\right).
\]

(3) The potential at the Earth's surface is
\[
\varphi=\int_R^\infty -E(r)\dd r=-\sigma_0E_0R^2\int_R^\infty\frac{\dd r}{r^2\left[\sigma_0+a(r-R)^2\right]}.
\]
By the integral formula given in the problem,
\[
\int_R^\infty\frac{\dd r}{r^2\left[\sigma_0+a(r-R)^2\right]}
=\int_R^\infty\frac{\dd r}{ar^2\left[\frac{\sigma_0}{a}+R^2-2Rr+r^2\right]}
\]
\begin{align}
&=\frac{1}{a}\left[\frac{-2R}{2\left(\frac{\sigma_0}{a}+R^2\right)^2}
\ln\frac{\frac{\sigma_0}{a}+R^2-2Rr+r^2}{r^2}\Bigg|_R^\infty
-\frac{1}{\left(\frac{\sigma_0}{a}+R^2\right)r}\Bigg|_R^\infty\right.\notag\\
&\qquad\left.+\left(\frac{4R^2}{2\left(\frac{\sigma_0}{a}+R^2\right)^2}-\frac{1}{\frac{\sigma_0}{a}+R^2}\right)
\frac{2}{2\sqrt{\frac{\sigma_0}{a}}}\tan^{-1}\frac{2(r-R)}{2\sqrt{\frac{\sigma_0}{a}}}\Bigg|_R^\infty\right]\notag\\
&=\frac{1}{a}\left[\frac{R}{\left(\frac{\sigma_0}{a}+R^2\right)^2}\ln\frac{\sigma_0}{aR^2}
+\frac{1}{\left(\frac{\sigma_0}{a}+R^2\right)R}
+\frac{R^2-\frac{\sigma_0}{a}}{\left(\frac{\sigma_0}{a}+R^2\right)^2}\cdot\frac{\frac{\pi}{2}}{\sqrt{\frac{\sigma_0}{a}}}\right]\notag\\
&=\frac{1}{a\left(\frac{\sigma_0}{a}+R^2\right)^2}\left[R\ln\frac{\sigma_0}{aR^2}
+\frac{\frac{\sigma_0}{a}+R^2}{R}+\left(R^2-\frac{\sigma_0}{a}\right)\frac{\pi}{2\sqrt{\frac{\sigma_0}{a}}}\right]
\end{align}
Substituting this into the expression for $\varphi$ gives
\[
\varphi=-\frac{\sigma_0E_0R^2}{a\left(\frac{\sigma_0}{a}+R^2\right)^2}
\left[R\ln\frac{\sigma_0}{aR^2}+\frac{\frac{\sigma_0}{a}+R^2}{R}
+\left(R^2-\frac{\sigma_0}{a}\right)\frac{\pi}{2\sqrt{\frac{\sigma_0}{a}}}\right].
\]
Inserting the data, and noting that $\dfrac{\sigma_0}{a}\ll R^2$, we have
\[
\varphi\approx-3.83\times10^5\ (\text{V}).
\]

(4) Let the total charge on the ground be $q$. The decrease of $q$ is caused by the current in the ground, so
\[
-\frac{\dd q}{\dd t}=\sigma_0E_0\cdot4\pi R^2=\sigma_0\frac{q}{4\pi\varepsilon_0R^2}\cdot4\pi R^2=\frac{\sigma_0}{\varepsilon_0}q,
\]
that is,
\[
\frac{\dd q}{q}=-\frac{\sigma_0}{\varepsilon_0}\dd t.
\]
Integrating, we obtain
\[
q=q_0\ee^{-\frac{\sigma_0}{\varepsilon_0}t}=q_0\ee^{-t/\tau},
\]
\[
\tau=\frac{\varepsilon_0}{\sigma_0}\approx294\ (\text{s})\approx300\ (\text{s}).
\]
This is roughly the time required for the charge on the Earth to be neutralised by the current. The mechanism which keeps the charge on the Earth from being neutralised and roughly constant is mainly lightning and thunderbolts.
\end{solution}

\begin{solution}{3.35}
(1) Let the distance between the two wires be $d$, let the current in the wires be $I$, and let the linear charge density carried by them be $\lambda$. Then, by Gauss's law, the electric field produced by one wire at the position of the other wire is $E=\dfrac{\lambda}{2\pi\varepsilon_0d}$, and the electrostatic (attractive) force acting on unit length of the wire is
\[
F_e=\frac{\lambda^2}{2\pi\varepsilon_0d}.
\]
Similarly, by Amp\`ere's circuital law, the magnetic field produced by the current in one wire at the position of the other wire is $B=\dfrac{\mu_0I}{2\pi d}$, and the magnetic (repulsive) force acting on unit length of the wire is
\[
F_m=\frac{\mu_0I^2}{2\pi d}.
\]
For no force to act, we must have $F_e=F_m$, that is,
\[
\frac{\mu_0I^2}{2\pi d}=\frac{\lambda^2}{2\pi\varepsilon_0d}.
\]
Hence
\[
I=\sqrt{\frac{1}{\varepsilon_0\mu_0}}\,\lambda=c\lambda
\qquad\left[\because\ \sqrt{\frac{1}{\varepsilon_0\mu_0}}=c\ (\text{the speed of light})\right].
\]
But
\[
\lambda=C^{*}V,
\]
where $C^{*}$ is the capacitance per unit length between the wires and $V$ is the voltage between the two wires. Hence the required power is
\[
P=IV=cC^{*}V^2=3\times10^8\times10^{-11}\times\left(33\times10^3\right)^2
\]
\[
=3.27\times10^6\ (\text{W}).
\]

(2) The electric energy $w_e$ and the magnetic energy $w_m$ stored per unit length are
\[
w_e=\frac{1}{2}C^{*}V^2=\frac{1}{2}\times10^{-11}\times\left(33\times10^3\right)^2=5.45\times10^{-3}\ (\text{J/m}),
\]
\[
w_m=\frac{1}{2}L^{*}I^2,
\]
where the self-inductance per unit length $L^{*}$ can be obtained from $C^{*}$. Let the radius of the two wires be $a$ $(a\ll d)$; then the voltage between the two wires is
\[
V=\int_a^d E\dd r=2\int_a^d\frac{\lambda}{2\pi\varepsilon_0r}\dd r=\frac{\lambda}{\pi\varepsilon_0}\ln\frac{d}{a},
\]
\[
\therefore\ C^{*}=\frac{\lambda}{V}=\frac{\pi\varepsilon_0}{\ln\dfrac{d}{a}}.
\]
The magnetic flux per unit length is
\[
\Phi=\int_a^d B\dd r=2\int_a^d\frac{\mu_0I}{2\pi r}\dd r=\frac{\mu_0I}{\pi}\ln\frac{d}{a},
\]
\[
L^{*}=\frac{\Phi}{I}=\frac{\mu_0}{\pi}\ln\frac{d}{a}.
\]
It is thus seen that
\[
L^{*}C^{*}=\varepsilon_0\mu_0 .
\]
\begin{align}
\therefore\ w_m&=\frac{1}{2}L^{*}I^2=\frac{1}{2}\frac{\varepsilon_0\mu_0}{C^{*}}I^2
=\frac{1}{2}\frac{\varepsilon_0\mu_0}{C^{*}}\cdot c^2\lambda^2\notag\\
&=\frac{1}{2}\frac{1}{C^{*}}C^{*2}V^2=\frac{1}{2}C^{*}V^2\notag\\
&=w_e=5.45\times10^{-3}\ (\text{J/m}).\notag
\end{align}
\end{solution}

\begin{solution}{3.36}
(1) The magnetic field produced at the point $O$ by a solenoid of length $\dd z$ can be obtained by superposing the fields of circular coils. By the Biot--Savart law (as shown in the figure),
\[
\dd B=\frac{\mu_0}{4\pi}\frac{n\dd z\,I\cdot2\pi R}{r^2}\cdot\frac{R}{r}
=\frac{\mu_0nIR^2}{2r^3}\dd z .
\]
But $\dd z\sin\theta=r\dd\theta$ and $r\sin\theta=R$, so the above expression becomes
\[
\dd B=\frac{\mu_0nI}{2}\frac{r^2\sin^2\theta\cdot r\dd\theta}{r^2\cdot r\sin\theta}
=\frac{\mu_0nI}{2}\sin\theta\,\dd\theta .
\]
\[
\therefore\ B=\int\dd B=\frac{\mu_0nI}{2}\int_{\theta_0}^{\pi/2}\sin\theta\,\dd\theta
=\frac{\mu_0nI}{2}\frac{L}{\sqrt{R^2+L^2}},
\]
(where $\theta_0=\arcsin\dfrac{R}{\sqrt{R^2+L^2}}$).

\begin{figbox}{3.36}
\begin{tikzpicture}[scale=0.92,>=latex]
  % solenoid
  \draw[thickline] (-5.2,1)--(0,1);
  \draw[thickline] (-5.2,-1)--(0,-1);
  \draw[thickline] (-5.2,-1)--(-5.2,1);
  \draw[thickline] (0,-1)--(0,1);
  % winding marks
  \foreach \x in {-5.0,-4.6,...,-0.2}{%
    \draw[thinline] (\x,1)--(\x-0.14,0.80);
    \draw[thinline] (\x,-1)--(\x-0.14,-0.80);}
  % axis
  \draw[axis] (-5.9,0)--(2.1,0) node[right] {$z$};
  \fill (0,0) circle (1.2pt);
  \node[below left] at (0,0) {$O$};
  % rays bounding the slice dz
  \draw[dashedline] (0,0)--(-1.30,1);
  \draw[dashedline] (0,0)--(-0.75,1);
  % theta and dtheta
  \draw[thickline] (180:0.80) arc (180:142.4:0.80);
  \node at (162:0.99) {$\theta$};
  \draw[thickline] (142.4:1.05) arc (142.4:126.9:1.05);
  \node at (134:1.02) {$\dd\theta$};
  % slice dz
  \draw[dashedline] (-1.30,1)--(-1.30,1.34);
  \draw[dashedline] (-0.75,1)--(-0.75,1.34);
  \draw[dim] (-1.30,1.28)--(-0.75,1.28) node[midway,above] {$\dd z$};
  % radius R
  \draw[dim] (-5.55,0)--(-5.55,1) node[midway,left] {$R$};
  % Gaussian cylinder about O
  \draw[dashedline] (0,-0.30)--(0.90,-0.30)--(0.90,0.30)--(0,0.30)--cycle;
  \fill (0.90,0) circle (1.2pt);
  \node[above right] at (0.90,0) {$O'$};
  \draw[dim] (0,-0.60)--(0.90,-0.60) node[midway,below] {$\dd z$};
  \node[right] at (0.94,0.24) {$\rho$};
\end{tikzpicture}
\end{figbox}

(2) Near the point $O$ construct a cylindrical surface of height $\dd z$ and radius $\rho$ $(\rho\ll R)$, as shown in the figure. According to the problem, now $L\to\infty$, and the magnetic field at $O$ is $B_O=\dfrac{\mu_0nI}{2}$. At the point $O'$ on the axis at a distance $\dd z$ from $O$, from the above result,
\[
B_{O'}=\frac{\mu_0nI}{2}\int_0^{\theta'}\sin\theta\,\dd\theta
=\frac{\mu_0nI}{2}\left(1-\cos\theta'\right)
\]
\[
\approx\frac{\mu_0nI}{2}\left(1-\frac{\dd z}{R}\right).
\]
Taking this cylindrical surface as a Gaussian surface, from the Gauss law for the magnetic field, $\oint\vec{B}\cdot\dd\vec{s}=0$, we obtain
\[
B_\rho\cdot2\pi\rho\dd z+B_{O'}\pi\rho^2=B_O\pi\rho^2,
\]
that is,
\[
B_\rho2\pi\rho\dd z=\left(B_O-B_{O'}\right)\pi\rho^2=\frac{\mu_0nI}{2}\frac{\dd z}{R}\pi\rho^2,
\]
\[
\therefore\ B_\rho=\frac{\mu_0nI}{4}\frac{\rho}{R}.
\]
\end{solution}

\begin{solution}{3.37}
When a dielectric moves with speed $v$ in a magnetic field, the charges in the dielectric are subject to the Lorentz force. This action is equivalent to an electric field $\vec{E}^{*}=\vec{v}\times\vec{B}$. Under the action of this field the dielectric becomes polarised, and the polarisation charges produced by the surface polarisation also excite an electric field; this field in turn affects the polarisation of the dielectric, and finally the dielectric is polarised under the joint action of these two fields. Let the polarisation intensity of the dielectric be $P$; then the surface density of the polarisation charge is $\sigma'=P$, and the magnitude of the electric field produced by the polarisation charges is $E'=\dfrac{\sigma'}{\varepsilon_0}=\dfrac{P}{\varepsilon_0}$, but its direction is opposite to that of $\vec{P}$. Substituting this into the relation between the polarisation intensity and the electric field gives
\[
P=(\varepsilon_r-1)\varepsilon_0\left(vB-\frac{P}{\varepsilon_0}\right).
\]
Solving this, we obtain
\begin{equation}
P=\left(1-\frac{1}{\varepsilon_r}\right)\varepsilon_0vB. \tag{1}
\end{equation}
The potential difference between the two plates is caused by the electric field excited by the charges, that is, $U=E'd=\dfrac{P}{\varepsilon_0}d$; hence
\begin{equation}
U=\left(1-\frac{1}{\varepsilon_r}\right)vBd. \tag{2}
\end{equation}
In the situation shown in the figure, the upper plate is at the higher potential and the lower plate at the lower potential.

\begin{figbox}{3.37}
\begin{tikzpicture}[scale=1.0,>=latex]
  % plates
  \draw[thickline] (-3.7,0.78)--(3.7,0.78);
  \draw[thickline] (-3.7,-0.78)--(3.7,-0.78);
  % dielectric slab
  \draw[thickline] (-3.1,0.56) rectangle (3.1,-0.56);
  % fields inside
  \node at (-2.45,0.20) {$\otimes\,\vec{B}$};
  \draw[force] (-1.85,-0.12)--(-1.05,-0.12) node[midway,below] {$\vec{v}$};
  \draw[force] (0.10,0.42)--(0.10,-0.42) node[midway,right] {$\vec{E}'$};
  \draw[force] (0.95,-0.42)--(0.95,0.42) node[midway,right] {$\vec{P}$};
  \draw[force] (2.05,-0.42)--(2.05,0.42) node[midway,right] {$\vec{E}^{*}=\vec{v}\times\vec{B}$};
  % thickness d
  \draw[dim] (-3.45,0.78)--(-3.45,-0.78) node[midway,left] {$d$};
\end{tikzpicture}
\end{figbox}
\end{solution}

\begin{solution}{3.38}
According to the discussion of Problem 3.37, the dielectric is polarised under the joint action of the Lorentz force and the field excited by the polarisation charges. Let the electric field in the dielectric (that is, the field excited by the polarisation charges) be $E$, and let the equivalent electric field corresponding to the Lorentz force be $vB=\omega rB$; then the polarisation intensity (note that $\varepsilon_r=1+\chi$) is
\begin{equation}
P=(\varepsilon_r-1)\varepsilon_0\left(E+\omega rB\right)=\chi\varepsilon_0\left(E+\omega rB\right). \tag{1}
\end{equation}
Here both $P$ and $E$ are taken as positive along the radial direction outwards, so that they agree with the direction of the equivalent field of the Lorentz force (see the figure). By relating them to the polarisation charges, $E$ and $P$ can be connected. Take a cylindrical surface of radius $r$ as a Gaussian surface; by Gauss's law,
\begin{equation}
E\cdot2\pi r=\frac{q'}{\varepsilon_0}, \tag{2}
\end{equation}
where $q'$ is the polarisation charge inside the cylindrical surface. By the relation between the polarisation intensity and the polarisation charge, the polarisation charge inside the cylindrical surface is determined by
\begin{equation}
q'=-\oint\vec{P}\cdot\dd\vec{s}=-P\cdot2\pi r. \tag{3}
\end{equation}
Substituting into (2) gives $E=-\dfrac{P}{\varepsilon_0}$; substituting this into (1) gives
\[
P=\chi\varepsilon_0\left(-\frac{P}{\varepsilon_0}+\omega rB\right),
\]
from which
\begin{equation}
P(r)=\frac{\chi}{1+\chi}\varepsilon_0\omega rB. \tag{4}
\end{equation}
Substituting the values of $P$ at the radii $r$ and $r+\Delta r$ respectively into (3), we obtain the polarisation charge density
\begin{equation}
\rho'=\frac{q'(r+\Delta r)-q'(r)}{2\pi r\Delta r}
=-2\,\frac{\chi}{1+\chi}\varepsilon_0\omega B. \tag{5}
\end{equation}

\begin{figbox}{3.38}
\begin{tikzpicture}[scale=0.95,>=latex]
  % cylinder
  \draw[thickline] (-1.25,-1.55)--(-1.25,1.55);
  \draw[thickline] (1.25,-1.55)--(1.25,1.55);
  \draw[thickline] (0,1.55) ellipse (1.25 and 0.34);
  \draw[thickline] (-1.25,-1.55) arc (180:360:1.25 and 0.34);
  \draw[dashedline] (1.25,-1.55) arc (0:180:1.25 and 0.34);
  % axis
  \draw[dashedline] (0,-1.95)--(0,2.30);
  % rotation
  \draw[force] (0.05,2.02) arc (90:-15:0.42);
  \node[right] at (0.62,2.15) {$\omega$};
  % radius
  \draw[dim] (0,0.30)--(1.25,0.30) node[midway,above] {$R$};
  % magnetic field
  \draw[force] (1.80,-0.60)--(1.80,0.60) node[midway,right] {$\vec{B}$};
\end{tikzpicture}
\end{figbox}

On the surface of the cylinder, however, there is a surface distribution of polarisation charge
\begin{equation}
\sigma'=P(R)=\frac{\chi}{1+\chi}\varepsilon_0\omega RB. \tag{6}
\end{equation}
\end{solution}

\begin{solution}{3.39}
(1) Let the field strength at a distance $r$ from the axis be $E$; then $E=E_0/r$, and
\[
U=\int_a^b E\dd r=E_0\int_a^b\frac{\dd r}{r}=E_0\ln\frac{b}{a},
\]
\[
\therefore\ E=\frac{U}{r\ln\dfrac{b}{a}}.
\]
By Ohm's law,
\[
j=\frac{1}{\rho}E=\frac{U}{\rho r\ln\dfrac{b}{a}},
\]
\[
I=2\pi rLj=\frac{2\pi LU}{\rho\ln\dfrac{b}{a}}.
\]

(2) When a magnetic field is present, let the drift velocity of the charge carriers be $u$; then
\[
\vec{j}=\frac{1}{\rho}\left(\vec{E}+\vec{u}\times\vec{B}\right)
=\frac{1}{\rho}\vec{E}+\frac{1}{\rho ne}\vec{j}\times\vec{B}.
\]
Let $\vec{j}=j_r\hat{r}+j_\theta\hat{\theta}$, $\vec{E}=E\hat{r}$ and $\vec{B}=B\hat{z}$; then
\begin{align}
j_r\hat{r}+j_\theta\hat{\theta}
&=\frac{1}{\rho}E\hat{r}+\frac{1}{\rho ne}\left(j_r\hat{r}+j_\theta\hat{\theta}\right)\times B\hat{z}\notag\\
&=\left(\frac{1}{\rho}E+\frac{1}{\rho ne}j_\theta B\right)\hat{r}
-\frac{1}{\rho ne}j_rB\hat{\theta}.\notag
\end{align}
Hence
\[
j_r=\frac{1}{\rho}E+\frac{1}{\rho ne}j_\theta B,
\]
\[
j_\theta=-\frac{1}{\rho ne}j_rB.
\]
Substituting the expression for $j_\theta$ into the expression for $j_r$, we obtain
\[
j_r=\frac{1}{\rho}E-\left(\frac{B}{\rho ne}\right)^2j_r .
\]
Solving this, we obtain
\[
j_r=\frac{E}{\rho\left[1+\left(\dfrac{B}{\rho ne}\right)^2\right]}.
\]
\[
\therefore\ I'=2\pi rLj_r
=\frac{2\pi rL}{\rho\left[1+\left(\dfrac{B}{\rho ne}\right)^2\right]}\cdot\frac{U}{r\ln\dfrac{b}{a}}
\]
\[
=\frac{2\pi LU}{\rho\ln\dfrac{b}{a}}\cdot\frac{1}{1+\left(\dfrac{B}{\rho ne}\right)^2}
=\frac{I}{1+\left(\dfrac{B}{\rho ne}\right)^2}.
\]
\end{solution}

\begin{solution}{3.40}
(1) By the Biot--Savart law, the magnetic field produced by the current-carrying coil at its centre is
\begin{equation}
B_0=\frac{\mu_0}{4\pi}\frac{I\cdot2\pi b}{b^2}=\frac{\mu_0I}{2b}. \tag{1}
\end{equation}
Since the sphere is very small, the dielectric sphere will be uniformly magnetised under the action of this field. Let the magnetisation intensity of the dielectric sphere be $M$; then the magnetisation surface current density is
\[
i_M=M\sin\theta .
\]
The meaning of $\theta$ is shown in Fig.~3.40(a). By the Biot--Savart law, the magnetic field produced by the magnetisation currents at the centre of the sphere can be obtained by integration (Fig.~3.40(a)):
\[
B'=\frac{\mu_0}{4\pi}\int_0^\pi\frac{i_M\cdot a\dd\theta\cdot2\pi a\sin\theta}{a^2}\sin\theta
=\frac{\mu_0M}{2}\int_0^\pi\sin^3\theta\,\dd\theta
\]
\[
=\frac{2}{3}\mu_0M .
\]
The resultant field is
\[
B=B_0+B'=B_0+\frac{2}{3}\mu_0M .
\]
This is also the field throughout the interior of the sphere. From the equation of state,
\[
M=(\mu_r-1)H=(\mu_r-1)\frac{B}{\mu_r\mu_0},
\]
and substituting this into the above expression gives
\[
B=B_0+\frac{2}{3}\mu_0\cdot\frac{\mu_r-1}{\mu_r\mu_0}B
=B_0+\frac{2(\mu_r-1)}{3\mu_r}B .
\]
Solving, we obtain
\begin{equation}
B=\frac{3\mu_r}{\mu_r+2}B_0 . \tag{2}
\end{equation}
Thus
\begin{equation}
M=\frac{\mu_r-1}{\mu_r\mu_0}B
=\frac{3(\mu_r-1)}{\mu_0(\mu_r+2)}B_0
=\frac{3(\mu_r-1)}{2(\mu_r+2)}\frac{I}{b}, \tag{3}
\end{equation}
where $B_0=\dfrac{\mu_0I}{2b}$ from (1) has been substituted.

\begin{figbox}{3.40(a)}
\begin{tikzpicture}[scale=1.0,>=latex]
  % sphere
  \draw[thickline] (0,0) circle (1.45);
  % axis and applied field
  \draw[axis] (-1.85,0)--(2.05,0);
  \draw[force] (2.05,0)--(3.00,0) node[right] {$\vec{B}$};
  % magnetisation
  \draw[force] (-0.40,-0.50)--(0.50,-0.50) node[below] {$\vec{M}$};
  % radii bounding the element
  \draw[thickline] (0,0)--(57:1.45);
  \draw[thickline] (0,0)--(68:1.45);
  \draw[thickline] (57:1.25) arc (57:68:1.25);
  \node at (72:1.42) {$\dd\theta$};
  % angle theta
  \draw[thickline] (0.70,0) arc (0:57:0.70);
  \node at (28:0.95) {$\theta$};
  % radius label
  \node at (1.18,0.80) {$a$};
  \node[below left] at (0,0) {$O$};
\end{tikzpicture}
\end{figbox}

(2) The presence of the dielectric sphere increases the magnetic flux through the coil, and hence increases the self-inductance. This increase of the flux comes from the additional magnetic field produced by the magnetisation currents. Since the flux produced by this additional field has both positive and negative parts in the region $r<b$, the net increase of the flux equals the flux in the region $r>b$ (Fig.~3.40(b)). Outside the dielectric, the field of the magnetisation currents is equivalent to the field of a magnetic dipole. Hence the flux increment is
\[
\Delta\Phi=\int_b^\infty\frac{\mu_0}{4\pi}\frac{m}{r^3}\cdot2\pi r\dd r
=\frac{\mu_0}{4\pi}\cdot2\pi m\int_b^\infty\frac{\dd r}{r^2}
=\frac{\mu_0m}{2b},
\]
where $m$ is the magnetic moment of the magnetised sphere. Substituting (3) of part (1), we have

\begin{figbox}{3.40(b)}
\begin{tikzpicture}[scale=0.80,>=latex]
  % dipole field lines (above the axis)
  \draw[thickline] (0.60,0.00) .. controls (2.20,0.65) and (-2.20,0.65) .. (-0.60,0.00);
  \draw[thickline] (0.60,0.00) .. controls (2.55,1.15) and (-2.55,1.15) .. (-0.60,0.00);
  \draw[thickline] (0.60,0.00) .. controls (2.90,1.70) and (-2.90,1.70) .. (-0.60,0.00);
  % dipole field lines (below the axis)
  \draw[thickline] (0.60,0.00) .. controls (2.20,-0.65) and (-2.20,-0.65) .. (-0.60,0.00);
  \draw[thickline] (0.60,0.00) .. controls (2.55,-1.15) and (-2.55,-1.15) .. (-0.60,0.00);
  \draw[thickline] (0.60,0.00) .. controls (2.90,-1.70) and (-2.90,-1.70) .. (-0.60,0.00);
  % coil seen edge-on
  \draw[thickline] (-3.10,0)--(3.10,0);
  \fill (-3.10,0) circle (2.4pt);
  \fill (3.10,0) circle (2.4pt);
  % magnetised sphere
  \fill[white] (0,0) circle (0.52);
  \draw[thickline] (0,0) circle (0.52);
  \draw[dim] (0,-0.42)--(3.10,-0.42) node[midway,below] {$b$};
\end{tikzpicture}
\end{figbox}

\[
m=\frac{4}{3}\pi a^3M=\frac{4}{3}\pi a^3\cdot\frac{3(\mu_r-1)I}{2(\mu_r+2)b}
=\frac{2\pi a^3(\mu_r-1)I}{(\mu_r+2)b}.
\]
Substituting into the expression for $\Delta\Phi$,
\[
\Delta\Phi=\frac{\mu_0}{2b}\cdot\frac{2\pi a^3(\mu_r-1)I}{(\mu_r+2)b}
=\frac{\pi a^3\mu_0(\mu_r-1)}{b^2(\mu_r+2)}I .
\]
Finally we obtain the increment of the self-inductance
\[
\Delta L=\frac{\Delta\Phi}{I}=\frac{\pi a^3\mu_0(\mu_r-1)}{b^2(\mu_r+2)}.
\]
\end{solution}

\begin{solution}{3.41}
(1) The surface current density produced on the surface of the shell by the motion of the charges is
\[
i=\sigma\omega a=\sigma kta .
\]
The magnetic field inside the shell is uniform, and its magnetic flux density is
\[
\vec{B}=\mu_0i\,\hat{z}=\mu_0\sigma akt\,\hat{z},
\]
with the direction along the axis (as shown in the figure).

(2) Since the magnetic field inside the shell varies with time, by Faraday's law the electric field $E$ inside the shell satisfies
\[
E\cdot2\pi r=-\frac{\dd B}{\dd t}\pi r^2 .
\]
Hence the electric field is
\[
\vec{E}=-\frac{r}{2}\frac{\dd B}{\dd t}\hat{\theta}
=-\frac{r}{2}\mu_0\sigma ak\,\hat{\theta}.
\]

(3)
\[
W_m=\frac{1}{2\mu_0}B^2\cdot\pi a^2l=\frac{1}{2}\mu_0\sigma^2a^4k^2t^2\pi l,
\]
\[
W_E=\int\frac{1}{2}\varepsilon_0E^2\dd\tau
=\frac{1}{2}\varepsilon_0\int_0^a\frac{r^2}{4}\mu_0^2\sigma^2a^2k^2\cdot l\cdot2\pi r\dd r
\]
\[
=\frac{\pi\varepsilon_0}{16}\mu_0^2\sigma^2a^6k^2l .
\]

\begin{figbox}{3.41}
\begin{tikzpicture}[scale=0.85,>=latex]
  % cylindrical shell
  \draw[thickline] (-3.2,0.9)--(3.2,0.9);
  \draw[thickline] (-3.2,-0.9)--(3.2,-0.9);
  \draw[thickline] (3.2,0.9) arc (90:-90:0.30 and 0.9);
  \draw[thickline] (-3.2,0.9) arc (90:270:0.30 and 0.9);
  % windings
  \foreach \x in {-2.8,-2.4,...,2.8}{%
    \draw[thinline] (\x,0.9) arc (90:-90:0.26 and 0.9);}
  % axis
  \draw[axis] (-3.75,0)--(4.15,0);
  \node[below] at (4.05,0) {$z$};
  % magnetic field
  \draw[force] (3.45,-0.55)--(4.35,-0.55) node[right] {$\vec{B}$};
  % rotation
  \draw[force] (-2.55,1.28) arc (150:30:0.34);
  \node[above] at (-2.55,1.42) {$\omega$};
  % surface current
  \draw[force] (-1.75,0.92)--(-1.75,0.42);
  \node[right] at (-1.72,0.62) {$i$};
\end{tikzpicture}
\end{figbox}

(4) Inside the shell take a cylindrical surface of radius $r$ $(r<a)$; the Poynting vector on this surface is
\[
\vec{S}=\frac{1}{\mu_0}\vec{E}\times\vec{B}
=\frac{1}{\mu_0}\left(-\frac{r}{2}\mu_0\sigma ak\hat{\theta}\right)\times\left(\mu_0\sigma akt\hat{z}\right)
\]
\[
=\frac{\mu_0r}{2}\sigma^2a^2k^2t\,(-\hat{r}),
\]
its direction being towards the inside of the cylinder. The energy flowing per second into the region inside the cylindrical surface is
\[
S\cdot2\pi rl=\mu_0\pi r^2\sigma^2a^2k^2tl .
\]
The electric field energy inside the cylindrical surface is constant, whereas the magnetic field energy depends on time:
\[
W_m(r)=\frac{1}{2\mu_0}B^2\pi r^2l=\frac{\mu_0}{2}\pi r^2l\sigma^2a^2k^2t^2 .
\]
\[
\frac{\dd W_m(r)}{\dd t}=\mu_0\pi r^2l\sigma^2a^2k^2t=S\cdot\pi r^2l,
\]
that is, the rate of increase of the magnetic field energy is supplied by the energy flow represented by the Poynting vector.
\end{solution}

\begin{solution}{3.42}
In this case the long rod will be uniformly magnetised along the $x$ direction, because for a rod magnetised uniformly in this way the magnetisation currents on it produce a uniform magnetic field inside the rod; this field, superposed on the uniform external field, is still a uniform magnetic field inside the rod, which is consistent with the uniform magnetisation of the rod.

Let the intensity of the uniform magnetisation of the rod along the $x$ direction be $M$; then on the surface of the rod there will be a magnetisation surface current along the axis ($z$) direction, and the surface current density is
\[
i_M=M\sin\theta,
\]
as shown in Fig.~3.42(a)(1). This surface current may be imagined as formed by two opposite volume currents along the direction of the rod, displaced slightly from each other in the direction perpendicular to $M$, as shown in Fig.~3.42(a)(2). Let the volume current density be $j$ and the displacement be $l$ $(l\ll R)$; then
\[
jl=M.
\]
According to Amp\`ere's circuital law, the magnetic flux density at any point $P$ inside the rod is the superposition of the magnetic fields of the two mutually displaced opposite volume currents. Let the vectors from the axes $O_+$ $(O_-)$ of the volume currents along the positive (negative) $z$ direction to the point $P$ be $\vec{r}_+$ $(\vec{r}_-)$; then the field at $P$ can readily be found from Amp\`ere's circuital law as
\[
\vec{B}'_P=\frac{\mu_0}{2}\vec{j}\times\left(\vec{r}_+-\vec{r}_-\right)
=-\frac{\mu_0}{2}\vec{j}\times\vec{l}.
\]
Here $\vec{j}$ is along the positive $z$ direction and $\vec{l}$ points from $O_-$ to $O_+$; hence $\vec{j}\times(-\vec{l})=\vec{M}$, and therefore
\[
\therefore\ \vec{B}'_{\text{in}}=\frac{\mu_0}{2}\vec{M},
\]
which is a uniform field. The resultant magnetic field inside the rod is $\vec{B}_1=\vec{B}_0+\vec{B}'_{\text{in}}$, and the equation of state is written as
\[
M=(\mu_r-1)\frac{1}{\mu_r\mu_0}\left(B_0+B'_{\text{in}}\right)
=(\mu_r-1)\frac{1}{\mu_r\mu_0}\left(B_0+\frac{\mu_0}{2}M\right).
\]
From this we obtain
\[
M=\frac{2(\mu_r-1)}{\mu_0(\mu_r+1)}B_0,
\]
\[
B_1=B_0+\frac{\mu_0}{2}M=\frac{2\mu_r}{\mu_r+1}B_0.
\]

\begin{figbox}{3.42(a)}
\begin{tikzpicture}[scale=0.88,>=latex]
  % ---------- panel (1) ----------
  \begin{scope}
    \draw[thickline] (0,0) circle (1.40);
    % magnetisation
    \draw[force] (-0.85,0.72)--(0.15,0.72) node[midway,above] {$\vec{M}$};
    % radius and angle
    \draw[thickline] (0,0)--(32:1.40);
    \draw[thickline] (0.70,0) arc (0:32:0.70);
    \node at (16:0.92) {$\theta$};
    % x axis
    \draw[axis] (0,0)--(1.95,0) node[right] {$x$};
    % surface current symbols
    \node at (62:1.40) {$\odot$};
    \node at (90:1.40) {$\odot$};
    \node at (118:1.40) {$\odot$};
    \node at (242:1.40) {$\otimes$};
    \node at (270:1.40) {$\otimes$};
    \node at (298:1.40) {$\otimes$};
    \node[right] at (1.05,1.10) {$i_M$};
    \node at (0,-1.95) {(1)};
  \end{scope}
  % ---------- panel (2) ----------
  \begin{scope}[xshift=5.6cm]
    \draw[thickline] (0,0) circle (1.40);
    \draw[axis] (0,0)--(1.95,0) node[right] {$x$};
    \fill (0,0.17) circle (1.5pt);
    \node[above left] at (0.02,0.19) {$O_+$};
    \fill (0,-0.17) circle (1.5pt);
    \node[below left] at (0.02,-0.19) {$O_-$};
    \draw[dim] (-0.55,0.17)--(-0.55,-0.17) node[midway,left] {$l$};
    \coordinate (P) at (38:1.40);
    \draw[force] (0,0.17)--(P);
    \draw[force] (0,-0.17)--(P);
    \node[right] at (P) {$P$};
    \node at (0.78,0.52) {$r_+$};
    \node at (0.95,0.10) {$r_-$};
    \node at (0,-1.95) {(2)};
  \end{scope}
\end{tikzpicture}
\end{figbox}

Outside the rod, the magnetic field produced by the magnetisation surface current can still be regarded as the superposition of the fields of the two opposite volume currents. For the positive volume current, by Amp\`ere's circuital law,
\[
\vec{B}'_{P+}=\frac{\mu_0}{2}\frac{R^2}{r_+}\vec{j}\times\hat{r}_+ .
\]
Similarly, for the opposite volume current,
\[
\vec{B}'_{P-}=-\frac{\mu_0}{2}\frac{R^2}{r_-}\vec{j}\times\hat{r}_- .
\]

\begin{figbox}{3.42(b)}
\begin{tikzpicture}[scale=0.90,>=latex]
  \draw[thickline] (0,0) circle (1.25);
  \draw[axis] (-1.55,0)--(3.45,0);
  \node[below] at (3.40,0) {$x$};
  \fill (0,0.30) circle (1.5pt);
  \node[above left] at (0,0.32) {$O_+$};
  \fill (0,-0.30) circle (1.5pt);
  \node[below left] at (0,-0.32) {$O_-$};
  \node[below left] at (0,0) {$O$};
  \coordinate (P) at (2.90,0.75);
  \draw[force] (0,0.30)--(P);
  \draw[force] (0,-0.30)--(P);
  \node[above right] at (P) {$P$};
  \node at (1.30,0.72) {$\vec{r}_+$};
  \node at (1.55,0.06) {$\vec{r}_-$};
  % angles at P
  \draw[thickline] ($(P)+(180:0.72)$) arc (180:188.8:0.72);
  \node at ($(P)+(186:1.08)$) {$\varphi_+$};
  \draw[thickline] ($(P)+(180:0.95)$) arc (180:199.9:0.95);
  \node at ($(P)+(196:1.32)$) {$\varphi_-$};
  % fields at P
  \draw[force] (P)--($(P)+(98.8:1.00)$) node[above left] {$\vec{B}'_{P+}$};
  \draw[force] (P)--($(P)+(-70.1:0.70)$) node[below right] {$\vec{B}'_{P-}$};
\end{tikzpicture}
\end{figbox}

Here the meanings of $\vec{r}_+$ and $\vec{r}_-$ are as before, while $\hat{r}_+$ and $\hat{r}_-$ are the unit vectors of $\vec{r}_+$ and $\vec{r}_-$. Taking the midpoint of $O_+$ and $O_-$ as the origin $O$, we have $r_+\approx r-\dfrac{l}{2}\sin\theta$ and $r_-\approx r+\dfrac{l}{2}\sin\theta$ (Fig.~3.42(b)). Resolving $\vec{B}'_P=\vec{B}'_{P+}+\vec{B}'_{P-}$ into $r$ and $\theta$ components, it is readily seen from the figure that (neglecting second-order small quantities)
\begin{align}
B'_r&=\frac{\mu_0}{2}\frac{R^2}{r}j\left(\sin\varphi_++\sin\varphi_-\right)
=\frac{\mu_0}{2}\frac{R^2}{r}j\left[\frac{\frac{l}{2}\cos\theta}{r}+\frac{\frac{l}{2}\cos\theta}{r}\right]\notag\\
&=\frac{\mu_0}{2}\frac{R^2}{r^2}jl\cos\theta
=\frac{\mu_0}{2}\frac{R^2}{r^2}M\cos\theta,\notag
\end{align}
\begin{align}
B'_\theta&=\frac{\mu_0}{2}R^2j\left(\frac{\cos\varphi_+}{r_+}-\frac{\cos\varphi_-}{r_-}\right)
\approx\frac{\mu_0}{2}R^2j\left(\frac{1}{r_+}-\frac{1}{r_-}\right)\notag\\
&=\frac{\mu_0}{2}R^2j\left[\frac{1}{r-\frac{l}{2}\sin\theta}-\frac{1}{r+\frac{l}{2}\sin\theta}\right]
=\frac{\mu_0}{2}\frac{R^2}{r^2}j\cdot l\sin\theta\notag\\
&=\frac{\mu_0}{2}\frac{R^2}{r^2}M\sin\theta.\notag
\end{align}
Noting that $M=\dfrac{2(\mu_r-1)}{\mu_0(\mu_r+1)}B_0$, the resultant magnetic field outside the rod is
\[
B_{2r}=B_0\cos\theta+B'_r=B_0\cos\theta\left[\frac{\mu_r-1}{\mu_r+1}\frac{R^2}{r^2}+1\right],
\]
\[
B_{2\theta}=-B_0\sin\theta+B'_\theta=B_0\sin\theta\left[\frac{\mu_r-1}{\mu_r+1}\frac{R^2}{r^2}-1\right].
\]
\end{solution}

\begin{solution}{3.43}
Set up plane polar coordinates, as shown in the figure; the electron must move in the plane perpendicular to the axis. Let the potential of the inner cylinder be zero; then the potential at a distance $r$ from the axis is $U(r)$. Since the magnetic field does no work on the electron, the energy of the electron is conserved:
\begin{equation}
\frac{1}{2}m\left(v_r^2+v_\theta^2\right)=eU(r) \tag{1}
\end{equation}
Here $v_r=\dot{r}$ and $v_\theta=r\dot{\theta}$ are respectively the components of the velocity of the electron along the radial direction $(\hat{r})$ and along the transverse direction $(\hat{\theta})$. Because the electron is acted on by the Lorentz force of the magnetic field, the angular momentum of the electron about the axis is not conserved. The angular momentum of the electron about the axis is
\begin{align}
\vec{l} &= \vec{r}\times m\vec{v}=\vec{r}\times m\left(\vec{v}_r+\vec{v}_\theta\right)\notag\\
&= m\vec{r}\times\vec{v}_\theta=mr^2\dot{\theta}\,\hat{z} \tag{2}
\end{align}
because $\hat{r}\times\hat{\theta}=\hat{z}$. The torque exerted on the electron by the Lorentz force (the torque of the electric force is obviously zero) is
\begin{align}
\vec{\tau} &= \vec{r}\times(-e\vec{v}\times\vec{B})=-\vec{r}\times e\left[\left(\vec{v}_r+\vec{v}_\theta\right)\times\vec{B}\right]\notag\\
&= eBrv_r\hat{z}=eBr\dot{r}\,\hat{z} \tag{3}
\end{align}
because only the Lorentz force caused by the $r$ component of the velocity exerts a torque on the electron. From the angular momentum theorem $\vec{\tau}=\dd\vec{l}/\dd t$ we obtain
\[
eBr\dot{r}=\frac{\dd}{\dd t}\left(mr^2\dot{\theta}\right),
\]
that is,
\begin{equation}
eBr\frac{\dd r}{\dd t}=\frac{\dd}{\dd t}\left(mr^2\dot{\theta}\right) \tag{4}
\end{equation}
Integrating both sides,
\[
\int_a^r eBr\,\dd r=\int_0^{\dot{\theta}}\dd\left(mr^2\dot{\theta}\right),
\]
that is,
\[
\frac{1}{2}eB\left(r^2-a^2\right)=mr^2\dot{\theta}.
\]
In order that the electron cannot reach the anode, in the critical case we must have $\dot{r}=0$ when $r=b$. Substituting this into the above expression gives
\begin{equation}
\frac{1}{2}eB\left(b^2-a^2\right)=mb^2\dot{\theta}\,\big|_{r=b} \tag{5}
\end{equation}
The value of $\dot{\theta}|_{r=b}$ can be obtained from conservation of energy. Substituting $v_r=\dot{r}=0$ into (1),
\[
\frac{1}{2}mb^2\left(\dot{\theta}\,\big|_{r=b}\right)^2=eU,
\]
we obtain
\begin{equation}
\dot{\theta}\,\big|_{r=b}=\frac{1}{b}\sqrt{\frac{2eU}{m}} \tag{6}
\end{equation}
Substituting (6) into (5), we obtain the required magnetic flux density
\[
B=\frac{2b}{b^2-a^2}\sqrt{\frac{2mU}{e}}.
\]

\begin{figbox}{3.43}
\begin{tikzpicture}[scale=1.0,>=latex]
  % outer (anode) cylinder and inner (cathode) cylinder in cross-section
  \draw[thickline] (0,0) circle (2.3);
  \draw[thickline] (0,0) circle (0.8);
  \fill (0,0) circle (1.1pt);
  \node at (0.28,-0.26) {$O$};
  \coordinate (P) at (125:2.3);
  \draw[thickline] (0,0)--(P);
  \fill (P) circle (1.4pt);
  % angle phi measured from the leftward horizontal
  \draw[dashedline] (0,0)--(-1.75,0);
  \draw[thickline] (1.4,0) arc (180:125:1.4);
  \node at (-0.95,0.42) {$\varphi$};
  % labels of the two radii
  \node at (200:2.10) {$b$};
  \node at (212:0.52) {$a$};
  \node at ($(0,0)!0.55!(P)+(0.05,0.20)$) {$r$};
  % unit vectors at the electron
  \draw[force] (P)--++(125:0.72) node[above left] {$\hat{r}$};
  \draw[force] (P)--++(35:0.72) node[above right] {$\hat{\theta}$};
  % magnetic flux density into the page
  \node at (1.45,-0.85) {$\otimes\,\vec{B}$};
\end{tikzpicture}
\end{figbox}
\end{solution}

\begin{solution}{3.44}
(1) Let the speed of the particle be $v$. From
\[
\frac{mv^2}{R}=qvB_z
\]
we obtain
\[
\omega_c=\frac{v}{R}=\frac{qB_z}{m}.
\]

(2) At the position of the particle, since the magnetic field weakens gradually along the radial direction, the lines of magnetic induction must be curved in the shape of a barrel, as shown in the figure. Hence at $z\neq0$ the magnetic field will have a radial component. At $z>0$ the radial component is negative, while at $z<0$ it is positive. If the particle carries a positive charge, then at the position shown in the figure the particle must move out of the plane of the page. Hence when it deviates upwards, under the action of the negative radial component of the magnetic field it experiences a downward restoring force; when it deviates downwards, under the action of the positive radial component of the magnetic field it experiences an upward restoring force. Therefore the particle will vibrate in the vertical direction.

In order to find the radial component of the magnetic field, take the rectangular loop shown in the figure, whose sides are $\Delta r$ and $z$. From $\oint\vec{B}\cdot\dd\vec{l}=0$ we obtain
\[
\left[B_z(r+\Delta r)-B_z(r)\right]z-B_r\cdot\Delta r=0,
\]
that is,
\[
\frac{\partial B_z}{\partial r}\Delta r\cdot z-B_r\Delta r=0,
\]
\[
\therefore\ B_r=\frac{\partial B_z}{\partial r}\cdot z\qquad\left(\text{note that }\frac{\partial B_z}{\partial r}<0\right).
\]

\begin{figbox}{3.44}
\begin{tikzpicture}[scale=1.05,>=latex]
  % axes
  \draw[axis] (0,-1.55)--(0,1.85) node[right] {$z$};
  \draw[axis] (0,0)--(5.2,0) node[right] {$r$};
  \node[below left] at (0,0) {$O$};
  % barrel-shaped field lines
  \draw[thickline] (0.85,1.45) .. controls (2.2,1.25) and (3.4,1.15) .. (4.95,1.45);
  \draw[thickline] (0.85,-1.15) .. controls (2.2,-0.95) and (3.4,-0.85) .. (4.95,-1.15);
  \draw[thickline] (0.85,0.50) .. controls (2.2,0.32) and (3.4,0.28) .. (4.95,0.50);
  \draw[thickline] (0.85,-0.50) .. controls (2.2,-0.32) and (3.4,-0.28) .. (4.95,-0.50);
  \draw[force] (4.55,1.43)--++(0.32,-0.05);
  \draw[force] (4.55,-1.13)--++(0.32,0.05);
  % particle going out of the page, on the circular orbit of radius R
  \coordinate (P) at (3.2,0);
  \draw[thickline] (P) circle (0.15);
  \fill (P) circle (0.05);
  % rectangular loop, sides Delta r and z
  \draw[thickline] (2.7,0.66)--(3.7,0.66)--(3.7,-0.66)--(2.7,-0.66)--cycle;
  \draw[force] (3.02,0.66)--(3.30,0.66);
  \draw[force] (3.7,0.16)--(3.7,-0.12);
  \draw[force] (3.38,-0.66)--(3.10,-0.66);
  \draw[force] (2.7,-0.16)--(2.7,0.12);
  % dimensions
  \draw[dim] (2.7,1.62)--(3.7,1.62) node[midway,above] {$\Delta r$};
  \draw[dashedline] (2.7,0.66)--(2.7,1.62);
  \draw[dashedline] (3.7,0.66)--(3.7,1.62);
  \draw[dim] (4.10,-0.66)--(4.10,0.66) node[above] {$z$};
  \draw[dashedline] (3.7,-0.66)--(4.10,-0.66);
  \draw[dashedline] (3.7,0.66)--(4.10,0.66);
  \draw[dim] (0,-1.45)--(3.2,-1.45) node[midway,below] {$R$};
\end{tikzpicture}
\end{figbox}

Hence the equation of motion of the particle in the $z$ direction is
\[
m\ddot{z}=qvB_r=qv\frac{\partial B_z}{\partial r}z.
\]
Since $\dfrac{\partial B_z}{\partial r}<0$, the particle performs simple harmonic motion.

(3) From the equation of motion, the angular frequency of the simple harmonic motion of the particle is
\[
\omega=\sqrt{-\frac{q}{m}v\frac{\partial B_z}{\partial r}}.
\]
But $\omega_c=\dfrac{qB_z}{m}$ and $mv=qB_zR$, so that
\begin{align*}
\omega &= \sqrt{-\frac{q}{m}\frac{qB_zR}{m}\frac{\partial B_z}{\partial r}}=\sqrt{-\frac{q^2B_z^2}{m^2}\frac{R}{B_z}\frac{\partial B_z}{\partial r}}\\
&= \omega_c\sqrt{-\frac{R}{B_z}\frac{\partial B_z}{\partial r}},
\end{align*}
\[
\therefore\ \alpha=\sqrt{-\frac{R}{B_z}\frac{\partial B_z}{\partial r}}.
\]
\end{solution}

\begin{solution}{3.45}
With the coordinates shown in the figure of the problem, the force on the particle (we may take $v_r\approx0$) is
\[
\vec{F}=q\vec{v}\times\vec{B}=q\begin{vmatrix}
\hat{r} & \hat{\theta} & \hat{z}\\
0 & v_\theta & v_z\\
B_r & 0 & B_z
\end{vmatrix}
\]
Hence
\begin{equation}
F_z=-qv_\theta B_r. \tag{1}
\end{equation}
The quantity $B_r$ can be related to $\dfrac{\partial B_z}{\partial z}$. By Gauss's law for the magnetic field, take as the Gaussian surface the cylindrical surface of radius $r$ and height $\dd z$ shown in the figure; then
\[
B_z(z)\Delta s=B_z(z+\Delta z)\Delta s+B_r\cdot2\pi r\,\dd z.
\]
Then
\[
B_r\cdot2\pi r\,\dd z=-\frac{\partial B_z}{\partial z}\dd z\,\Delta s,
\]
\[
B_r=-\frac{r}{2}\frac{\partial B_z}{\partial z}.
\]

\begin{figbox}{3.45}
\begin{tikzpicture}[scale=0.95,>=latex]
  % cylinder of radius r and height dz
  \draw[thickline] (0,0) ellipse (1.5 and 0.38);
  \draw[thickline] (-1.5,0)--(-1.5,1.7);
  \draw[thickline] (1.5,0)--(1.5,1.7);
  \draw[thickline] (0,1.7) ellipse (1.5 and 0.38);
  \draw[dashedline] (-1.5,0) arc (180:360:1.5 and 0.38);
  \draw[dashedline] (-1.5,1.7) arc (180:360:1.5 and 0.38);
  % axis
  \draw[dashedline] (0,-0.85)--(0,2.3);
  \draw[axis] (0,2.3)--(0,2.95) node[right] {$z$};
  % field lines through the cylinder
  \draw[thickline,->] (-0.85,-0.55)--(-0.85,2.25);
  \draw[thickline,->] (0,-0.45)--(0,2.15);
  \draw[thickline,->] (0.85,-0.55)--(0.85,2.25);
  % heights of the two faces
  \draw[dashedline] (-1.5,0)--(-2.15,0);
  \node[left] at (-2.15,0) {$z$};
  \draw[dashedline] (-1.5,1.7)--(-2.15,1.7);
  \node[left] at (-2.15,1.7) {$z+\dd z$};
  % height dz of the cylinder
  \draw[dim] (1.85,0)--(1.85,1.7) node[midway,right] {$\dd z$};
  \draw[dashedline] (1.5,0)--(1.85,0);
  \draw[dashedline] (1.5,1.7)--(1.85,1.7);
  % radius of the cylinder
  \draw[thickline] (0,0)--(1.5,0) node[midway,above] {$r$};
  % area of the upper face
  \node[left] at (-3.05,2.62) {$\Delta S=\pi r^2$};
  \draw[->] (-2.45,2.50)--(-0.95,2.10);
\end{tikzpicture}
\end{figbox}

Substituting into (1), and noting that $v_\theta=-v_\perp$, we have
\[
F_z=-qv_\perp\frac{r}{2}\frac{\partial B_z}{\partial z}.
\]
But $r=\dfrac{mv_\perp}{qB_z}$; substituting this into the above expression gives
\begin{equation}
F_z=-qv_\perp\cdot\frac{1}{2}\frac{mv_\perp}{qB_z}\frac{\partial B_z}{\partial z}
=-\frac{1}{2}\frac{mv_\perp^2}{B_z}\frac{\dd B_z}{\dd z}. \tag{2}
\end{equation}

And from Newton's law,
\begin{equation}
F_z=m\frac{\dd v_z}{\dd t}=m\frac{\dd v_z}{\dd z}\frac{\dd z}{\dd t}=mv_z\frac{\dd v_z}{\dd z}. \tag{3}
\end{equation}
And from conservation of energy,
\begin{equation}
v_z^2+v_\perp^2=v_0^2. \tag{4}
\end{equation}
Substituting (4) and (3) into (2) gives
\[
mv_z\frac{\dd v_z}{\dd z}=-\frac{1}{2}m\left(v_0^2-v_z^2\right)\frac{1}{B_z}\frac{\dd B_z}{\dd z}.
\]
But $v_z\dfrac{\dd v_z}{\dd z}=\dfrac{1}{2}\dfrac{\dd v_z^2}{\dd z}$, so the above expression becomes
\[
\frac{\dd v_z^2}{v_0^2-v_z^2}=-\frac{\dd B_z}{B_z},
\]
that is,
\[
\frac{\dd\left(v_0^2-v_z^2\right)}{v_0^2-v_z^2}=\frac{\dd B_z}{B_z}.
\]
Integrating both sides, and supposing that the particle is reflected at the place where the magnetic field is $B$, at which $v_z=0$, we have
\[
\ln\left(v_0^2-v_z^2\right)\big|_0^B=\ln B_z\big|_{B_1}^{B},
\]
that is,
\[
\ln\frac{v_0^2}{v_0^2-v_z^2}=\ln\frac{B}{B_1}.
\]
Noting that $v_0^2-v_z^2=v_\perp^2$, the above expression is
\[
\frac{v_0^2}{v_\perp^2}=\frac{B}{B_1}.
\]
But $v_\perp/v_0=\sin\varphi_0$, so the above expression is
\[
\sin^2\varphi_0=\frac{B_1}{B}\geqslant\frac{B_1}{B_2},
\]
or
\[
\sin\varphi_0\geqslant\sqrt{\frac{B_1}{B_2}}=\sin\varphi_c.
\]
The required $\varphi_c$ is
\[
\varphi_c=\arcsin\sqrt{\frac{B_1}{B_2}}.
\]
\end{solution}

\begin{solution}{3.46}
Take the positive direction of circulation of the coil to be the same as the direction of the magnetic field.

(1) The equation of motion of the coil is
\begin{equation}
m\frac{\dd v}{\dd t}=IBb \tag{1}
\end{equation}
The circuit equation of the coil is
\begin{equation}
IR=-vBb \tag{2}
\end{equation}
Substituting (2) into (1) gives
\[
m\frac{\dd v}{\dd t}=-\frac{B^2b^2}{R}v,
\]
that is,
\[
\frac{\dd v}{v}=-\frac{B^2b^2}{mR}\dd t.
\]
Integrating both sides, we have
\[
\int_{v_0}^{v}\frac{\dd v}{v}=\int_0^t-\frac{B^2b^2}{mR}\,\dd t
\quad\Rightarrow\quad \ln\frac{v}{v_0}=-\frac{B^2b^2}{mR}t,
\]
that is,
\begin{equation}
v=v_0\,\ee^{-\frac{B^2b^2}{mR}t}. \tag{3}
\end{equation}
Integrating once more with respect to time, we have
\begin{align}
x &= \int_0^t v\,\dd t=\int_0^t v_0\,\ee^{-\frac{B^2b^2}{mR}t}\,\dd t\notag\\
&= \frac{mR}{B^2b^2}v_0\left(1-\ee^{-\frac{B^2b^2}{mR}t}\right). \tag{4}
\end{align}
After the coil enters the magnetic field its speed gradually decreases, and finally it stops.

(2) The equation of motion of the coil is the same as (1), while the circuit equation becomes
\begin{equation}
-L\frac{\dd I}{\dd t}-vBb=0. \tag{5}
\end{equation}
Differentiating (1), we obtain
\[
m\frac{\dd^2v}{\dd t^2}=Bb\frac{\dd I}{\dd t}.
\]
Substituting (5) into this, the above expression becomes
\[
m\frac{\dd^2v}{\dd t^2}=Bb\left(-\frac{Bbv}{L}\right)=-\frac{B^2b^2}{L}v.
\]
It is seen that $v$ is a simple harmonic function; let
\begin{equation}
v=V\cos(\omega t+\varphi), \tag{6}
\end{equation}
where
\[
\omega=\sqrt{\frac{B^2b^2}{mL}}=\frac{Bb}{\sqrt{mL}}.
\]
Integrating (6),
\[
x=\int_0^t v\,\dd t=\frac{V}{\omega}\left[\sin(\omega t+\varphi)-\sin\varphi\right].
\]
Here $V$ and $\varphi$ can be determined from the initial conditions:
\[
x_0=\frac{V}{\omega}(\sin\varphi-\sin\varphi)=0,\qquad v_0=V\cos\varphi.
\]
We may take $\varphi=0$; then $V=v_0$, and
\[
\therefore\ x=\frac{v_0}{\omega}\sin\omega t
\qquad\text{when } t<\frac{\pi}{\omega}=\pi\frac{\sqrt{mL}}{Bb}=\frac{T}{2},
\]
\[
x=-v_0\left(t-\frac{T}{2}\right)
\qquad\text{when } t>\frac{T}{2}.
\]
The coil behaves as if it were repelled back by the magnetic field.
\end{solution}

\begin{solution}{3.47}
(1) Let the time taken by the ring to fall down be $T$, and let the moment of inertia of the ring about the pivot be $I$. Then the kinetic energy of the ring when it has fallen down is estimated as
\begin{align*}
W_k &= \frac{1}{2}I\omega^2=\frac{1}{2}I\left(\frac{\pi/2}{T}\right)^2
=\frac{1}{2}\left(\frac{1}{2}mr^2+mr^2\right)\cdot\frac{\pi^2}{4T^2}\\
&= \frac{3\pi^2}{16}\frac{mr^2}{T^2}.
\end{align*}
The Joule heat is estimated as
\begin{align*}
W_h &= i^2RT=\left(\frac{\varepsilon}{R}\right)^2RT=\frac{\varepsilon^2}{R}T=\frac{(\Phi/T)^2}{R}T\\
&= \frac{\Phi^2}{RT}=\frac{(B\cdot\pi r^2)^2}{RT}=\frac{B^2\pi^2r^4}{RT}.
\end{align*}
The ratio of the two is
\[
\frac{W_k}{W_h}=\frac{3\pi^2mr^2}{16T^2}\frac{RT}{B^2\pi^2r^4}=\frac{3mR}{16B^2Tr^2}.
\]
With
\[
m=\rho\cdot2\pi r\cdot\pi r_1^2,\qquad
R=\frac{1}{\sigma}\frac{2\pi r}{\pi r_1^2}=\frac{1}{\sigma}\frac{2r}{r_1^2}
\]
substituted into the above expression,
\[
\frac{W_k}{W_h}=\frac{3}{16}\frac{\rho\cdot2\pi^2rr_1^2}{B^2Tr^2}\cdot\frac{1}{\sigma}\frac{2r}{r_1^2}
=\frac{3\pi^2}{4}\frac{\rho}{\sigma B^2T}.
\]
But $T\geqslant T_0$, where $T_0$ is the time taken by the ring to fall down in the absence of a magnetic field. It is not difficult to see that
\[
T_0\approx\frac{\pi}{2}\sqrt{\frac{I}{mgr}}=\frac{\pi}{2}\sqrt{\frac{\frac{3}{2}mr^2}{mgr}}=\frac{\pi}{2}\sqrt{\frac{3r}{2g}}.
\]
\[
\therefore\ \frac{W_k}{W_h}=\frac{3\pi^2}{4}\frac{\rho}{\sigma B^2T}
<\frac{3\pi^2}{4}\frac{\rho}{\sigma B^2T_0}
=\frac{3\pi^2}{4}\frac{\rho}{\sigma B^2}\frac{2}{\pi}\sqrt{\frac{2g}{3r}}
=\frac{3\pi}{2}\frac{\rho}{\sigma B^2}\sqrt{\frac{2g}{3r}}.
\]
Substituting the data, we obtain
\[
\frac{W_k}{W_h}<\frac{3\pi}{2}\frac{19.3\times10^3}{4.25\times10^7\times1^2}\sqrt{\frac{2\times9.8}{3\times0.01}}=5.5\times10^{-2}.
\]
It is seen that the released potential energy is mainly converted into Joule heat.

(2) Neglecting the kinetic energy, the Ampere torque is at every instant balanced with the gravitational torque during the falling of the ring. During the falling of the ring, when the plane of the ring makes an angle $\theta$ with the vertical plane, the magnetic flux through the ring is
\[
\Phi=BS\sin\theta,
\]
where $S=\pi r^2$. Then the electromotive force is
\[
\varepsilon=\left|\frac{\dd\Phi}{\dd t}\right|=BS\cos\theta\cdot\dot{\theta}.
\]
The current is
\[
i=\frac{\varepsilon}{R}=\frac{BS\cos\theta\cdot\dot{\theta}}{R}.
\]
The Ampere torque is
\[
\tau_B=iSB\cos\theta=\frac{B^2S^2\cos^2\theta}{R}\dot{\theta}.
\]
The gravitational torque is
\[
\tau_G=mgr\sin\theta.
\]
Setting $\tau_B=\tau_G$,
\[
\frac{B^2S^2\cos^2\theta}{R}\frac{\dd\theta}{\dd t}=mgr\sin\theta,
\]
so
\[
\frac{\dd\theta}{\dd t}=\frac{mgr\sin\theta\,R}{B^2S^2\cos^2\theta},
\]
\[
\dd t=\frac{B^2S^2}{mgrR}\frac{\cos^2\theta}{\sin\theta}\dd\theta.
\]
The time taken to fall down is
\[
t=\frac{B^2S^2}{mgrR}\int_{0.1}^{\frac{\pi}{2}}\frac{\cos^2\theta}{\sin\theta}\dd\theta,
\]
where
\begin{align*}
\int_{0.1}^{\frac{\pi}{2}}\frac{\cos^2\theta}{\sin\theta}\dd\theta
&= \int_{0.1}^{\frac{\pi}{2}}\frac{\dd\theta}{\sin\theta}-\int_{0.1}^{\frac{\pi}{2}}\sin\theta\,\dd\theta
=\left[\ln\tan\frac{\theta}{2}+\cos\theta\right]_{0.1}^{\frac{\pi}{2}}\\
&\approx 2.99-0.99=2.
\end{align*}
And
\begin{align*}
\frac{B^2S^2}{mgrR}
&= \frac{B^2\left(\pi r^2\right)^2}{\rho\cdot2\pi r\cdot\pi r_1^2gr\cdot\frac{1}{\sigma}\cdot\dfrac{2\pi r}{\pi r_1^2}}
=\frac{\sigma B^2r}{4\rho g}\\
&= \frac{4.25\times10^7\times1^2\times0.01}{4\times19.3\times10^3\times9.8}=0.56\ (\mathrm{s}).
\end{align*}
Substituting into the expression for $t$, we obtain
\[
t=0.56\times2=1.12\ (\mathrm{s}).
\]
\end{solution}

\begin{solution}{3.48}
(1) The magnetic field of the small bar magnet may be regarded as the field of a magnetic dipole. Take the centre of the bar as the origin and take the vertical axis of the bar magnet as the $z$-axis, as shown in the figure of the problem. On the axis the magnetic flux density may be written as
\begin{equation}
B_z=\frac{B_0}{z^3}. \tag{1}
\end{equation}
For, by the statement of the problem, when $z=1$, $B=B_0$. Since the coil is superconducting and its resistance is $R=0$, from the law of electromagnetic induction and Ohm's law, for the circuit of the coil we have
\[
-\frac{\dd\Phi}{\dd t}=IR=0.
\]
Integrating gives
\begin{equation}
\Phi=\text{constant}=0, \tag{2}
\end{equation}
because when the coil was originally far away, $\Phi=0$. When the coil is in equilibrium at $z=z_0$, let the current in the coil be $I_0$; then
\begin{equation}
\frac{B_0}{z_0^3}\pi a^2-I_0L=0. \tag{3}
\end{equation}
From this we obtain
\begin{equation}
I_0=\frac{B_0\pi a^2}{Lz_0^3}. \tag{4}
\end{equation}
The direction of $I_0$ is as shown in the figure.

In equilibrium the weight of the coil is balanced by the Ampere force, and the vertical component of the Ampere force acting on the coil is determined by the component $B_r$ of the magnetic field at the position of the coil along the radial direction of the coil. This radial component can be found from the rate of change of the $z$ component by means of Gauss's law for the magnetic field. At $z$ construct a cylindrical surface of height $\Delta z$ and radius $a$ as the Gaussian surface, as shown in the figure; then by Gauss's law,
\[
B_z(z)\pi a^2=B_z(z+\Delta z)\pi a^2+B_r\cdot2\pi a\cdot\Delta z,
\]
that is,
\[
-\frac{\partial B_z}{\partial z}\Delta z\cdot\pi a^2=B_r\cdot2\pi a\cdot\Delta z.
\]
Substituting (1) into the above expression, we find
\begin{equation}
B_r=\frac{3}{2}\frac{B_0a}{z^4}. \tag{5}
\end{equation}

\begin{figbox}{3.48}
\begin{tikzpicture}[scale=0.95,>=latex]
  % axis
  \draw[dashedline] (0,-3.5)--(0,2.55);
  \draw[axis] (0,2.55)--(0,3.2) node[right] {$z$};
  % bar magnet
  \fill[black!10] (-0.45,-3.00) rectangle (0.45,-2.25);
  \draw[thickline] (-0.45,-3.00) rectangle (0.45,-1.50);
  \draw[thickline] (-0.45,-2.25)--(0.45,-2.25);
  \node at (0,-1.95) {$\mathrm{N}$};
  \node at (0,-2.62) {$\mathrm{S}$};
  % coil drawn as a short solenoid, faces at z and z+dz
  \draw[thickline] (0,0.40) ellipse (1.00 and 0.27);
  \draw[thickline,postaction={decorate},
        decoration={markings,mark=at position 0.25 with {\arrow{>}}}]
        (0,1.20) ellipse (1.00 and 0.27);
  \draw[thickline] (-1.00,0.40)--(-1.00,1.20);
  \draw[thickline] (1.00,0.40)--(1.00,1.20);
  \draw[dashedline] (-1.00,0.40) arc (180:360:1.00 and 0.27);
  % labels of the two faces
  \draw[dashedline] (-1.00,0.40)--(-1.55,0.40);
  \node[left] at (-1.55,0.40) {$z$};
  \draw[dashedline] (-1.00,1.20)--(-1.55,1.20);
  \node[left] at (-1.55,1.20) {$z+\Delta z$};
  % current and height of the coil
  \node[right] at (1.05,1.52) {$I_0$};
  \draw[dim] (1.55,0.40)--(1.55,1.20) node[midway,right] {$\Delta z$};
  \draw[dashedline] (1.00,0.40)--(1.55,0.40);
  \draw[dashedline] (1.00,1.20)--(1.55,1.20);
\end{tikzpicture}
\end{figbox}

Substituting (4) and (5) into the force equilibrium equation of the coil,
\begin{equation}
I_0B_r(z_0)2\pi a=mg, \tag{6}
\end{equation}
we obtain
\[
\frac{B_0\pi a^2}{Lz_0^3}\cdot\frac{3}{2}\frac{B_0a}{z_0^4}\cdot2\pi a=mg.
\]
From this we find
\begin{equation}
z_0=\sqrt[7]{\frac{3B_0^2\pi^2a^4}{Lmg}}. \tag{7}
\end{equation}

(2) When the coil is moved upwards from the equilibrium position by a small amount $z'$, let the current in the coil become $I_0+i$. From $\Phi=0$, equation (2), we obtain
\[
\frac{B_0}{(z_0+z')^3}\pi a^2-L(I_0+i)=0.
\]
Retaining first-order small quantities in $z'/z_0$, the above expression becomes
\[
\frac{B_0\pi a^2}{z_0^3}\left(1-3\frac{z'}{z_0}\right)-LI_0-Li=0.
\]
Using (3) and (4), we obtain
\begin{equation}
i=-3\frac{z'}{z_0}\frac{B_0\pi a^2}{Lz_0^3}=-3I_0\frac{z'}{z_0}. \tag{8}
\end{equation}
The force on the coil at this time is
\[
F=(I_0+i)B_r(z_0+z')2\pi a-mg.
\]
From (5),
\[
B_r(z_0+z')=\frac{3}{2}\frac{B_0a}{(z_0+z')^4}=\frac{3}{2}B_r(z_0)\left(1-4\frac{z'}{z_0}\right).
\]
Substituting into the above expression,
\[
F=I_0\left(1+\frac{i}{I_0}\right)B_r(z_0)\left(1-4\frac{z'}{z_0}\right)2\pi a-mg.
\]
Substituting (8) into the above expression, neglecting small quantities of second and higher order, and using (6), we obtain
\[
F=-7mg\frac{z'}{z_0}.
\]
Then the equation of motion of the coil is
\[
m\ddot{z}'=F=-7mg\frac{z'}{z_0}.
\]
It is seen that the coil performs simple harmonic motion, with angular frequency
\[
\omega=\sqrt{\frac{7g}{z_0}},
\]
and period
\[
T=2\pi\sqrt{\frac{z_0}{7g}}.
\]
\end{solution}

\begin{solution}{3.49}
When the magnet falls, induced currents (eddy currents) are produced in the metal tube by electromagnetic induction. According to Lenz's law, the polarity of the magnetic field produced by the induced current is opposite to that of the magnet, which hinders the fall of the magnet, reduces the acceleration of the falling magnet, and finally makes the magnet fall uniformly with a certain speed. The constant speed of fall of the magnet can be found from energy relations.

To simplify the calculation, the magnet is regarded as a magnetic dipole of magnetic moment $\mu$, and the metal tube is imagined to consist of a series of identical circular coils arranged coaxially (with $\mu$ on the axis). Let the radius of each coil be $a$, its resistance be $R$, and the density of arrangement of the coils along the axial direction (that is, the number of turns per unit length) be $n$. Suppose that the direction of the magnetic moment $\mu$ of the magnetic dipole is vertically downwards. At a certain instant when $\mu$ falls with constant speed $v$, consider the induced current in a certain coil. Take the axis of the metal tube as the $z$-axis, with the downward direction positive and the origin at the position of $\mu$; the centre of the coil is at a distance $z$ below $\mu$, as shown in the figure. Let the distance of the wire of the coil from the centre of $\mu$ (that is, from the origin) be $r$, and let the angle between the vector $r$ directed from the origin to the wire and the $z$-axis be $\theta$. Then the induced electromotive force $\varepsilon$ in the coil is equal to the ratio of the magnetic flux $\dd\Phi$ of the field of the magnetic moment $\mu$ through the annular strip of circumference $2\pi a$ and width $\dd z=v\dd t$ to $\dd t$, that is,
\[
|\varepsilon|=\frac{\dd\Phi}{\dd t}=\left(B_r\sin\theta+B_\theta\cos\theta\right)2\pi a\,\dd z/\dd t.
\]
For the field of a magnetic dipole,
\[
B_r=\frac{\mu_0}{4\pi}\frac{2\mu\cos\theta}{r^3},
\]
\[
B_\theta=\frac{\mu_0}{4\pi}\frac{\mu\sin\theta}{r^3}.
\]
And
\[
\dd z/\dd t=v.
\]
Substituting into the above expression, we obtain
\begin{align*}
|\varepsilon| &= \frac{\mu_0}{4\pi}\left(\frac{2\mu\cos\theta}{r^3}\sin\theta+\frac{\mu\sin\theta}{r^3}\cos\theta\right)\cdot2\pi av\\
&= \frac{3}{2}\mu_0\mu av\frac{\sin\theta\cos\theta}{r^3}.
\end{align*}
Then the Joule heating power dissipated in the coil is
\[
p=I^2R=\varepsilon^2/R=\frac{9\mu_0^2\mu^2a^2v^2}{4R}\frac{\sin^2\theta\cos^2\theta}{r^6}.
\]

\begin{figbox}{3.49}
\begin{tikzpicture}[scale=0.95,>=latex]
  % axis
  \draw[axis] (0,0.25)--(0,-4.05) node[right] {$z$};
  \fill (0,0) circle (1.1pt) node[above left] at (-0.02,0.06) {$O$};
  % magnetic moment
  \draw[force] (0,0)--(0,-0.95) node[midway,right] {$\mu$};
  % coil seen in perspective
  \draw[dashedline] (1.9,-2.7) arc (0:180:1.9 and 0.42);
  \draw[thickline] (1.9,-2.7) arc (0:-180:1.9 and 0.42);
  % radius a of the coil
  \draw[thickline] (0,-2.7)--(1.9,-2.7) node[midway,above] {$a$};
  % r from origin to the wire
  \draw[dashedline] (0,0)--(-1.9,-2.7) node[pos=0.42,above left=-1pt] {$r$};
  \draw[thickline] (0,-1.2) arc (-90:-125:1.2);
  \node at (-0.52,-1.42) {$\theta$};
  % height z
  \draw[dashedline] (0,0)--(-3.5,0);
  \draw[dashedline] (-1.9,-2.7)--(-3.5,-2.7);
  \draw[dim] (-3.5,0)--(-3.5,-2.7) node[midway,left] {$z$};
  % width dz = v dt
  \draw[dim] (2.60,-2.94)--(2.60,-2.46) node[above right] {$\dd z=v\,\dd t$};
  \draw[dashedline] (1.9,-2.7)--(2.60,-2.7);
  % field components at the wire
  \draw[force] (-1.9,-2.7)--++(-0.50,-0.12);
  \node[below left] at (-2.42,-2.82) {$B_r$};
  \draw[force] (-1.9,-2.7)--++(-0.12,-0.50);
  \node[below left] at (-2.02,-3.22) {$B_\theta$};
\end{tikzpicture}
\end{figbox}

In the interval of height $\dd z$ there are altogether $\dd N=n\dd z=n\dfrac{r\,\dd\theta}{\sin\theta}$ turns of coils, and the Joule heating power dissipated in these coils is
\begin{align*}
\dd P &= p\,\dd N=pn\,\dd z=\frac{9\mu_0^2\mu^2a^2v^2}{4R}\frac{\sin^2\theta\cos^2\theta}{r^6}\cdot n\frac{r\,\dd\theta}{\sin\theta}\\
&= \frac{9\mu_0^2\mu^2a^2v^2}{4R}\frac{\sin\theta\cos^2\theta}{r^5}n\,\dd\theta\\
&= \frac{9\mu_0^2\mu^2v^2n}{4Ra^3}\sin^6\theta\cos^2\theta\,\dd\theta.
\end{align*}
Here the relation $r=a/\sin\theta$ has been used. The Joule heating power dissipated in all the coils is the integral of the above expression. For simplicity we may suppose that the coils are arranged to an infinite length, so that the limits of integration are taken as $0$ and $\pi$, that is,
\[
P=\int\dd P=\int_0^{\pi}\frac{9\mu_0^2\mu^2v^2n}{4Ra^3}\sin^6\theta\cos^2\theta\,\dd\theta.
\]
But
\[
\int_0^{\pi}\sin^6\theta\cos^2\theta\,\dd\theta
=\int_0^{\pi}\sin^6\theta\left(1-\sin^2\theta\right)\dd\theta
=\int_0^{\pi}\left(\sin^6\theta-\sin^8\theta\right)\dd\theta,
\]
and from the integration formula given in the problem, we have
\begin{align*}
\int_0^{\pi}\sin^6\theta\,\dd\theta
&= -\frac{\sin^5\theta\cos\theta}{6}\bigg|_0^{\pi}+\frac{5}{6}\int_0^{\pi}\sin^4\theta\,\dd\theta\\
&= \frac{5}{6}\left[-\frac{\sin^3\theta\cos\theta}{4}\bigg|_0^{\pi}+\frac{3}{4}\int_0^{\pi}\sin^2\theta\,\dd\theta\right]\\
&= \frac{3}{5}\left[\frac{1}{2}\theta-\frac{1}{4}\sin2\theta\right]_0^{\pi}=\frac{3\pi}{10}.
\end{align*}
\begin{align*}
\int_0^{\pi}\sin^8\theta\,\dd\theta
&= -\frac{\sin^7\theta\cos\theta}{8}\bigg|_0^{\pi}+\frac{7}{8}\int_0^{\pi}\sin^6\theta\,\dd\theta\\
&= \frac{7}{8}\cdot\frac{3\pi}{10}.
\end{align*}
\[
\therefore\ \int_0^{\pi}\sin^6\theta\cos^2\theta\,\dd\theta
=\int_0^{\pi}\left(\sin^6\theta-\sin^8\theta\right)\dd\theta
=\left(1-\frac{7}{8}\right)\frac{3\pi}{10}=\frac{3\pi}{80}.
\]
Substituting into the expression for $P$, we obtain the total thermal power dissipated in all the coils:
\[
P=\frac{9\mu_0^2\mu^2v^2n}{4Ra^3}\cdot\frac{3\pi}{80}.
\]
From conservation of energy, when the magnet (magnetic dipole) falls steadily, the power of the work done on it by gravity must be equal to the total thermal power dissipated in all the coils, that is,
\[
P=\frac{9\mu_0^2\mu^2v^2n}{4Ra^3}\cdot\frac{3\pi}{80}=mgv.
\]
From this we obtain the speed of fall of the magnetic dipole:
\[
v=\frac{320Ra^3mg}{27\pi\mu_0^2\mu^2n}.
\]
\end{solution}

\begin{solution}{3.50}
Let the current through the appliance be $I_L$; then $I_L$ lags behind $U$ in phase by $\varphi$. Let the capacitance connected in parallel be $C$; then the current through the capacitor is
\[
I_C=\frac{U}{1/(C\omega)}=UC\omega,
\]
and its phase leads $U$ by $\pi/2$. The vector relation among $I_L$, $I_C$ and $U$ is as shown in the figure. In order to raise the power factor to $1$, the current $I=I_L+I_C$ must be in phase with $U$. Hence
\[
UC\omega=I_L\sin\varphi.
\]
From the conditions given,
\[
P=I_LU\cos\varphi.
\]
Dividing the two expressions gives
\[
C=\frac{P\sin\varphi}{\omega U^2\cos\varphi}
=\frac{P\sqrt{1-\cos^2\varphi}}{\omega U^2\cos\varphi}.
\]
Substituting the data,
\[
C=\frac{100\sqrt{1-0.6^2}}{100\pi\times220^2\times0.6}
=8.77\times10^{-6}\unitv{F}=8.77\,\mu\mathrm{F}.
\]

\begin{figbox}{3.50}
\begin{tikzpicture}[scale=0.95,>=latex]
  \coordinate (O) at (0,0);
  \draw[force] (O)--(4.40,0) node[right] {$U$};
  \draw[force] (O)--(0,2.40) node[above] {$I_C=UC\omega$};
  \draw[force] (O)--(3.20,-2.40) node[below right] {$I_L$};
  \draw[dashedline,->] (0,2.40)--(3.20,2.40) node[midway,above] {$I$};
  \draw[dashedline] (3.20,2.40)--(3.20,-2.40);
  \draw[thickline] (1.00,0) arc (0:-36.87:1.00);
  \node at (1.22,-0.38) {$\varphi$};
\end{tikzpicture}
\end{figbox}
\end{solution}

\begin{solution}{3.51}
Let the currents in the four loops be as shown in the figure. Since $\omega=1/\sqrt{LC}$ and $R=\sqrt{L/C}$, the complex impedances of the resistor, the inductor and the capacitor are respectively
\[
\tilde{Z}_R=R,
\]
\[
\tilde{Z}_L=j\omega L=j\sqrt{L/C}=jR,
\]
\[
\tilde{Z}_C=\frac{1}{j\omega C}=-\sqrt{L/C}=-jR.
\]
\begin{figbox}{3.51}
\begin{tikzpicture}[scale=0.95,>=latex]
  \coordinate (A) at (0,2.35);
  \coordinate (B) at (-2.25,-1.30);
  \coordinate (C) at (2.25,-1.30);
  \coordinate (D) at (0,-0.10);
  \draw[thickline] (A)--(B)--(C)--cycle;
  \draw[thickline] (A)--(D)--(B);
  \draw[thickline] (C)--(D);
  \node[above] at (A) {$A$};
  \node[left] at (B) {$B$};
  \node[right] at (C) {$C$};
  \node[above right] at (D) {$D$};
  \node at (-1.42,0.62) {$C$};
  \node at (1.42,0.62) {$L$};
  \node[right] at (0.10,1.20) {$R$};
  \node[below] at (0.55,-1.38) {$R$};
  \node[left] at (-1.12,-0.80) {$L$};
  \node[right] at (1.12,-0.80) {$C$};
  % loop currents
  \draw[force] (-1.05,-0.62) arc (150:30:0.62);
  \node at (-0.22,0.10) {$I_1$};
  \draw[force] (-1.55,0.62) arc (140:20:0.55);
  \node at (-1.05,1.10) {$I_2$};
  \draw[force] (1.55,0.62) arc (40:-80:0.55);
  \node at (1.02,1.10) {$I_3$};
  % source connected across resistor BC
  \draw[thickline] (B)--(-2.25,-2.95)--(-0.35,-2.95);
  \draw[thickline] (C)--(2.25,-2.95)--(0.35,-2.95);
  \draw[thickline,fill=white] (0,-2.95) circle (0.35);
  \node at (0,-2.95) {$\sim$};
  \node[below] at (0,-3.36) {$V_0$};
  \draw[force] (1.30,-1.72) arc (70:-60:0.55);
  \node at (1.62,-1.62) {$I_4$};
\end{tikzpicture}
\end{figbox}

According to Kirchhoff's laws, the loop equations are as follows:
\[
(I_1-I_4)R+(I_1-I_2)(jR)+(I_1-I_3)(-jR)=0,
\]
\[
I_2(-jR)+(I_2-I_3)R+(I_2-I_1)(jR)=0,
\]

\[
I_3(jR)+(I_3-I_1)(-jR)+(I_3-I_2)R=0,
\]
\[
(I_4-I_1)R=V_0.
\]

Simplifying gives
\[
I_1-I_4+j(I_3-I_2)=0,
\]
\[
I_2-I_3-jI_1=0,
\]
\[
I_3-I_2+jI_1=0,
\]
\[
I_4-I_1=V_0/R.
\]
The second and third equations are in fact the same. Multiplying the second equation by $j$ and adding it to the first gives $I_1=I_4/2$; substituting this into the fourth equation gives the required total current
\[
I_4=2V_0/R.
\]
Thus the equivalent impedance of the circuit is $R/2$: the five branches above are equivalent to a single resistance $R$, connected in parallel with the resistance $R$ at the bottom.
\end{solution}

\begin{solution}{3.52}
(1) The equivalent circuit is as shown in the figure, where $\tilde{Z}_1=\dfrac{1}{j\omega C}$ and $\tilde{Z}_2=jL\omega$. Let the complex impedance between the terminals $A$ and $B$ be $\tilde{Z}$. By the conditions of the problem,
\[
\tilde{Z}=\tilde{Z}_1+\frac{\tilde{Z}_2\tilde{Z}}{\tilde{Z}+\tilde{Z}_2},
\]
that is,
\[
\tilde{Z}^2-\tilde{Z}_1\tilde{Z}-\tilde{Z}_1\tilde{Z}_2=0.
\]
Solving gives
\[
\tilde{Z}=\frac{\tilde{Z}_1}{2}\left(1+\sqrt{1+4\frac{\tilde{Z}_2}{\tilde{Z}_1}}\right)
=\frac{1}{2j\omega C}\left(1+\sqrt{1-4LC\omega^2}\right).
\]
When $4LC\omega^2\le1$, i.e. $\omega\le\dfrac{1}{2\sqrt{LC}}$, $\tilde{Z}$ is purely imaginary. When $4LC\omega^2>1$, i.e. $\omega>\dfrac{1}{2\sqrt{LC}}$, $\tilde{Z}$ has both imaginary and real parts. Since the real part cannot be negative, only the ``$+$'' sign can be taken before the square root in the expression for $\tilde{Z}$.

(2) The average power is
\[
\bar{P}=\frac{1}{2}\mathrm{Re}\left[V_0I_0^*\right]
=\frac{1}{2}|V_0|^2\mathrm{Re}\left[\frac{1}{\tilde{Z}}\right].
\]
\begin{figbox}{3.52}
\begin{tikzpicture}[scale=0.95,>=latex]
  \draw[thickline] (-3.60,1.50)--(-1.20,1.50);
  \draw[thickline,fill=white] (-2.35,1.34) rectangle (-1.55,1.66);
  \node[above] at (-1.95,1.70) {$\tilde{Z}_1$};
  \node[left] at (-3.68,1.50) {$A$};
  \draw[thickline] (-1.20,1.50)--(2.70,1.50);
  \draw[thickline] (1.40,1.50)--(1.40,0.95);
  \draw[thickline,fill=white] (1.24,0.55) rectangle (1.56,0.95);
  \node[right] at (1.60,0.75) {$\tilde{Z}_2$};
  \draw[thickline] (1.40,0.55)--(1.40,-1.50);
  \draw[thickline] (2.70,1.50)--(2.70,0.95);
  \draw[thickline,fill=white] (2.54,0.45) rectangle (2.86,0.95);
  \node[right] at (2.92,0.70) {$\tilde{Z}$};
  \draw[thickline] (2.70,0.45)--(2.70,-1.50);
  \draw[thickline] (-3.60,-1.50)--(2.70,-1.50);
  \node[left] at (-3.68,-1.50) {$B$};
  \draw[thickline,fill=white] (-3.60,1.50) circle (0.055);
  \draw[thickline,fill=white] (-3.60,-1.50) circle (0.055);
  \fill (-1.20,1.50) circle (0.04);
  \fill (2.70,1.50) circle (0.04);
\end{tikzpicture}
\end{figbox}

Hence when
\[
\omega\le\frac{1}{2\sqrt{LC}},
\]
\[
\bar{P}=0,
\]
and when
\[
\omega>\frac{1}{2\sqrt{LC}},
\]

\[
\bar{P}=\frac{1}{2}V_0^2\,\mathrm{Re}\left[\frac{1}{\tilde{Z}}\right].
\]
From part (1),
\[
\frac{1}{\tilde{Z}}
=\frac{2j\omega C}{1+j\sqrt{4LC\omega^2-1}}
=\frac{2j\omega C\left(1-j\sqrt{4LC\omega^2-1}\right)}{4LC\omega^2},
\]
\[
\therefore\ \mathrm{Re}\left[\frac{1}{\tilde{Z}}\right]
=\frac{C\omega}{2LC\omega^2}\sqrt{4LC\omega^2-1}
=\frac{1}{2L\omega}\sqrt{4LC\omega^2-1},
\]
\[
\therefore\ \bar{P}=\frac{V_0^2}{4L\omega}\sqrt{4LC\omega^2-1}.
\]
\end{solution}

\begin{solution}{3.53}
(1) Since the quasi-steady condition is satisfied, at any instant the charge on the two plates is distributed uniformly. The electric field is
\[
E=\frac{V}{d}=\frac{V_0}{d}\sin\omega t.
\]
The magnetic field is excited by the total current. By the Amp\`ere circuital law,
\[
B\cdot2\pi r=\mu_0\varepsilon_0\frac{\partial E}{\partial t}\pi r^2.
\]
From the expression for $E$,
\[
\frac{\partial E}{\partial t}=\frac{V_0\omega}{d}\cos\omega t.
\]
Substituting into the above expression gives
\[
B=\frac{\mu_0\varepsilon_0V_0\omega}{2d}\,r\cos\omega t.
\]

(2) At the edge of the capacitor, $r=a$, we have
\[
E=\frac{V_0}{d}\sin\omega t,
\]
\[
B=\frac{\mu_0\varepsilon_0V_0\omega}{2d}\,a\cos\omega t.
\]
The energy flux density is
\[
S=EH=\frac{1}{\mu_0}EB=\frac{\varepsilon_0V_0^2\omega a}{2d^2}\sin\omega t\cos\omega t.
\]
The energy flux entering the capacitor through its edge is
\[
I=S\cdot2\pi ad=\frac{\pi a^2}{d}\varepsilon_0V_0^2\omega\sin\omega t\cos\omega t.
\]

(3) The electric field energy between the plates is
\[
W_e=\int\frac{1}{2}\varepsilon_0E^2\,\dd\tau=\frac{1}{2}\varepsilon_0E^2\cdot\pi a^2d,
\]
\[
\frac{\dd W_e}{\dd t}=\varepsilon_0E\frac{\partial E}{\partial t}\pi a^2d
=\varepsilon_0\frac{V_0}{d}\sin\omega t\cdot\frac{V_0\omega}{d}\cos\omega t\cdot\pi a^2d
\]
\[
=\frac{\pi a^2}{d}\varepsilon_0V_0^2\omega\sin\omega t\cos\omega t
=I.
\]
The magnetic field energy is
\begin{align*}
W_m&=\int\frac{1}{2\mu_0}B^2\,\dd\tau
=\frac{\mu_0\varepsilon_0^2V_0^2\omega^2}{8\mu_0d^2}\cos^2\omega t\,r^2\cdot2\pi r\,\dd r\cdot d\\
&=\frac{\mu_0\varepsilon_0^2V_0^2\omega^2\pi}{4d}\cos^2\omega t\int_0^a r^3\,\dd r\\
&=\frac{\mu_0\varepsilon_0^2V_0^2\omega^2\pi a^4}{16d}\cos^2\omega t,\\
\frac{\dd W_m}{\dd t}&=-\frac{\mu_0\varepsilon_0^2V_0^2\omega^3\pi a^4}{8d}\sin\omega t\cos\omega t.
\end{align*}
According to the conditions of the problem, $a\ll c/\omega$, and noting that $\varepsilon_0\mu_0=1/c^2$, we have
\[
\left|\frac{\dd W_m/\dd t}{\dd W_e/\dd t}\right|
=\frac{1}{8}\frac{\omega^2a^2}{c^2}\ll1,
\]
\[
\therefore\ \frac{\dd W_e}{\dd t}+\frac{\dd W_m}{\dd t}\approx\frac{\dd W_e}{\dd t}=I.
\]
It is thus seen that under the quasi-steady condition the energy flowing in through the edge of the capacitor is equal to the rate of increase of the electromagnetic field energy between the two plates.
\end{solution}

\begin{solution}{3.54}
(1) Solve by the complex method. Let the current flowing through $R_1$ be $\tilde{I}_1$. The input impedance of the circuit is
\[
\tilde{Z}=R_1+\frac{\left(-j\dfrac{1}{\omega C}\right)R_2}{R_2-j\dfrac{1}{\omega C}},
\]
\[
\tilde{I}_1=\frac{\tilde{V}}{\tilde{Z}},\qquad
\tilde{V}_1=\tilde{I}_1R_1=\frac{R_1}{\tilde{Z}}\tilde{V}.
\]
Substituting the expression for $\tilde{Z}$,
\[
\tilde{V}_1=\frac{R_1}{R_1+\dfrac{\left(-j\dfrac{1}{\omega C}\right)R_2}{R_2-j\dfrac{1}{\omega C}}}\,\tilde{V}
=\frac{R_1\left(R_2-j\dfrac{1}{\omega C}\right)}{R_1R_2-j\dfrac{R_1+R_2}{\omega C}}\,\tilde{V}
\]
\[
=\frac{R_1\sqrt{R_2^2+\dfrac{1}{\omega^2C^2}}}
{\sqrt{R_1^2R_2^2+\left(\dfrac{R_1+R_2}{\omega C}\right)^2}}
\cdot\frac{\ee^{-j\varphi_1}}{\ee^{-j\varphi_2}}\,\tilde{V}
=\frac{R_1\sqrt{R_2^2+\dfrac{1}{\omega^2C^2}}}
{\sqrt{R_1^2R_2^2+\left(\dfrac{R_1+R_2}{\omega C}\right)^2}}
\,\ee^{j(\varphi_2-\varphi_1)}\tilde{V}.
\]
\[
\therefore\ V_1=V_0\frac{R_1\sqrt{R_2^2+\dfrac{1}{\omega^2C^2}}}
{\sqrt{R_1^2R_2^2+\left(\dfrac{R_1+R_2}{\omega C}\right)^2}}\cos(\omega t+\varphi),
\]
where
\begin{align*}
\tan\varphi&=\tan(\varphi_2-\varphi_1)=\frac{\tan\varphi_2-\tan\varphi_1}{1+\tan\varphi_2\tan\varphi_1}\\
&=\frac{\dfrac{R_1+R_2}{\omega C}\cdot\dfrac{1}{R_1R_2}-\dfrac{1}{\omega CR_2}}
{1+\dfrac{R_1+R_2}{\omega C}\cdot\dfrac{1}{R_1R_2}\cdot\dfrac{1}{\omega CR_2}}\\
&=\frac{\omega CR_2^2}{R_1+R_2+\omega^2C^2R_1R_2^2}.
\end{align*}

(2) When the pulse is just applied, $C$ may be regarded as a short circuit (the proof is given below); hence $C$ is charged with a current $U_0/R_1$, and after a duration $\Delta t$ the charge accumulated at the two ends of the capacitor is
\[
q_0=\frac{U_0}{R_1}\Delta t.
\]
Afterwards the capacitor discharges through the parallel resistors $R_1$ and $R_2$, and its charge varies with time as
\[
q=q_0\ee^{-t'/\frac{R_1R_2}{R_1+R_2}C},
\]
\[
\therefore\ V_1=\frac{q}{C}=\frac{q_0}{C}\ee^{-t'\frac{R_1+R_2}{R_1R_2C}}
=\frac{U_0\Delta t}{R_1C}\ee^{-t'\frac{R_1+R_2}{R_1R_2C}}.
\]
Here $t'$ is the time of discharge of the capacitor through the parallel resistors $R_1$ and $R_2$. In summary, the result is
\[
V_1=\begin{cases}
0, & t<0,\\[4pt]
U_0, & 0<t<\Delta t,\\[4pt]
-\dfrac{U_0\Delta t}{R_1C}\ee^{-(t-\Delta t)\frac{R_1+R_2}{R_1R_2C}}, & t>\Delta t.
\end{cases}
\]
The negative sign indicates that at this time the potential at the right end of $R_1$ is higher than that at the left end.

The proof that $C$ may be regarded as a short circuit when $V(t)$ is a short pulse is as follows. When $V=U_0$ (constant), we have (see the figure)
\begin{align}
U_0&=i_1R_1+i_2R_2 \tag{1}\\
i_1-i_2&=\frac{\dd q}{\dd t} \tag{2}\\
\frac{q}{C}&=i_2R_2 \tag{3}
\end{align}
Differentiating (3) and substituting into (2) gives
\[
i_1-i_2=R_2C\frac{\dd i_2}{\dd t}.
\]
Multiplying both sides by $R_1$ and using (1) gives
\begin{equation}
U_0-(R_1+R_2)i_2=R_1R_2C\frac{\dd i_2}{\dd t}, \tag{4}
\end{equation}
that is,
\[
\frac{\dd i_2}{U_0-(R_1+R_2)i_2}=\frac{\dd t}{R_1R_2C}.
\]
Integrating the above expression gives
\begin{equation}
i_2=\frac{U_0}{R_1+R_2}\left(1-\ee^{-t/RC}\right), \tag{5}
\end{equation}
where $R=\dfrac{R_1R_2}{R_1+R_2}$. It is seen that as $t\to0$, $i_2\to0$. Substituting (5) into (1) gives
\[
i_1=\frac{U_0}{R_1}-\frac{U_0R_2}{R_1(R_1+R_2)}\left(1-\ee^{-t/RC}\right).
\]
It is seen that as $t\to0$, $i_1\approx\dfrac{U_0}{R_1}$, and the charging current is
\[
i_c=i_1-i_2\approx i_1=\frac{U_0}{R_1}.
\]
This proves the statement.

\begin{figbox}{3.54}
\begin{tikzpicture}[scale=0.95,>=latex]
  \draw[thickline] (-3.60,1.20)--(-3.60,-1.20);
  \draw[thickline,fill=white] (-3.60,1.20) circle (0.06);
  \draw[thickline,fill=white] (-3.60,-1.20) circle (0.06);
  \draw[thickline] (-3.60,1.20)--(-2.95,1.20);
  \draw[thickline,fill=white] (-2.95,1.03) rectangle (-2.15,1.37);
  \node[above] at (-2.55,1.40) {$R_1$};
  \draw[thickline] (-2.15,1.20)--(-0.40,1.20);
  \draw[force] (-1.95,1.20)--(-1.30,1.20);
  \node[above] at (-1.62,1.24) {$i_1$};
  \fill (-1.30,1.20) circle (0.04);
  \draw[force] (-0.25,1.20)--(0.60,1.20);
  \node[above] at (0.18,1.24) {$i_2$};
  \draw[thickline] (-0.40,1.20)--(1.80,1.20);
  \draw[thickline] (1.80,1.20)--(1.80,0.95);
  \draw[thickline,fill=white] (1.64,0.35) rectangle (1.96,0.95);
  \node[right] at (2.02,0.65) {$R_2$};
  \draw[thickline] (1.80,0.35)--(1.80,-1.20);
  \draw[thickline] (-3.60,-1.20)--(1.80,-1.20);
  \draw[thickline,fill=white] (-0.40,1.20) circle (0.05);
  \draw[thickline] (-0.40,1.20)--(-0.40,0.30);
  \draw[thickline] (-0.78,0.30)--(-0.02,0.30);
  \draw[thickline] (-0.78,0.02)--(-0.02,0.02);
  \draw[thickline] (-0.40,0.02)--(-0.40,-1.20);
  \node[left] at (-0.84,0.16) {$C$};
  \node[left] at (-0.86,0.72) {$q$};
  \draw[force] (-0.40,1.05)--(-0.40,0.52);
  \node[right] at (-0.30,0.80) {$i_1-i_2$};
\end{tikzpicture}
\end{figbox}
\end{solution}

\begin{solution}{3.55}
(1) After $K$ is closed, let the charges on the capacitors be $Q_1$ and $Q_2$ respectively, and let the current flowing through the resistor and the inductor be $I$ (see the figure). Then the circuit equation is
\[
IR+L\frac{\dd I}{\dd t}+\frac{Q_2}{C_2}=\frac{Q_1}{C_1}.
\]
\begin{figbox}{3.55}
\begin{tikzpicture}[scale=0.95,>=latex]
  \draw[thickline] (-3.20,1.50)--(-3.20,0.24);
  \draw[thickline] (-3.55,0.24)--(-2.85,0.24);
  \draw[thickline] (-3.55,-0.24)--(-2.85,-0.24);
  \draw[thickline] (-3.20,-0.24)--(-3.20,-1.50);
  \node[left] at (-3.60,0) {$C_1$};
  \node[right] at (-2.80,0.62) {$Q_1$};
  \node[right] at (-2.80,-0.62) {$-Q_1$};
  \draw[thickline] (3.20,1.50)--(3.20,0.24);
  \draw[thickline] (2.85,0.24)--(3.55,0.24);
  \draw[thickline] (2.85,-0.24)--(3.55,-0.24);
  \draw[thickline] (3.20,-0.24)--(3.20,-1.50);
  \node[right] at (3.60,0) {$C_2$};
  \node[left] at (2.80,0.62) {$Q_2$};
  \node[left] at (2.80,-0.62) {$-Q_2$};
  \draw[thickline] (-3.20,1.50)--(3.20,1.50);
  \draw[thickline,fill=white] (-1.30,1.33) rectangle (-0.50,1.67);
  \node[above] at (-0.90,1.72) {$R$};
  \foreach \x in {0.35,0.65,0.95,1.25} {
    \draw[thickline] (\x-0.15,1.50) arc (180:0:0.15);
  }
  \node[above] at (0.80,1.68) {$L$};
  \draw[force] (-2.60,1.50)--(-1.95,1.50);
  \node[above] at (-2.28,1.55) {$I$};
  \draw[force] (1.85,1.50)--(2.50,1.50);
  \node[above] at (2.18,1.55) {$I$};
  \draw[thickline] (-3.20,-1.50)--(3.20,-1.50);
  \draw[thickline] (0,-1.50)--(0,-1.85);
  \draw[thickline] (-0.32,-1.85)--(0.32,-1.85);
  \draw[thickline] (-0.20,-1.97)--(0.20,-1.97);
  \draw[thickline] (-0.09,-2.09)--(0.09,-2.09);
\end{tikzpicture}
\end{figbox}

Differentiating both sides with respect to $t$, and noting that $\dfrac{\dd Q_2}{\dd t}=I$ and $-\dfrac{\dd Q_1}{\dd t}=I$, we obtain
\[
L\frac{\dd^2I}{\dd t^2}+R\frac{\dd I}{\dd t}+\left(\frac{1}{C_1}+\frac{1}{C_2}\right)I=0.
\]
The roots of the characteristic equation corresponding to this equation are
\[
D=-\frac{R}{2L}\pm\sqrt{\frac{R^2}{4L^2}-\frac{1}{L}\left(\frac{1}{C_1}+\frac{1}{C_2}\right)}.
\]
According to the conditions of the problem, the process should be in the critically damped state; hence
\[
\frac{R^2}{4L^2}=\frac{1}{L}\left(\frac{1}{C_1}+\frac{1}{C_2}\right).
\]
Solving gives
\[
R=R_0=2\sqrt{L\left(\frac{1}{C_1}+\frac{1}{C_2}\right)}.
\]
(2) At this time the solution of the original equation is
\[
I=(a+bt)\ee^{-\frac{R}{2L}t}.
\]
Here $a$ and $b$ are determined by the boundary conditions. When $t=0$, $I=0$, so $a=0$, and
\[
I=bt\,\ee^{-\frac{R}{2L}t}.
\]
To find the maximum of $I$, set $\dfrac{\dd I}{\dd t}=0$:
\[
\frac{\dd I}{\dd t}=b\ee^{-\frac{R_0}{2L}t}+bt\left(-\frac{R_0}{2L}\right)\ee^{-\frac{R_0}{2L}t}
=b\left(1-\frac{R_0t}{2L}\right)\ee^{-\frac{R_0}{2L}t}=0,
\]
which gives
\[
t=\frac{2L}{R_0}=\sqrt{\frac{LC_1C_2}{C_1+C_2}}.
\]
When $t=\dfrac{2L}{R_0}$,
\[
\frac{\dd^2I}{\dd t^2}
=-\frac{R_0b}{2L}\ee^{-\frac{R_0}{2L}t}-\frac{R_0b}{2L}\ee^{-\frac{R_0}{2L}t}
+\frac{R_0^2}{4L^2}bt\,\ee^{-\frac{R_0}{2L}t}
\]
\[
=-\frac{R_0b}{L}\ee^{-\frac{R_0}{2L}t}+\frac{R_0^2}{4L^2}bt\,\ee^{-\frac{R_0}{2L}t}.
\]
Substituting $t=\dfrac{2L}{R_0}$ into the above expression gives
\[
\left.\frac{\dd^2I}{\dd t^2}\right|_{t=\frac{2L}{R_0}}
=\left(-\frac{R_0b}{L}+\frac{R_0^2}{4L^2}b\frac{2L}{R_0}\right)\ee^{-1}
=-\frac{R_0b}{2L}\frac{1}{\ee}<0.
\]
It is thus seen that $I$ has a maximum when $t=\dfrac{2L}{R_0}$.

To find the maximum value $I_m$ of $I$, the value of $b$ must be found. This can be obtained from the integral and the final-state condition:
\begin{align*}
\int_0^\infty I\,\dd t&=\int_0^\infty bt\,\ee^{-\frac{R_0}{2L}t}\,\dd t
=b\left(-\frac{2L}{R_0}\right)^2\int_0^\infty\left(-\frac{R_0t}{2L}\right)\ee^{-\frac{R_0t}{2L}}\,\dd\left(-\frac{R_0t}{2L}\right)\\
&=b\frac{4L^2}{R_0^2}\left[-\frac{R_0t}{2L}\ee^{-\frac{R_0t}{2L}}-\ee^{-\frac{R_0t}{2L}}\right]_0^\infty
=b\frac{4L^2}{R_0^2}.
\end{align*}
On the other hand, $\displaystyle\int_0^\infty I\,\dd t=Q_2=Q_0\frac{C_2}{C_1+C_2}$; hence
\[
Q_0\frac{C_2}{C_1+C_2}=b\frac{4L^2}{R_0^2}.
\]
Substituting the expression for $R_0$ gives
\[
b=Q_0\frac{C_2}{C_1+C_2}\frac{R_0^2}{4L^2}=\frac{Q_0}{LC_1}.
\]
Substituting into the expression for $I$ gives
\[
I=\frac{Q_0}{LC_1}t\,\ee^{-\frac{R_0}{2L}t}.
\]
Substituting $t=\dfrac{2L}{R_0}$ into the above expression, and inserting the expression for $R_0$, the maximum value of the current is obtained:
\[
I_m=\frac{Q_0}{LC_1}\frac{2L}{R_0}\frac{1}{\ee}
=\frac{2Q_0}{\ee R_0C_1}
=\frac{Q_0}{\ee C_1\sqrt{L\left(\dfrac{1}{C_1}+\dfrac{1}{C_2}\right)}}.
\]
\end{solution}

\begin{solution}{3.56}
(1) Obviously, the propagation factor is
\[
\alpha=\frac{\dfrac{1}{j\omega C}}{j\omega L+\dfrac{1}{j\omega C}}
=\frac{1}{1-LC\omega^2}.
\]
For $|\alpha|<1$ we must have $LC\omega^2>2$, that is,
\[
\omega>\sqrt{\frac{2}{LC}}.
\]

(2) According to the conditions of the problem, we must have
\[
jL\omega+\frac{1}{j\omega C+\dfrac{1}{\tilde{Z}_0}}=\tilde{Z}_0,
\]
that is,
\[
\tilde{Z}_0^2-jL\omega\tilde{Z}_0-\frac{L}{C}=0.
\]
Solving gives
\[
\tilde{Z}_0=\frac{1}{2}jL\omega+\sqrt{\frac{L}{C}-\frac{\omega^2L^2}{4}}.
\]
When $\dfrac{L}{C}>\dfrac{\omega^2L^2}{4}$, the quantity under the square root is real. Since the real part cannot be negative, only the ``$+$'' sign can be taken before the square root.

(3) Consider the propagation factor $\alpha_n=U_{n+1}/U_n$ between two adjacent sections ($n$ being the serial number from left to right). Obviously,
\[
\alpha_n=\frac{U_{n+1}}{U_n}=\frac{\tilde{Z}_0-jL\omega}{\tilde{Z}_0}.
\]
Substituting the expression for $\tilde{Z}_0$,
\[
\alpha_n=\frac{\sqrt{\dfrac{L}{C}-\dfrac{\omega^2L^2}{4}}-\dfrac{1}{2}jL\omega}
{\sqrt{\dfrac{L}{C}-\dfrac{\omega^2L^2}{4}}+\dfrac{1}{2}jL\omega}.
\]
When $\dfrac{L}{C}>\dfrac{\omega^2L^2}{4}$, i.e. $\omega<\dfrac{2}{\sqrt{LC}}$, $|\alpha_n|=1$, and the electromagnetic disturbance passes through the ladder network without attenuation.
\[
\therefore\ \omega_c=\frac{2}{\sqrt{LC}}.
\]

(4) When $\omega<\omega_c$, we have
\[
\alpha_n=\frac{\sqrt{\dfrac{L}{C}-\dfrac{\omega^2L^2}{4}}-\dfrac{1}{2}jL\omega}
{\sqrt{\dfrac{L}{C}-\dfrac{\omega^2L^2}{4}}+\dfrac{1}{2}jL\omega}
=\ee^{-j2\delta_1}=\ee^{-j\delta}.
\]
Here
\[
\tan\delta_1=\frac{\dfrac{1}{2}L\omega}{\sqrt{\dfrac{L}{C}-\dfrac{\omega^2L^2}{4}}},\qquad \delta=2\delta_1,
\]
so
\[
\tan\delta=\frac{2\tan\delta_1}{1-\tan^2\delta_1}
=\frac{\dfrac{L\omega}{\sqrt{\dfrac{L}{C}-\dfrac{\omega^2L^2}{4}}}}
{1-\dfrac{\dfrac{1}{4}\omega^2L^2}{\dfrac{L}{C}-\dfrac{\omega^2L^2}{4}}}
=\frac{L\omega\sqrt{\dfrac{L}{C}-\dfrac{\omega^2L^2}{4}}}{\dfrac{L}{C}-\dfrac{1}{2}\omega^2L^2}.
\]
When $\omega\ll\omega_c$, $\dfrac{\omega^2L^2}{4}\ll\dfrac{L}{C}$, and
\[
\tan\delta\approx\frac{L\omega\sqrt{\dfrac{L}{C}}}{\dfrac{L}{C}}=\omega\sqrt{LC}\ll1,
\]
\[
\therefore\ \delta\approx\omega\sqrt{LC}.
\]
The time required for the disturbance to propagate from one section to the next is
\[
\Delta t=\frac{\delta}{\omega}=\sqrt{LC}.
\]
Hence the propagation speed of the signal along the network is
\[
v=\frac{\Delta r}{\Delta t}=\frac{\Delta r}{\sqrt{LC}}.
\]
\end{solution}

\begin{solution}{3.57}
Using the $Y$-$\Delta$ transformation, the original circuit can be converted into the circuit shown in Fig.~3.57(a), in which the unprimed quantities are the transformed complex impedances of the lower T-shaped circuit (two $C$'s and one $R/2$) in the original circuit, and the primed quantities are the transformed complex impedances of the upper T-shaped circuit (two $R$'s and one $C'$) in the original circuit. By symmetry, $\tilde{Z}_2=\tilde{Z}_3$ and $\tilde{Z}_2'=\tilde{Z}_3'$. Moreover, the parallel combination of $\tilde{Z}_2$ and $\tilde{Z}_2'$ cannot be zero (otherwise the input would be short-circuited); hence, for the output voltage $U_{\mathrm{out}}$ to be zero, the parallel impedance of $\tilde{Z}_1$ and $\tilde{Z}_1'$ can only be made infinite. The parallel impedance of $\tilde{Z}_1$ and $\tilde{Z}_1'$ is
\[
\tilde{Z}=\frac{\tilde{Z}_1\tilde{Z}_1'}{\tilde{Z}_1+\tilde{Z}_1'}.
\]
For $\tilde{Z}\to\infty$, we can only have $\tilde{Z}_1+\tilde{Z}_1'=0$.

\begin{figbox}{3.57(a)}
\begin{tikzpicture}[scale=0.92,>=latex]
  \coordinate (A) at (0,2.20);
  \coordinate (B) at (4.80,2.20);
  \coordinate (M) at (2.40,-1.30);
  % top branch: two parallel paths
  \draw[thickline] (A)--(0,2.55)--(0.55,2.55);
  \draw[thickline,fill=white] (0.55,2.39) rectangle (1.45,2.71);
  \node[above] at (1.00,2.75) {$\tilde{Z}_1'$};
  \draw[thickline] (1.45,2.55)--(4.80,2.55)--(B);
  \draw[thickline] (A)--(0,1.95)--(1.15,1.95);
  \draw[thickline,fill=white] (1.15,1.79) rectangle (2.05,2.11);
  \node[below] at (1.60,1.75) {$\tilde{Z}_1$};
  \draw[thickline] (2.05,1.95)--(4.80,1.95)--(4.80,2.55);
  % left arm: two parallel branches
  \draw[thickline] (A)--(2.50,-1.23);
  \draw[thickline] (A)--(2.30,-1.37);
  \begin{scope}[rotate around={-55.6:(1.30,0.52)}]
    \draw[thickline,fill=white] (1.05,0.41) rectangle (1.55,0.63);
  \end{scope}
  \begin{scope}[rotate around={-55.6:(1.10,0.38)}]
    \draw[thickline,fill=white] (0.85,0.27) rectangle (1.35,0.49);
  \end{scope}
  \node at (1.95,0.92) {$\tilde{Z}_2'$};
  \node at (0.72,0.02) {$\tilde{Z}_2$};
  % right arm: two parallel branches
  \draw[thickline] (B)--(2.50,-1.37);
  \draw[thickline] (B)--(2.30,-1.23);
  \begin{scope}[rotate around={55.6:(3.70,0.38)}]
    \draw[thickline,fill=white] (3.45,0.27) rectangle (3.95,0.49);
  \end{scope}
  \begin{scope}[rotate around={55.6:(3.50,0.52)}]
    \draw[thickline,fill=white] (3.25,0.41) rectangle (3.75,0.63);
  \end{scope}
  \node at (3.02,0.92) {$\tilde{Z}_3'$};
  \node at (4.10,-0.02) {$\tilde{Z}_3$};
  % bottom rail and ports
  \draw[thickline] (-0.60,-1.30)--(5.40,-1.30);
  \draw[thickline] (A)--(-0.60,2.20);
  \draw[thickline] (B)--(5.40,2.20);
  \fill (-0.60,2.20) circle (0.05);
  \fill (-0.60,-1.30) circle (0.05);
  \fill (5.40,2.20) circle (0.05);
  \fill (5.40,-1.30) circle (0.05);
  \draw[<->] (-1.20,2.20)--(-1.20,-1.30) node[midway,left] {$U_{\mathrm{in}}$};
  \draw[<->] (6.00,2.20)--(6.00,-1.30) node[midway,right] {$U_{\mathrm{out}}$};
\end{tikzpicture}
\end{figbox}

\begin{figbox}{3.57(b)}
\begin{tikzpicture}[scale=0.92,>=latex]
  % Y network
  \coordinate (O) at (-1.20,0.00);
  \coordinate (T3) at (-2.30,1.15);
  \coordinate (T2) at (-0.10,1.15);
  \coordinate (T1) at (-1.20,-1.15);
  \draw[thickline] (O)--(T3);
  \draw[thickline] (O)--(T2);
  \draw[thickline] (O)--(T1);
  \node[above left] at (T3) {$3$};
  \node[above right] at (T2) {$2$};
  \node[below] at (T1) {$1$};
  \fill (T3) circle (0.05);
  \fill (T2) circle (0.05);
  \fill (T1) circle (0.05);
  \node at (-2.00,0.45) {$\tilde{z}_3$};
  \node at (-0.42,0.45) {$\tilde{z}_2$};
  \node at (-1.42,-0.55) {$\tilde{z}_1$};
  % Delta network
  \coordinate (D3) at (1.90,1.15);
  \coordinate (D2) at (3.60,1.15);
  \coordinate (D1) at (2.75,-1.15);
  \draw[thickline] (D3)--(D2)--(D1)--cycle;
  \node[above left] at (D3) {$3$};
  \node[above right] at (D2) {$2$};
  \node[below] at (D1) {$1$};
  \fill (D3) circle (0.05);
  \fill (D2) circle (0.05);
  \fill (D1) circle (0.05);
  \node[above] at (2.75,1.15) {$\tilde{Z}_1$};
  \node[right] at (2.98,0.45) {$\tilde{Z}_2$};
  \node[left] at (2.52,0.45) {$\tilde{Z}_3$};
\end{tikzpicture}
\end{figbox}

Corresponding to the d.c. $Y$-$\Delta$ transformation (Fig.~3.57(b)), we have
\begin{align*}
\tilde{Z}_1&=\frac{\tilde{z}_1\tilde{z}_2+\tilde{z}_2\tilde{z}_3+\tilde{z}_3\tilde{z}_1}{\tilde{z}_1}
=\tilde{z}_2+\tilde{z}_3+\frac{\tilde{z}_2\tilde{z}_3}{\tilde{z}_1}\\
&=\frac{1}{j\omega C}+\frac{1}{j\omega C}
+\frac{\dfrac{1}{j\omega C}\cdot\dfrac{1}{j\omega C}}{R/2}\\
&=-j\frac{2}{\omega C}-\frac{2}{R\omega^2C^2}.
\end{align*}
Similarly,
\[
\tilde{Z}_1'=R+R+\frac{R^2}{\dfrac{1}{j\omega C'}}=2R+jR^2\omega C'.
\]
For $\tilde{Z}_1+\tilde{Z}_1'=0$, we have
\[
2R=\frac{2}{R\omega^2C^2},
\]
\[
R^2\omega C'=\frac{2}{\omega C}.
\]
From these we obtain
\[
\omega=\frac{1}{RC},
\]
\[
C'=2C.
\]
At this time,
\begin{align*}
\tilde{Z}_2&=\frac{\tilde{z}_1\tilde{z}_2+\tilde{z}_2\tilde{z}_3+\tilde{z}_3\tilde{z}_1}{\tilde{z}_2}
=\tilde{z}_1+\tilde{z}_3+\frac{\tilde{z}_1\tilde{z}_3}{\tilde{z}_2}\\
&=\frac{R}{2}+\frac{1}{j\omega C}
+\frac{\dfrac{R}{2}\cdot\dfrac{1}{j\omega C}}{\dfrac{1}{j\omega C}}
=R-j\frac{1}{\omega C}=\tilde{Z}_3.
\end{align*}

Similarly,
\[
\tilde{Z}_2'=\frac{1}{j\omega C'}+R+\frac{\dfrac{1}{j\omega C'}\cdot R}{R}
=R-j\frac{2}{\omega C'}=\tilde{Z}_3'.
\]
The parallel impedance of $\tilde{Z}_2$ and $\tilde{Z}_2'$, and that of $\tilde{Z}_3$ and $\tilde{Z}_3'$, are both
\[
\frac{1}{2}\left(R-j\frac{2}{\omega C'}\right)\neq0.
\]
\end{solution}

\begin{solution}{3.58}
(1) The magnetic field of the incident wave must be along the $y$ direction, and the magnetic field strength can be written as
\[
\vec{H}_i(z,t)=\sqrt{\frac{\varepsilon_0}{\mu_0}}E_{i0}\cos(\omega t-kz)\,\hat{j}.
\]
Here use has been made of
\[
H=\frac{1}{\mu_0}B=\frac{1}{\mu_0c}E=\sqrt{\frac{\varepsilon_0}{\mu_0}}E.
\]
The reflected wave must propagate along the negative $z$ direction and can be written as
\[
\vec{E}_r(z,t)=E_{r0}\cos(\omega t+kz)\,\hat{i},
\]
\[
\vec{H}_r(z,t)=-H_{r0}\cos(\omega t+kz)\,\hat{j}.
\]
The minus sign in front of the expression for $H_r$ is due to the fact that the Poynting vector $\vec{S}_r=\vec{E}_r\times\vec{H}_r$ of the reflected wave must be along the negative $z$ direction. After the incident and reflected waves are superposed at the conductor surface, since $E=0$ inside an ideal conductor while the tangential component of the electric field must be continuous, we obtain
\[
E_{i0}+E_{r0}=0,
\]
that is,
\[
E_{r0}=-E_{i0}.
\]
Hence the total electric field outside the conductor is
\begin{align*}
E_x(z,t)&=E_i(z,t)+E_r(z,t)\\
&=E_{i0}\cos(\omega t-kz)+E_{r0}\cos(\omega t+kz)\\
&=E_{i0}\cos(\omega t-kz)-E_{i0}\cos(\omega t+kz)\\
&=E_{i0}\bigl[(\cos\omega t\cos kz+\sin\omega t\sin kz)-(\cos\omega t\cos kz-\sin\omega t\sin kz)\bigr]\\
&=2E_{i0}\sin kz\sin\omega t.
\end{align*}
The total magnetic field outside the conductor is then
\begin{align*}
H_y(z,t)&=H_i(z,t)+H_r(z,t)\\
&=H_{i0}\cos(\omega t-kz)-H_{r0}\cos(\omega t+kz)\\
&=\sqrt{\frac{\varepsilon_0}{\mu_0}}E_{i0}\cos(\omega t-kz)-\sqrt{\frac{\varepsilon_0}{\mu_0}}E_{r0}\cos(\omega t+kz)\\
&=\sqrt{\frac{\varepsilon_0}{\mu_0}}E_{i0}\cos(\omega t-kz)+\sqrt{\frac{\varepsilon_0}{\mu_0}}E_{i0}\cos(\omega t+kz)\\
&=\sqrt{\frac{\varepsilon_0}{\mu_0}}E_{i0}\bigl[(\cos\omega t\cos kz+\sin\omega t\sin kz)+(\cos\omega t\cos kz-\sin\omega t\sin kz)\bigr]\\
&=2\sqrt{\frac{\varepsilon_0}{\mu_0}}E_{i0}\cos kz\cos\omega t.
\end{align*}

(2) On the two sides of the conductor surface $H$ is discontinuous, which shows that there is a surface current on the conductor surface. Applying the Amp\`ere circuital law to the loop $abcda$ in the figure, we have
\[
\oint\vec{H}\cdot\dd\vec{l}=H\cdot\overline{ab}
=2\sqrt{\frac{\varepsilon_0}{\mu_0}}E_{i0}\cos\omega t\cdot\overline{ab}
=i_c\cdot\overline{ab},
\]
\[
\therefore\ i_c=2\sqrt{\frac{\varepsilon_0}{\mu_0}}E_{i0}\cos\omega t\,\hat{i},
\]
with its direction along the $x$ axis. The force acting on unit area of the conductor surface is
\begin{align*}
\vec{f}&=\vec{i}_c\times\bar{\vec{B}}=\vec{i}_c\times\mu_0\bar{\vec{H}}
=\frac{1}{2}\vec{i}_c\times\mu_0\vec{H}_{\text{surface}}\\
&=\frac{1}{2}\cdot2\sqrt{\frac{\varepsilon_0}{\mu_0}}E_{i0}\cos\omega t\,\hat{i}
\times\mu_0\cdot2\sqrt{\frac{\varepsilon_0}{\mu_0}}E_{i0}\cos\omega t\,\hat{j}\\
&=2\varepsilon_0E_{i0}^2\cos^2\omega t\,\hat{k}.
\end{align*}
Because the average magnetic field acting on the current should be half the magnetic field at the conductor surface, the expression for $\vec{f}$ has an extra factor ``$\frac{1}{2}$'' in front. Since the time average of $\cos^2\omega t$ equals $1/2$, the average value of $\vec{f}$ is
\begin{align*}
\bar{f}&=2\varepsilon_0E_{i0}^2\overline{\cos^2\omega t}\,\hat{k}\\
&=\varepsilon_0E_{i0}^2\hat{k}.
\end{align*}

\begin{figbox}{3.58}
\begin{tikzpicture}[scale=0.95,>=latex]
  \draw[thickline] (0,-1.30)--(0,1.30);
  \fill[hatch] (0,-1.30) rectangle (0.60,1.30);
  \draw[thickline] (0.60,-1.30)--(0.60,1.30);
  \draw[thickline] (-0.90,0.85)--(-0.90,-0.85)--(0.35,-0.85)--(0.35,0.85)--cycle;
  \node[above left] at (-0.90,0.85) {$a$};
  \node[below left] at (-0.90,-0.85) {$b$};
  \node[below right] at (0.35,-0.85) {$c$};
  \node[above right] at (0.35,0.85) {$d$};
  \node[right] at (1.05,0.95) {conductor};
  \draw[->] (1.00,0.90)--(0.58,0.62);
  \draw[axis,->] (0,-1.55)--(1.55,-1.55) node[right] {$z$};
  \draw[axis,->] (0,-1.45)--(0,-2.30) node[below] {$y$};
\end{tikzpicture}
\end{figbox}

\begin{commentary}
In this problem, since the electromagnetic wave is completely reflected by the metal surface, the force per unit area on the conductor surface equals twice the product of the average momentum density $G$ of the electromagnetic field and the speed of light $c$, i.e. $\bar{f}=2Gc$; it also equals twice the ratio of the Poynting vector to the speed of light, i.e. $\bar{f}=2S/c$, and is numerically exactly twice the average energy density of the electromagnetic field.
\end{commentary}
\end{solution}

